\chapter{Chiral Hochschild cohomology and deformation theory}
\label{chap:hochschild-calculus}\label{ch:hochschild-cohomology}

\index{Hochschild cohomology!classical}
\index{cyclic homology!chiral}
\index{three Hochschild theories!disjoint discipline}

\begin{remark}[Hochschild theories are separated by their input]
\label{rem:hh-three-theories-discipline}

Hochschild notation in this chapter is typed by its input. The
curve-level cochain theory
$\ChirHoch^\bullet(\cA)$ takes a chiral algebra on a smooth curve~$X$
and is the object controlled by Theorem~H through a family support
datum $H_H(\cA;S)$.  The ordinary algebraic cochain
complex of a dg algebra~$A$ is
$\mathrm{HH}^\bullet_{\mathrm{alg}}(A)=\RHom_{A^e}(A,A)$. The
categorical Hochschild cochains of a dg category~$\mathcal C$ are
the derived endomorphisms of its identity functor. The topological
Hochschild homology $\mathrm{THH}(R)$ takes an $E_1$-ring spectrum
and returns an $S^1$-equivariant spectrum.

The later sections use these theories on their native surfaces:
the cyclic section uses Hochschild \emph{chains} and Connes' mixed
complex; the derived-centre section uses categorical/algebraic
cochains through a compact generator; the annulus section uses
algebraic/categorical chains; spectral $\mathrm{THH}$ appears only
after spectrification.  Theorem~H has source
$C^\bullet_{\mathrm{ch}}(\cA,\cA)$; ordinary algebraic, categorical,
topological, and product-ambient Hall Hochschild theories enter through
named comparison functors.
\end{remark}

For a factorisation module model of a chiral algebra~$\cA$ on a
curve, the shifted complex $\ChirHoch^\bullet(\cA,\cA)[1]$ is the
tangent complex at the Maurer--Cartan point~$\Theta_\cA$. It controls
chiral deformations through the chiral derived centre
$Z^{\mathrm{der}}_{\mathrm{ch}}(\cA)$, not through the ordinary
centre except on $H^0$.

The chiral Hochschild complex carries a degree $-1$ Lie bracket
(the chiral Gerstenhaber bracket, Theorem~\ref{thm:gerstenhaber-structure}),
computed by OPE residues on FM configuration spaces.
In the classical limit (fiber at a point), it recovers
Gerstenhaber's bracket~\cite{Ger63} on associative Hochschild
cochains. The chiral bracket is the strict shadow of a full
$L_\infty$ bracket on the chiral Hochschild complex. In the
classical algebraic strict model,
$\mathrm{HH}^2_{\mathrm{alg}}$ classifies first-order
deformations and $\mathrm{HH}^3_{\mathrm{alg}}$ contains
obstructions;
the $L_\infty$ extension records coherent homotopies governing
convergence of formal deformations and transfer of Koszul properties
across quasi-isomorphisms. The cyclic structure (Connes'
operator~$B$, the mixed complex, the SBI sequence) is the algebraic
shadow of the modular operad's gluing maps. The shadow obstruction
tower of Chapter~\ref{sec:koszul-across-genera} projects onto the
cyclic deformation complex
$\mathrm{Def}_{\mathrm{cyc}}^{\mathrm{mod}}(\cA)$ defined here.

For a PBW chiral datum on a smooth curve, Theorem~H computes the
cohomological support from a supplied family retract
$H_H(\cA;S)$.  Bakalov--De Sole--Kac compute the bounded Virasoro
support $\{0,2,3\}$ and the rank-one even superboson dimensions
$(2,1,0,\ldots)$; a bounded-to-chart quasi-isomorphism transports
these values to the curve-level complex.  Principal $W$-algebras,
simple quotients, and product-ambient Calabi--Yau objects require
their own chain models and comparison maps.

The fiberwise curvature $\dfib^{\,2} = \kappa(\cA) \cdot \omega_1$ at
genus~$1$ (Chapter~\ref{ch:heisenberg-frame}) shows that deformations
of the bar complex are controlled by a cohomological invariant.
Chiral Hochschild cohomology organizes these deformations and obstructions
into a single graded object equipped with a Gerstenhaber bracket;
the chiral enhancement incorporates OPE data, the geometry of the
curve, and the $A_\infty$ structure of the bar complex.

\begin{remark}[Conventions for HH/HC/HP]
\label{rem:hh-hc-canonical-source}
\index{Hochschild cohomology!canonical conventions}
\index{cyclic homology!canonical conventions}
\index{periodic cyclic homology!canonical conventions}
The definitions and sign conventions for classical algebraic
Hochschild cohomology $\mathrm{HH}^\bullet_{\mathrm{alg}}$,
classical algebraic Hochschild homology
$\mathrm{HH}_{\bullet,\mathrm{alg}}$, cyclic homology
$\mathrm{HC}_*$, negative cyclic homology $\mathrm{HC}_*^{\mathrm{neg}}$,
and periodic cyclic homology $\mathrm{HP}_*$ are as follows. In particular:
\begin{itemize}
\item the chiral Hochschild complex
 (Definition~\ref{def:chiral-hochschild}),
\item the cyclic operator and Connes' operator~$B$
 (Definitions~\ref{def:cyclic-operator} and~\ref{def:connes-B}),
\item the mixed complex $(d_{\mathrm{HH}}, B)$ and the SBI exact
 sequence (Theorem~\ref{thm:connes-exact-sequence},
 Corollary~\ref{cor:connes-SBI}),
\end{itemize}
are all defined here. These definitions govern the chiral Hochschild,
topological Hochschild, and cyclic-invariant conventions used by the
deformation theory and Poincar\'e duality chapters.
\end{remark}

\begin{remark}[Chiral Kontsevich formality: $\ChirHoch^\bullet$ and
chiral polyvector fields]
\label{rem:hh-chiral-kontsevich-formality-pointer}
\index{Kontsevich formality!chiral!pointer from Hochschild chapter}%
The chiral Hochschild complex $\ChirHoch^\bullet(\mathbf H_{\Delta_5})$
fits Kontsevich's $1997$ formality framework in the BD
factorisation-algebra setting through
Theorem~\ref{thm:chiral-kontsevich-formality} of
Chapter~\ref{chap:mc5-class-m-chain-level-platonic}: on the
Humbert-free complement
$\mathcal U_{\mathrm{Humb}}^{\mathrm{c}}=
\overline{\mathcal A}_2\setminus\bigcup_n H_n$ there is an
$L_\infty$-quasi-isomorphism
$T^\bullet_{\mathrm{ch}}(\overline{\mathcal A}_2,\mathbf H_{\Delta_5})
\xrightarrow{\sim}\ChirHoch^\bullet(\mathbf H_{\Delta_5})$
between chiral polyvector fields and chiral Hochschild cochains,
transporting Kontsevich's admissible-graph weights along the KZ
pullback of the Climax Theorem. The formality obstruction on the
full Humbert locus $\bigcup_n H_n$ is the Bruinier Heegner Chern
class $c_n(\mathcal L^{\Delta_5}|_{H_n})$ of Bruinier 2002
\textsc{LNM} 1780 Prop.~5.1, equal to the Costello--Gwilliam
$1$-loop BV counterterm and to the Deligne--Malcev bi-unipotent
classifying cocycle of
Remark~\ref{rem:formality-obstruction-malcev-ladder}. Strict
chain-level bar--cobar inversion
$\Omega^{\mathrm{ch}}(B^{\mathrm{ch}}(\mathbf H_{\Delta_5}))
\xrightarrow{\sim}\mathbf H_{\Delta_5}$ on
$\mathcal U_{\mathrm{Humb}}^{\mathrm{c}}$ is then a corollary of
chiral formality through the Kontsevich--Soibelman transfer lemma,
paired against the Koszul pairing
$\mathbf H_{\Delta_5}\otimes\mathbf H_{\Delta_5}^!\to k[-2]$ of
Theorem~\ref{thm:main-koszul-hoch}. The unified form of the
theorem--five classical ingredients (Kontsevich admissible graphs,
Costello--Gwilliam factorisation polyvector fields, Francis--Gaitsgory
BD chiralisation, bi-unipotent factorisation compatibility,
Malcev ladder cocycle) assembled into one chain-level
statement with the per-wall obstruction
$\mathrm{ob}^{\mathrm{Kont}}_n\in H^3(H_n,\mathrm{Sym}^2
T^{\mathrm{poly}}_{\mathrm{ch}}|_{H_n})$ identified with the $n$-th
Malcev ladder rung classifying cocycle
$[\mathrm{gr}_n^{\mathrm{Malcev}}]$--is
Theorem~\ref{thm:chiral-kontsevich-formality-unified}. Primary:
Kontsevich 1997 \textsc{arXiv}:q-alg/9709040 Thm.~4.1;
Costello--Gwilliam \cite[\S\S 6--8]{CG17}; Francis 2013 \textit{Adv.\
Math.}\ 239 Thm.~1.1; Bruinier 2002 \textsc{LNM} 1780 Prop.~5.1;
Hain--Matsumoto 2003 \textit{JAMS}\ 16 Thm.~5.1; Cattani--Kaplan--Schmid
1986 \textit{Ann.\ Math.}\ 123 Thm.~3.3 \S3.7.
\end{remark}

\section{Classical to chiral}

\subsection{Review of classical Hochschild}

For an associative algebra $A$ over $\mathbb{C}$ and an
$A$-bimodule~$M$, the classical algebraic Hochschild cohomology
$\mathrm{HH}^\bullet_{\mathrm{alg}}(A, M) = \mathrm{Ext}^*_{A^e}(A, M)$
classifies the deformation theory of~$A$:
$\mathrm{HH}^0_{\mathrm{alg}}$ is the center,
$\mathrm{HH}^1_{\mathrm{alg}}$ classifies outer derivations,
$\mathrm{HH}^2_{\mathrm{alg}}$ classifies first-order deformations,
and $\mathrm{HH}^3_{\mathrm{alg}}$ contains the obstructions to
extending them. The full structure is computed by the bar resolution
$\cdots \to A^{\otimes(n+2)} \xrightarrow{b} A^{\otimes(n+1)} \to
\cdots$ with differential $b(a_0 \otimes \cdots \otimes a_{n+1})
= \sum_{i=0}^n (-1)^i a_0 \otimes \cdots \otimes a_ia_{i+1}
\otimes \cdots \otimes a_{n+1}$, and the resulting cohomology
carries a Gerstenhaber bracket governing the Maurer--Cartan equation
for formal deformations \textup{(}Gerstenhaber~\cite{Ger63}\textup{)}.

\subsection{Chiral enhancement}

What replaces the bimodule $A \otimes A^{\mathrm{op}}$ when the algebra lives on a curve rather than a point? For chiral algebras, each ingredient of the classical theory acquires a geometric upgrade: the bimodule structure is replaced by the factorization structure on~$\operatorname{Ran}(X)$, the bar resolution becomes the geometric bar complex on Fulton--MacPherson configuration spaces, and the Gerstenhaber bracket lifts to a chiral bracket governed by the OPE. The result is a chiral Hochschild complex whose cohomology computes chiral deformations and obstructions; its degree-$0$ part is the ordinary chiral centre, while the full complex is the derived chiral centre.

\begin{definition}[Chiral Hochschild complex, derived-functor form]
\label{def:chiral-hochschild-derived}
For a chiral algebra $\mathcal{A}$ on $X$, the chiral Hochschild
cohomology with coefficients in a chiral $\mathcal{A}$-module
$\mathcal{M}$ is
\[\ChirHoch^*(\mathcal{A}, \mathcal{M})
= \operatorname{RHom}_{\mathcal{D}_X}(\barBgeom(\mathcal{A}), \mathcal{M}).\]
The chain-level model is Definition~\ref{def:chiral-hochschild}.
\end{definition}

\begin{theorem}[Comparison with classical theory {\cite{BD04}}; \ClaimStatusProvedElsewhere]\label{thm:hochschild-classical-comparison}
There is a spectral sequence:
\[
E_2^{p,q}
= \mathrm{HH}^p_{\mathrm{alg}}(\mathcal{A}_0, H^q(\Omega^*_X))
\Rightarrow \ChirHoch^{p+q}(\mathcal{A})
\]
where $\mathcal{A}_0$ is the fiber at a point.
\end{theorem}

\section{Periodicity phenomena}

\subsection{Virasoro chiral Hochschild cohomology}

\begin{theorem}[Virasoro bounded cohomology and chart comparison;
\ClaimStatusConditional]\label{thm:virasoro-hochschild}
\textup{[Type signature: Open quadrant; bounded vertex-cohomology
presentation and curve-level chart presentation; Beilinson level~$3$;
hypothesis package: the Bakalov--De Sole--Kac complex and a
bounded-to-chart quasi-isomorphism $\chi_{\operatorname{Vir}_c}^{\mathrm{bd}}$.]}
For every central charge~$c$, Bakalov--De Sole--Kac compute
\[
 \dim H^n_{\mathrm{ch},b}(\operatorname{Vir}_c,
 \operatorname{Vir}_c)
 =\begin{cases}
 1,&n\in\{0,2,3\},\\
 0,&n\in\mathbb Z_{\geq0}\setminus\{0,2,3\}.
 \end{cases}
\]
Thus the bounded Poincar\'e polynomial is $1+t^2+t^3$.  If
\[
 \chi_{\operatorname{Vir}_c}^{\mathrm{bd}}\colon
 C^\bullet_{\mathrm{ch},b}(\operatorname{Vir}_c,
 \operatorname{Vir}_c)
 \xrightarrow{\ \sim\ }
 C^\bullet_{\mathrm{ch}}(\cA_{\operatorname{Vir}_c},
 \cA_{\operatorname{Vir}_c})
\]
is a quasi-isomorphism, the same dimensions compute the curve-level
chart complex.
\end{theorem}

\begin{proof}
The bounded calculation is \cite[Theorem~7.2]{BDSK21}.  The chart
statement is transport of cohomology along the displayed
quasi-isomorphism.
\end{proof}

\begin{remark}[Relation to Gel'fand--Fuchs cohomology]
\label{rem:gf-vs-chirhoch}
The continuous Lie algebra cohomology
$H^*_{\mathrm{cont}}(L_1; \mathbb{C}) = \mathbb{C}[\Theta]$ is a
polynomial ring on a degree-$2$ generator
(Gel'fand--Fuchs~\cite{Fuks86}).  It is the cohomology of a
topological Lie algebra, whereas
$\ChirHoch^*(\mathrm{Vir}_c)$ is the factorization Ext-group on the
curve~$X$.  The latter has support $\{0,2,3\}$ after the
Bakalov--De Sole--Kac bounded-to-chart comparison.  The two theories
are related by a
spectral sequence
(Theorem~\ref{thm:hochschild-classical-comparison}).  Its convergence
and edge maps determine the range in which one theory computes the
other.  The chart support itself is the output of a family datum
$H_H(\cA;S)$.
\end{remark}

\subsection{Affine Kac--Moody periodicity}

\begin{theorem}[Critical level Lie algebra cohomology; \ClaimStatusConditional]\label{thm:critical-level-cohomology}
The continuous Lie algebra cohomology of $\widehat{\mathfrak{g}}$ at critical level $k = \critLevel$ is the exterior-polynomial algebra
\[H^*_{\mathrm{cont}}(\widehat{\mathfrak{g}}_{\critLevel},
\widehat{\mathfrak{g}}_{\critLevel,0};\, \mathbb{C})
\;\cong\; \Lambda^*(P_1,\ldots,P_r) \otimes \mathbb{C}[\Theta_1,\ldots,\Theta_r]\]
with $\deg P_i = 2m_i + 1$ and $\deg \Theta_i = 2(m_i+1)$, where $m_1,\ldots,m_r$ are the exponents of~$\mathfrak{g}$ (Fuks--Feigin--Tsygan~\cite{FT87}, via BD comparison~\cite{BD04}). For rank~$1$ ($\mathfrak{sl}_2$): $H^{n+4} \cong H^n$ (strict periodicity). For rank $r > 1$: polynomial growth $O(n^{r-1})$ without strict periodicity (Theorem~\ref{thm:affine-periodicity-critical}).
\end{theorem}

\begin{remark}[Critical level: Lie cohomology vs chiral Hochschild]
\label{rem:critical-level-lie-vs-chirhoch}
At generic level, chiral Hochschild cohomology and continuous Lie
algebra cohomology are distinct functors.  A comparison requires the
Beilinson--Drinfeld chain map and its family-specific quasi-isomorphism
range.  Theorem~H computes the chart support from $H_H(\cA;S)$.

At critical level ($k=-h^\vee$), the Ext algebra carries the
Feigin--Frenkel centre $\mathrm{Fun}(\mathrm{Op})$ in bar degree
$p=0$ and positive conformal weights.  The critical chart therefore
requires a support datum adapted to this centre.

Under the Beilinson--Drinfeld chiral-envelope comparison hypothesis
(\cite[Theorem~4.8.1]{BD04}, specialised to the affine Lie-$*$ algebra
$\fg \otimes \omega_X$ at critical level and combined with the
Feigin--Frenkel identification of the chiral centre
\cite{Feigin-Frenkel}) identifies the chiral Hochschild
cohomology with the continuous Lie algebra cohomology:
\[
\ChirHoch^*(\widehat{\fg}_{-h^\vee})
\;\cong\;
H^*_{\mathrm{Lie,cont}}(\fg \otimes t\bC[[t]];\, \bC)
\;\cong\;
\Lambda^*(P_1, \ldots, P_r) \otimes \bC[\Theta_1, \ldots, \Theta_r],
\]
whose polynomial generators produce classes in arbitrarily large
cohomological degree with polynomial growth
(Theorem~\ref{thm:affine-periodicity-critical}).  The
Feigin--Frenkel centre supplies the separate degree-zero chiral-centre
action.

At generic and critical levels alike, the chart support is carried by
its explicit comparison complex.  The critical-level comparison above
is a separate conditional statement whose degree-zero term contains
the Feigin--Frenkel centre.
\end{remark}

The degrees in the exterior--polynomial Lie-cohomology formula are
determined by the exponents of~$\mathfrak g$; the corresponding
critical chiral centre is the Feigin--Frenkel oper algebra
\cite{Feigin-Frenkel}.  For simply-laced~$\mathfrak g$ one has
$h=h^\vee$; the two Coxeter numbers differ in the non-simply-laced
types.

\subsection{W-algebra chiral Hochschild cohomology}

For a principal $W$-algebra, free strong generation supplies a PBW
presentation of the associated graded object.  The cohomological
support requires the collision differential and its family retract.

\begin{theorem}[Family support criterion for principal $W$-algebras;
\ClaimStatusConditional]\label{thm:w-algebra-hochschild}
\textup{[Type signature: Open quadrant; principal $W$-algebra chart
presentation; Beilinson level~$3$; hypothesis package
$H_H(\Walg^k(\mathfrak g);S_W)$.]}
Let
\(\Walg^k(\mathfrak g)=\Walg^k(\mathfrak g,f_{\mathrm{prin}})\)
and assume a family support datum
\(H_H(\Walg^k(\mathfrak g);S_W)\).  Then
\[
 \ChirHoch^n(\Walg^k(\mathfrak g))
 \cong H^n(K_{\Walg^k(\mathfrak g),S_W}),
 \qquad
 \operatorname{Supp}\ChirHoch^\bullet(\Walg^k(\mathfrak g))
 \subseteq S_W.
\]
The determination of $S_W$ is the family-specific collision
calculation.  De Sole--Kac variational Poisson cohomology is
nonvanishing in unboundedly many degrees for Virasoro-type Poisson
vertex algebras \cite{BDSK21}, so this concentration lives strictly
on the stated Koszul/bounded locus, with the bounded-to-chart
comparison part of \(H_H(\Walg^k(\mathfrak g);S_W)\).
\end{theorem}

\begin{proof}
This is Theorem~\ref{thm:main-koszul-hoch} applied to the displayed
family support datum.
\end{proof}

\begin{remark}[Gel'fand--Fuchs cohomology vs chiral Hochschild]
\label{rem:gf-walg-vs-chirhoch}
The continuous relative cohomology of the topological Lie algebra
$\mathrm{Lie}(\Walg^k(\mathfrak{g}))$ with $r$ generating currents is a
polynomial ring
$H^*_{\mathrm{cont}}(\mathrm{Lie}(\Walg),
\mathrm{Lie}(\Walg)_0;\, \mathbb{C})
= \mathbb{C}[\Theta_1, \ldots, \Theta_r]$
with $\Theta_i \in H^{h_i}$
(Fuchs~\cite{Fuks86}, Ch.~2; Feigin--Fuchs~\cite{FF84} for $r = 1$).
Gel'fand--Fuchs cohomology is the continuous Lie algebra cohomology of
the mode algebra.  Chiral Hochschild cohomology is the curve-level
factorization Ext.  A comparison between them consists of an explicit
chain map and its quasi-isomorphism range; the support of the chart
complex is supplied by $H_H(\Walg^k;S_W)$.
For the Virasoro ($r = 1$): $H^*_{\mathrm{cont}}(L_1) = \mathbb{C}[\Theta]$
has Hilbert series $1/(1-t^2)$, while the bounded vertex-cohomology
complex has Poincar\'e polynomial $1+t^2+t^3$.  Transport of the latter
to the curve chart uses $\chi_{\operatorname{Vir}_c}^{\mathrm{bd}}$.
\end{remark}

\begin{remark}[Scope]
Free strong generation of the universal principal algebra supplies
its PBW associated graded.  A curve-level support calculation also
requires incidence-compatible collision retracts, completion,
averaging, and a comparison with the chart complex.  Admissible
quotients, simple quotients, critical level, and non-principal
nilpotent orbits each define a separate family datum
$H_H(\Walg^k(\mathfrak g,f);S_{k,f})$.  At critical level the
Feigin--Frenkel centre enters this datum in degree zero.
\end{remark}

\begin{remark}[Curve-level model for Theorem~H]
\label{rem:theorem-h-model}
\index{chiral Hochschild cohomology!model presentation}
The curve-level chiral derived centre is computed from the two-sided
diagonal bar object
\[
 \bar{B}_{\mathrm{HH}}(\cA)
 =
 \bigl(\cA\otimes T^c(s^{-1}\bar\cA)\otimes\cA,\,d\bigr),
\]
viewed as a model for
$R\!\operatorname{End}_{\cA\text{-}\cA}(\cA)$.  The bar--cobar counit
resolves the diagonal bimodule under the Ran reconstruction package.
A quadratic presentation supplies
\[
 \cA^{\mathrm i}
 \xrightarrow{\ q_\cA\ }
 \barB^{\mathrm{ch}}(\cA),
 \qquad
 \cA^{\mathrm i}
 \xrightarrow{\mathbb D_X}
 \cA^!,
\]
and Koszulness makes $q_\cA$ a quasi-isomorphism.  These are the full
bar coalgebra, the quadratic coalgebra, and its Verdier-dual algebra.

Theorem~H adds a family support datum
$H_H(\cA;S)$ to this resolution.  Its incidence-compatible retract
identifies the chart cohomology with $H^\bullet(K_{\cA,S})$.  The
shifted chart complex governs the chiral Maurer--Cartan deformation
problem; its degree-three group carries the primary obstruction.
For principal $W$-algebras, determining that group is part of the
family-specific collision calculation.
\end{remark}

\begin{remark}[Non-principal nilpotent orbits]
For $\Walg^k(\mathfrak g,f)$, the Kazhdan grading associated to the
$\mathfrak{sl}_2$-triple determines the conformal weights of the
generators.  A support theorem is the construction of
$H_H(\Walg^k(\mathfrak g,f);S_{k,f})$.  Its degree-zero term records
the centre, its degree-one term records fixed-product derivations, and
its higher terms record deformations and their obstructions.  Each
nilpotent orbit therefore carries its own collision and comparison
problem.
\end{remark}

\section{Deformation theory}

\subsection{Infinitesimal deformations}

\begin{theorem}[Low-degree chiral deformation classes;
\ClaimStatusConditional]\label{thm:deformation-classification}
Assume the chiral deformation functor of $\mathcal A$ is governed by
the shifted dg-Lie algebra
$\ChirHoch^\bullet(\mathcal A,\mathcal A)[1]$ with the chiral
Gerstenhaber bracket of Theorem~\textup{\ref{thm:gerstenhaber-structure}}.
Then:
\begin{enumerate}
\item $\ChirHoch^1(\mathcal{A})$ classifies \emph{outer derivations} $\mathrm{Der}(\mathcal{A})/\mathrm{Inn}(\mathcal{A})$ (not infinitesimal automorphisms, which are all derivations).
\item $\ChirHoch^2(\mathcal{A})$ classifies first-order deformations.
\item $\ChirHoch^3(\mathcal{A})$ contains the primary obstructions to extending deformations.
\item In conformal-field-theoretic applications, exactly marginal
operators are weight-one classes in $\ChirHoch^2(\mathcal A)$ whose
primary obstruction in $\ChirHoch^3(\mathcal A)$ vanishes.
\end{enumerate}
\end{theorem}

\begin{example}[Marginal deformations of beta-gamma; standard
free-field normalisation]
For the standard $\beta\gamma$ system in the free-field normalisation:
\[\ChirHoch^2(\beta\gamma) = \mathbb{C} \cdot [\beta\gamma]\]
The class $[\beta\gamma]$ corresponds to the exactly marginal operator changing the conformal weights.
\end{example}

\subsection{Formal deformation theory}

The formal deformation space is controlled by the differential graded Lie algebra
$\mathfrak{g}_{\mathcal{A}} = \ChirHoch^*(\mathcal{A}, \mathcal{A})[1]$
with the Gerstenhaber bracket.

\begin{theorem}[Maurer--Cartan equation {\cite{Kon03,KontsevichSoibelman}};
\ClaimStatusConditional]\label{thm:maurer-cartan-deformations}
Under the governing-dg-Lie hypothesis of
Theorem~\textup{\ref{thm:deformation-classification}}, formal
deformations correspond to solutions of:
\[d\alpha + \frac{1}{2}[\alpha, \alpha] = 0, \quad \alpha \in \mathfrak{g}_{\mathcal{A}}^1\]
\end{theorem}

\section{Physical applications}

\subsection{Marginal operators and RG flow}

In CFT, marginal operators have dimension $(1,1)$ and correspond to
$\ChirHoch^2_{\text{marginal}}(\mathcal{A}) = \{\omega \in \ChirHoch^2 : h(\omega) = 1\}$.
The primary obstruction to extending such a deformation is the
corresponding class in $\ChirHoch^3$; vanishing of all higher
renormalisation obstructions requires the full deformation problem,
not only Theorem~H.

\subsection{String field theory}

The conjectural identification of the bar complex with BRST cohomology (Conjecture~\ref{conj:master-bv-brst}) takes the form:
\[H^*_{\text{BRST}}(\text{String}[\mathcal{A}]) \cong \ChirHoch^*(\mathcal{A})\]

String vertices are encoded in the $A_\infty$ structure: $m_2$ gives the three-string vertex, $m_3$ the four-string contact term, and higher $m_k$ the multi-string interactions.

\section{Computational tools}

\subsection{Spectral sequences}

\begin{theorem}[Bar spectral sequence {\cite{BD04,CG17}}; \ClaimStatusProvedElsewhere]\label{thm:bar-spectral-sequence-hochschild}
\[E_1^{p,q} = H^q(\ConfigSpace{p}(X), \mathcal{A}^{\boxtimes p}) \Rightarrow \ChirHoch^{p+q}(\mathcal{A})\]
\end{theorem}

\subsection{Explicit computations}

Bakalov--De Sole--Kac compute bounded vertex cohomology for two basic
families.  For the rank-one even free superboson (B_{\mathfrak h}),
\cite[Theorem~7.4]{BDSK21} gives
\[
 \dim H^n_{\mathrm{ch},b}(B_{\mathfrak h},B_{\mathfrak h})
 =(2,1,0,0,\ldots).
\]
For the Virasoro vertex algebra, \cite[Theorem~7.2]{BDSK21} gives one
class in each of degrees (0,2,3) and zero in every other degree.
A bounded-to-chart quasi-isomorphism
\[
 \chi_V^{\mathrm{bd}}\colon
 C^\bullet_{\mathrm{ch},b}(V,V)
 \longrightarrow C^\bullet_{\mathrm{ch}}(\cA_V,\cA_V)
\]
transports either calculation to the curve-level chart complex.

For the free-fermion OPE
\(psi(z)\psi^*(w)\sim(z-w)^{-1}\), the rescaling cochain
\[
 N(\psi)=\psi,\qquad N(\psi^*)=0
\]
has coboundary equal to the corresponding rescaling of the mixed OPE.
Thus the degree-two cochain represented by
\(\psi\psi^*\) is exact in this model.  A free-fermion chart
calculation requires a family support datum
\(H_H(\cA_{\mathcal F};S_{\mathcal F})\) or a proved
bounded-to-chart comparison, as in
Example~\ref{ex:HH-fermion-complete}.

\section{Chiral Hochschild-cyclic spectral sequence for chiral algebras}
\label{sec:hochschild-cyclic-spectral}

\begin{remark}[Physical origin]
\emph{Heuristic.} In 2d~CFT on a genus~$g$ surface $\Sigma_g$, the partition function
$Z[\Sigma_g] = \operatorname{Tr}_{\mathcal{H}}(q^{L_0-c/24})$ is
read by the chiral Hochschild-cyclic spectral sequence in the
following physical mnemonic: the $E_1$ page
gives the tree-level (genus~$0$) contribution, $E_2$ incorporates
one-loop corrections, and higher pages capture multi-loop effects.
For chirally Koszul algebras at generic level, the spectral sequence
degenerates at~$E_2$ by the Ext-vanishing argument
(Theorem~\ref{thm:w-algebra-hochschild}).
\end{remark}

\begin{construction}[Topological cyclic comparison; conditional]
The chiral Hochschild chain model used for cyclic homology has
chain groups
\[
\mathrm{CH}^{\mathrm{ch}}_n(\mathcal{A}) =
\Gamma(C_n(X),
 \mathcal{A}^{\boxtimes n} \otimes \det(\Omega^1_{C_n(X)/X}))\]
before compactification and derived completion.

In a locally constant topological comparison model, cyclic rotation of
configuration points gives an $S^1$-action and a comparison map toward:
\[
\mathrm{CC}^{\mathrm{ch,top}}_n(\mathcal{A}) \longrightarrow
 [\mathrm{Maps}(S^1, LX) \otimes_{\mathcal{A}} \mathcal{A}^{\otimes n}]^{S^1}
\]
where $LX = \text{Maps}(S^1, X)$ is the loop space.

Under the hypotheses of that comparison, the cyclic spectral sequence
maps to the spectral sequence for $S^1$-equivariant loop homology.
This is not the curve-level Theorem~H statement. Theorem~H concerns
the chiral cochain complex $\ChirHoch^\bullet(\cA)$ on a curve, with
the conventions of Appendix~\ref{app:hochschild-conventions}.
The classical loop-space identification for associative algebras is
\cite{Loday98}; the factorization comparison requires the additional
locally constant/topological hypotheses of \cite[Chapter~5]{CG17}.
\end{construction}

\begin{computation}[Locally constant zero-mode model; \ClaimStatusConditional]
In the locally constant zero-mode comparison model, the $E_2$ page for
standard examples is:

\emph{Heisenberg algebra.}
\[\text{HH}_*^{\text{ch}}(\mathcal{H}) = \mathbb{C}[c] \otimes \Lambda(\sigma)\]
where $|c| = 2$ (central charge) and $|\sigma| = 1$ (desuspension of $S^1$).

The Connes operator $B: \text{HH}_n \to \text{HH}_{n+1}$ (degree $+1$) satisfies $B(c) = 0$,
$B(\sigma) = c$, giving:
\[\text{HC}_*^{\text{ch}}(\mathcal{H}) = \mathbb{C}[c]\]
(polynomials in the central charge).

\emph{Free fermion $\beta\gamma$ system.}
\[\text{HH}_*^{\text{ch}}(\beta\gamma) = \Lambda(\beta_0, \gamma_0) \otimes \mathbb{C}[c]\]
where $\beta_0, \gamma_0$ are zero modes. Cyclic homology adds periodicity:
\[\text{HC}_*^{\text{ch}}(\beta\gamma) = \mathbb{C}[c] \otimes
 \mathbb{C}[\beta_0^{-1}, \gamma_0^{-1}]_{\text{formal}}\]
(the localization $\mathbb{C}[\beta_0, \gamma_0][\beta_0^{-1}, \gamma_0^{-1}]$ at the zero-mode variables, completed formally).

The zero-mode periodicity in $\text{HC}_*^{\text{ch}}(\beta\gamma)$ reflects the symplectic pairing of the $\beta\gamma$ system: the localized zero modes span the deformation space of the Lagrangian splitting.
\end{computation}

\begin{remark}[Functoriality]
The chiral Hochschild-cyclic spectral sequence is functorial: a morphism
$f\colon \cA \to \mathcal{B}$ induces $f_*\colon E_r(\cA) \to E_r(\mathcal{B})$,
tensor products satisfy $E_r(\cA \otimes \mathcal{B}) \cong E_r(\cA) \otimes E_r(\mathcal{B})$
up to the Kunneth convergence hypotheses for the chosen completed
chain model. Koszul duality acts on the cyclic spectral sequence only
after a trace/Poincare duality identifies the chain model with the
dual of the chiral cochain model; without that extra datum no
pagewise equality $E_r(\cA^!)\cong E_r(\cA)^\vee$ is asserted.
The representing object for deformations is the cochain complex
$\ChirHoch^\bullet(\cA,\cA)$; the chain complex
$\mathrm{CH}_*(\cA)$ represents traces and cyclic invariants.
\end{remark}

\subsection{Chiral Hochschild complex for chiral algebras}

\begin{definition}[Chiral Hochschild chain model for cyclic homology]
\label{def:chiral-hochschild}
\index{Hochschild cohomology!chiral|textbf}
For a chiral algebra $\mathcal{A}$ on a smooth curve $X$, the
chain-side model used below is
\[\text{CH}_n(\mathcal{A}) = \Gamma\left(\overline{C}_{n+1}(X),
 \mathcal{A}^{\boxtimes(n+1)} \otimes \det\left(\Omega^1_{\overline{C}_{n+1}(X)/X}\right)\right)\]

with differential:
\[d_{\text{HH}}(\omega) = \sum_{i=0}^{n} (-1)^i \left[\text{omit } i\text{-th factor}\right]\]

More explicitly, using the chiral product $\mu_{ij}$:
\begin{align*}
d_{\text{HH}}(a_0 \otimes \cdots \otimes a_n)
&= \sum_{i=0}^{n-1} (-1)^i (a_0 \otimes \cdots \otimes \mu(a_i, a_{i+1})
 \otimes \cdots \otimes a_n) \\
&\quad + (-1)^n (\mu(a_n, a_0) \otimes a_1 \otimes \cdots \otimes a_{n-1})
\end{align*}

The last term implements the \emph{cyclic structure}: $a_n$ wraps around to multiply $a_0$.
This is not the cochain complex
$\ChirHoch^\bullet(\mathcal A)$ of Theorem~H; it is the chain model
whose homology feeds Connes' cyclic construction.
\end{definition}

\begin{theorem}[Chiral Hochschild chain model is a chain complex; \ClaimStatusProvedHere]
\label{thm:hochschild-chain-complex}
The identity $d_{\text{HH}}^2 = 0$ is the chiral Hochschild analogue
of the Arnold-relation mechanism in the degree-$2$ bar calculation:
adjacent face maps cancel in pairs. In the chiral setting, the same
cancellation acquires OPE corrections from the higher poles.
The differential $d_{\text{HH}}$ satisfies $d_{\text{HH}}^2 = 0$, making
$(\text{CH}_*(\mathcal{A}), d_{\text{HH}})$ a chain complex.
\end{theorem}

\begin{proof}
We must show that for any $\omega \in \text{CH}_n(\mathcal{A})$:
\[d_{\text{HH}}^2(\omega) = 0\]

\emph{Step 1: Expand $d_{\text{HH}}^2$}

Applying $d_{\text{HH}}$ twice gives:
\begin{align*}
d_{\text{HH}}^2(a_0 \otimes \cdots \otimes a_n)
&= d_{\text{HH}}\left(\sum_{i=0}^{n-1} (-1)^i (\cdots \otimes \mu(a_i, a_{i+1})
 \otimes \cdots)\right. \\
&\quad\left. + (-1)^n (\mu(a_n, a_0) \otimes a_1 \otimes \cdots)\right)
\end{align*}

\emph{Step 2: Identify canceling pairs}

After expanding, terms come in two types:
\begin{enumerate}
\item \emph{Type A}: Apply $\mu$ at positions $(i,i+1)$ then $(j,j+1)$ with $j \neq i, i+1$
\item \emph{Type B}: Apply $\mu$ twice at adjacent triples $(i,i+1,i+2)$
\end{enumerate}

\emph{Type A terms cancel} because applying $\mu$ at position~$i$ first reduces the chain length by~1, so the second application at position~$j$ (with $j > i+1$) picks up index $j-1$:
\[(-1)^i (-1)^{j-1} (\ldots) +
 (-1)^j (-1)^{i} (\ldots) = (-1)^{i+j-1} + (-1)^{i+j} = 0\]
(the index shift from the first contraction produces the sign difference)

\emph{Type B terms cancel} by the Borcherds identity (the chiral analog of associativity).
For three elements $a_i, a_{i+1}, a_{i+2}$ at pairwise distinct points $z_i, z_{i+1}, z_{i+2}$,
the Type~B contributions are:
\begin{align*}
&(-1)^i(-1)^{i+1}(\cdots \otimes \mu_{z_i, z_{i+2}}(\mu_{z_i, z_{i+1}}(a_i,a_{i+1}),a_{i+2}) \otimes \cdots) \\
&+ (-1)^{i+1}(-1)^i (\cdots \otimes \mu_{z_i, z_{i+2}}(a_i,\mu_{z_{i+1}, z_{i+2}}(a_{i+1},a_{i+2})) \otimes \cdots)
\end{align*}

The signs give $(-1)^{2i+1} = -1$ for the first and $(-1)^{2i+1} = -1$ for the second,
so the total contribution is:
\[-\bigl(\mu(\mu(a_i,a_{i+1}),a_{i+2}) - \mu(a_i,\mu(a_{i+1},a_{i+2}))\bigr)\]

By the Borcherds identity \cite[Section~3.4.7]{BD04}
(equivalently: the factorization algebra structure
encodes \emph{homotopy} associativity via the codimension-1 boundary strata
of $\overline{C}_3$), the iterated products
$\mu(\mu(a_i,a_{i+1}),a_{i+2})$ and $\mu(a_i,\mu(a_{i+1},a_{i+2}))$
agree up to the contribution of the boundary stratum $\partial \overline{C}_3$,
which is accounted for by the chiral Hochschild complex structure
(the face maps encode precisely these boundary identifications).
Thus the two terms cancel in the chiral Hochschild complex, giving $d_{\text{HH}}^2 = 0$.

\emph{Note}: Unlike the classical case, the chiral product is \emph{not} strictly associative;
the Borcherds identity provides associativity only up to total derivatives
(equivalently, up to homotopy). The cancellation above relies on the
factorization algebra axioms rather than strict associativity.
\end{proof}

\begin{definition}[Chiral Hochschild homology]
\label{def:chiral-HH}
The \emph{chiral Hochschild homology} of $\mathcal{A}$ is:
\[\text{HH}_n^{\text{ch}}(\mathcal{A}) = H_n(\text{CH}_*(\mathcal{A}), d_{\text{HH}})\]
\end{definition}

\subsection{\texorpdfstring{Cyclic structure and $S^1$-action}{Cyclic structure and S 1-action}}

\begin{definition}[Cyclic operator]
\label{def:cyclic-operator}
The \emph{cyclic operator} $t: \text{CH}_n(\mathcal{A}) \to \text{CH}_n(\mathcal{A})$
is defined by:
\[t(a_0 \otimes a_1 \otimes \cdots \otimes a_n) = (-1)^n (a_n \otimes a_0 \otimes a_1
 \otimes \cdots \otimes a_{n-1})\]

This generates a $\mathbb{Z}/(n+1)$-action on $\text{CH}_n(\mathcal{A})$.
\end{definition}

\begin{lemma}[Cyclic operator commutes with the chiral Hochschild differential; \ClaimStatusProvedHere]
\label{lem:cyclic-commutes}
The cyclic operator commutes with the chiral Hochschild differential:
\[t \circ d_{\text{HH}} = d_{\text{HH}} \circ t\]
\end{lemma}

\begin{proof}
Direct computation using the definition of $d_{\text{HH}}$:
\begin{align*}
t(d_{\text{HH}}(a_0 \otimes \cdots \otimes a_n))
&= t\left(\sum_{i=0}^{n-1} (-1)^i (\cdots \otimes \mu(a_i,a_{i+1}) \otimes \cdots)
 + (-1)^n (\mu(a_n,a_0) \otimes \cdots)\right)
\end{align*}

After applying $t$ (cyclic permutation), each term shifts indices: $i \to i+1 \pmod{n+1}$.

Similarly:
\[d_{\text{HH}}(t(a_0 \otimes \cdots \otimes a_n)) =
 d_{\text{HH}}((-1)^n(a_n \otimes a_0 \otimes \cdots \otimes a_{n-1}))\]

After accounting for Koszul signs from moving $a_n$ past $a_0, \ldots, a_{n-1}$,
the two expressions are identical.
\end{proof}

\begin{definition}[Connes operator B]
\label{def:connes-B}
\index{cyclic homology!Connes operator|textbf}
The \emph{Connes operator} $B: \mathrm{CH}_n(\mathcal{A}) \to \mathrm{CH}_{n+1}(\mathcal{A})$ is defined as
\[B = (1-t_{n+1}) \circ s \circ N_n\]
where $N_n = \sum_{i=0}^{n} t^i$ is the norm operator and $s: \mathrm{CH}_n \to \mathrm{CH}_{n+1}$ is the extra degeneracy $s(a_0 \otimes \cdots \otimes a_n) = 1 \otimes a_0 \otimes \cdots \otimes a_n$. The identities $B^2 = 0$ and $Bd + dB = 0$ hold (Theorem~\ref{thm:connes-exact-sequence}).
\end{definition}

\begin{theorem}[Connes mixed-complex structure {\cite{Connes85,Loday98}}; \ClaimStatusProvedElsewhere]
\label{thm:connes-exact-sequence}
The pair $(d_{\mathrm{HH}}, B)$ equips $\mathrm{CH}_*(\mathcal{A})$ with the
structure of a \emph{mixed complex}: $d^2 = B^2 = dB + Bd = 0$.
This gives rise to a short exact sequence of complexes:
\[0 \to \mathrm{CC}_*(\mathcal{A}) \to \mathrm{CC}_*^{\mathrm{per}}(\mathcal{A})
\to \mathrm{CC}_*^{\mathrm{neg}}(\mathcal{A})[2] \to 0\]
where:
\begin{itemize}
\item $\mathrm{CC}_* = \mathrm{CH}_* \otimes k[u]/(d + uB)$ is the cyclic complex
 ($|u| = -2$)
\item $\mathrm{CC}_*^{\mathrm{per}} = \mathrm{CH}_* \otimes k((u))/(d+uB)$ is the
 periodic cyclic complex
\item $\mathrm{CC}_*^{\mathrm{neg}}$ is the negative cyclic complex
\end{itemize}
\end{theorem}

\begin{remark}[Provenance and citation]
This theorem is imported and treated as \ClaimStatusProvedElsewhere. Standard
presentations of Connes' mixed-complex formalism and the associated exact
sequence are in \cite{Connes85,Loday98}.
\end{remark}

\begin{corollary}[Connes SBI exact sequence {\cite{Connes85,Loday98}}; \ClaimStatusProvedElsewhere]
\label{cor:connes-SBI}
For a chiral Hochschild chain model carrying the mixed-complex
structure of Theorem~\textup{\ref{thm:connes-exact-sequence}}, there
is a long exact sequence relating chiral Hochschild and cyclic homology:
\[\cdots \to \mathrm{HH}_n(\mathcal{A}) \xrightarrow{I} \mathrm{HC}_n(\mathcal{A}) \xrightarrow{S} \mathrm{HC}_{n-2}(\mathcal{A}) \xrightarrow{B} \mathrm{HH}_{n-1}(\mathcal{A}) \to \cdots\]
where $S$ is the periodicity operator, $B$ is Connes' operator, and $I$ is the natural inclusion. In particular, $S$ is an isomorphism in degrees where $\mathrm{HH}_*$ vanishes, yielding $2$-periodicity of periodic cyclic homology $\mathrm{HP}_* = \varprojlim_S \mathrm{HC}_{*+2k}$.
\end{corollary}

\subsection{The chiral Hochschild-cyclic spectral sequence}

\begin{theorem}[Chiral Hochschild-cyclic spectral sequence {\cite{Connes85,Loday98}}; \ClaimStatusProvedElsewhere]
\label{thm:HC-spectral-sequence}
\index{spectral sequence!Hochschild-cyclic}
Let
\[
C=(\mathrm{CH}^{\mathrm{ch}}_*(\mathcal A),d_{\mathrm{HH}},B)
\]
be the mixed complex of Theorem~\textup{\ref{thm:connes-exact-sequence}}.
The $u$-adic filtration on the negative cyclic complex
$C[[u]]$ and the induced filtration on the periodic complex $C((u))$
give spectral sequences whose first pages are
\[
E_1^{\mathrm{neg}}\cong \mathrm{HH}^{\mathrm{ch}}_*(\mathcal A)[[u]],
\qquad
E_1^{\mathrm{per}}\cong \mathrm{HH}^{\mathrm{ch}}_*(\mathcal A)((u)),
\qquad |u|=-2,
\]
with first differential induced by $uB$. They converge, under the
usual completeness and boundedness hypotheses for the chosen completed
chain model, to $\mathrm{HC}^{\mathrm{neg,ch}}_*(\mathcal A)$ and
$\mathrm{HP}^{\mathrm{ch}}_*(\mathcal A)$ respectively. The ordinary
cyclic spectral sequence is obtained from the quotient
$C((u))/uC[[u]]$.
\end{theorem}

\begin{remark}[Provenance and citation]
This theorem is imported and treated as \ClaimStatusProvedElsewhere. The
spectral sequence attached to the mixed complex $(\mathrm{CH}_*,d_{\mathrm{HH}},B)$
is standard in the cyclic-homology literature; see \cite{Connes85,Loday98}.
The theorem asserts the mixed-complex spectral sequence. Any two-column
display below records the first two $u$-columns or a zero-mode
truncation; the full cyclic spectral sequence keeps the complete
$u$-adic filtration.
\end{remark}

\subsection{\texorpdfstring{$E_2$ page: explicit computation}{E2 page: explicit computation}}

\begin{theorem}[Second-page formula {\cite{Loday98}}; \ClaimStatusProvedElsewhere]
\label{thm:E2-page-formula}
For the mixed complex
$C=(\mathrm{CH}^{\mathrm{ch}}_*(\mathcal A),d_{\mathrm{HH}},B)$, the
second page of the $u$-adic cyclic spectral sequence is
\[
E_2^{\mathrm{neg}}
\cong
H\!\left(\mathrm{HH}^{\mathrm{ch}}_*(\mathcal A)[[u]],\,uB\right),
\qquad
E_2^{\mathrm{per}}
\cong
H\!\left(\mathrm{HH}^{\mathrm{ch}}_*(\mathcal A)((u)),\,uB\right).
\]
The ordinary cyclic version is the corresponding homology for
$\mathrm{HH}^{\mathrm{ch}}_*(\mathcal A)((u))/u\mathrm{HH}^{\mathrm{ch}}_*(\mathcal A)[[u]]$.
When a locally constant trace model identifies the mixed complex with
an $S^1$-equivariant chain model, these pages are the algebraic
counterparts of the associated homotopy fixed-point, orbit, and Tate
spectral sequences. No invariant/coinvariant formula is used without
that trace comparison.
\end{theorem}

\begin{remark}[Provenance and citation]
This theorem is imported and treated as \ClaimStatusProvedElsewhere. It
is the standard second-page calculation for a mixed complex filtered by
the cyclic parameter $u$; see \cite{Loday98}.
\end{remark}

\begin{example}[Second page for the Heisenberg algebra]
\label{ex:E2-heisenberg}
For the Heisenberg chiral algebra $\mathcal{H}$ at level $k$:

\emph{Step 1: Compute chiral Hochschild homology}
\[\text{HH}_*(\mathcal{H}) = \mathbb{C}[a_{-n}]_{n \geq 1} \otimes \mathbb{C} \cdot \mathbf{1}\]
where $a_{-n}$ are negative modes of the Heisenberg field $a(z) = \sum_n a_n z^{-n-1}$.

In more invariant terms:
\[\text{HH}_*(\mathcal{H}) = \text{Sym}^*(H_1(S^1, \mathbb{C})) \otimes \Lambda(H_0(S^1, \mathbb{C})) = \mathbb{C}[c] \otimes \Lambda(\sigma)\]
where $|c| = 2$ (from the $S^1$ homology) and $|\sigma| = 1$ (from the constant functions), consistent with the Serre duality computation.

\emph{Step 2: $S^1$-action}

The $S^1$-action on $\text{HH}_*(\mathcal{H})$ is trivial on the center, so:
\[\text{HH}_*(\mathcal{H})^{S^1} = \mathbb{C}[c]\]
\[\text{HH}_*(\mathcal{H})_{S^1} = \mathbb{C}[c]\]

\emph{Step 3: first two $u$-columns of the $E_2$ page}
\[E_2^{0,q} = \mathbb{C}[c] \cap \{|*| = q\} = \begin{cases}
\mathbb{C} & \text{if } q = 0, 2, 4, \ldots \\
0 & \text{otherwise}
\end{cases}\]

\[E_2^{1,q} = \mathbb{C}[c] \cap \{|*| = q-1\} = \begin{cases}
\mathbb{C} & \text{if } q = 1, 3, 5, \ldots \\
0 & \text{otherwise}
\end{cases}\]

The $E_2$ page is therefore:
\[E_2 = \begin{array}{c|cc}
q \backslash p & 0 & 1 \\
\hline
0 & \mathbb{C} & 0 \\
1 & 0 & \mathbb{C} \\
2 & \mathbb{C} & 0 \\
3 & 0 & \mathbb{C} \\
\vdots & \vdots & \vdots
\end{array}\]

This display is the first-two-column truncation of the $u$-adic
mixed-complex spectral sequence. The full negative or periodic cyclic
spectral sequence has further $u$-columns, generated by multiplication
by the cyclic parameter $u$. Degeneration at $E_2$ is asserted only
after the chosen
locally constant zero-mode model has no further $uB$-homology
extension data.

\emph{Comparison with classical Hochschild cohomology.}
The polynomial $E_2$ page above should not be confused with the
ordinary classical Hochschild cohomology
$\mathrm{HH}^*(\mathcal{H})_{\mathrm{alg}}$, which is computed in
Example~\ref{ex:HH-heisenberg-complete}. The latter uses the Koszul
resolution (length~$2$) and gives
$\mathrm{HH}^n_{\mathrm{alg}}(\mathcal{H}) = \mathbb{C}$ for
$n = 0, 1, 2$ and
$\mathrm{HH}^n_{\mathrm{alg}}(\mathcal{H}) = 0$ for $n \geq 3$.
The difference is that the $E_2$ page here is computed from chiral
Hochschild \emph{homology}
$\mathrm{HH}^{\mathrm{ch}}_*(\mathcal{H})
= \mathbb{C}[c] \otimes \Lambda(\sigma)$, which is polynomial and
does not vanish in high degrees. Classical Hochschild cohomology
\textup{(}the $\mathrm{Ext}$-computation\textup{)} and chiral
Hochschild homology \textup{(}the $\mathrm{Tor}$-computation\textup{)}
are distinct invariants; the former detects deformations, while the
latter feeds the cyclic spectral sequence.
\end{example}

\begin{example}[Second page for affine sl3]
\label{ex:E2-sl3}
\index{$\mathfrak{sl}_3$!Hochschild-cyclic spectral sequence}
For $\widehat{\mathfrak{sl}}_3$ at generic level~$k$:

\emph{Step~1: Chiral Hochschild homology.}
By the continuous cohomology of current algebras, $\mathrm{HH}_*^{\mathrm{ch}}(\widehat{\mathfrak{sl}}_3)$ is determined by the Lie algebra cohomology $H^*(\mathfrak{sl}_3; \mathbb{C}) = \Lambda(x_3, x_5)$ (exterior algebra on generators of degrees $3$ and $5$, corresponding to the exponents $m_1 = 1$, $m_2 = 2$ of $\mathfrak{sl}_3$) tensored with the polynomial algebra from the loop direction. This gives:
\[\mathrm{HH}_*^{\mathrm{ch}}(\widehat{\mathfrak{sl}}_3)
= \mathbb{C}[\Theta_2, \Theta_3] \otimes \Lambda(\sigma_3, \sigma_5)\]
where $|\Theta_j| = j$ (from the Casimir operators of degrees $2$ and $3$, corresponding to the exponents $m_1 = 1$, $m_2 = 2$) and $|\sigma_{2i+1}| = 2i+1$ (from $H^*(\mathfrak{sl}_3)$).

\emph{Step~2: $S^1$-action and first two $u$-columns.}
The $S^1$-action on the polynomial part is trivial (Casimirs are invariant).
The exterior generators $\sigma_3, \sigma_5$ lie in odd degrees and are
$S^1$-coinvariants. In the first two $u$-columns:
\[E_2^{0,*} = \mathbb{C}[\Theta_2, \Theta_3], \qquad
E_2^{1,*} = \mathbb{C}[\Theta_2, \Theta_3].\]

\emph{Step~3: Truncated Hilbert function.}
At generic~$k$, the locally constant zero-mode truncation has no
additional $uB$-extension data after the displayed columns. The
Hilbert function of that truncation,
$\mathrm{HC}_n^{\mathrm{tr}} = E_2^{0,n} \oplus E_2^{1,n-1}$, gives:
\[\dim \mathrm{HC}_n^{\mathrm{ch}}(\widehat{\mathfrak{sl}}_3)
= \#\{(a,b) : 2a + 3b = n\} + \#\{(a,b) : 2a + 3b = n-1\}\]
which grows linearly in~$n$ (rank~$2$), in contrast to the constant
$\dim \mathrm{HC}_n = 1$ for $\widehat{\mathfrak{sl}}_2$ (rank~$1$).
\end{example}

\begin{example}[Second page for lattice VOAs]
\label{ex:E2-lattice}
\index{lattice vertex algebra!Hochschild-cyclic}
For a lattice VOA $V_\Lambda$ of rank~$d$, the underlying Heisenberg
algebra $\mathcal{H}^{\otimes d}$ gives:
\[\mathrm{HH}_*^{\mathrm{ch}}(V_\Lambda) = \mathbb{C}[c_1, \ldots, c_d]
\otimes \Lambda(\sigma_1, \ldots, \sigma_d)\]
with $|c_i| = 2$, $|\sigma_i| = 1$. The lattice vertex operators do not
contribute additional chiral Hochschild classes (they are inner automorphisms of the
Heisenberg algebra at generic level). Thus the first two $u$-columns
are
$E_2^{0,*} = \mathbb{C}[c_1, \ldots, c_d]$
and $E_2^{1,*} = \mathbb{C}[c_1, \ldots, c_d]$ in the generic
zero-mode model.
\end{example}

\subsection{Physical interpretation: anomalies and partition functions}

\begin{theorem}[Partition function as cyclic trace in the locally constant model {\cite{FT87,BZFN10}}; \ClaimStatusConditional]
\label{thm:partition-cyclic}
Let $\mathcal A$ be a chiral algebra whose factorisation theory on
$\Sigma_g$ is locally constant after the chosen topological
specialisation, and assume the trace over the module category is
defined and convergent. Then the cyclic factorisation trace computes
the chiral partition function:
\[
Z[\Sigma_g]
= \operatorname{Tr}_{\int_{\Sigma_g}\mathcal A}(q^{L_0})
= \operatorname{Tr}_{\mathrm{HC}^{\mathrm{ch}}_*(\mathcal A)}(q^{L_0}).
\]

The physical loop-order filtration is read from the chiral
Hochschild-cyclic spectral sequence by:
\begin{itemize}
\item $E_1$ = tree-level (genus 0 contribution)
\item $E_2$ = one-loop (genus 1 + quantum corrections)
\item $E_r$ ($r \geq 3$) = multi-loop corrections (genus $\geq 2$)
\end{itemize}
\end{theorem}

\begin{remark}[Provenance]
Feigin--Tsygan~\cite{FT87} identify cyclic homology with the cyclic
bar trace in the algebraic setting. Ben-Zvi--Francis--Nadler
\cite{BZFN10} identify factorisation homology on closed manifolds
with categorical traces under the hypotheses of their theorem. The
displayed formula uses these two inputs after the locally constant
specialisation; it is not a statement about spectral $\mathrm{THH}$
and not a consequence of Theorem~H.
\end{remark}

\subsection{Master computation table}

\begin{table}[H]
\centering
\caption{\texorpdfstring{First two $u$-columns of the $E_2$}{First two u-columns of the E2} page in standard zero-mode models}
\begin{tabular}{|l|c|c|c|}
\hline
\textbf{Chiral Algebra} & $\mathbf{E_2^{0,*}}$ & $\mathbf{E_2^{1,*}}$ &
\emph{Status} \\
\hline
Heisenberg $\mathcal{H}_k$ & $\mathbb{C}[c]$ (even deg) & $\mathbb{C}[c]$ (odd deg) &
zero-mode model \\
\hline
$\beta\gamma$ bosons & $\mathbb{C}[c](1 \oplus \beta_0\gamma_0)$ & $\mathbb{C}[c]$ &
conditional model \\
\hline
$bc$ ghosts & Similar to $\beta\gamma$ & Similar to $\beta\gamma$ & conditional model \\
\hline
Virasoro $\mathcal{V}_c$ (generic $c$) & $\mathbb{C}[\Theta]$ ($|\Theta| = 2$) & 0 &
generic model${}^\dagger$ \\
\hline
$\widehat{\mathfrak{sl}}_2$ level $k$ & $\mathbb{C}[c]$ & $\mathbb{C}[c]$ & generic model \\
\hline
$\widehat{\mathfrak{sl}}_3$ level $k$ & $\mathbb{C}[\Theta_2, \Theta_3]$ & $\mathbb{C}[\Theta_2, \Theta_3]$ & generic model \\
\hline
$W_3$ algebra & $\mathbb{C}[\Theta_2, \Theta_3]$ & $\mathbb{C}[\Theta_2]$ & Depends on $c$ \\
\hline
$W_N$ algebra ($N \geq 4$) & $\mathbb{C}[\Theta_2, \ldots, \Theta_N]$ & $\mathbb{C}[\Theta_2, \ldots, \Theta_{N-1}]$ & generic model \\
\hline
Lattice $V_\Lambda$ (rank $d$) & $\mathbb{C}[c]^{\otimes d}$ & $\mathbb{C}[c]^{\otimes d}$ & generic model \\
\hline
$Y(\mathfrak{sl}_2)^{\mathrm{ch}}$ & $\mathbb{C}[\Theta_2]$ & $\mathbb{C}$ & Conjectural \\
\hline
\end{tabular}

\smallskip\noindent
${}^\dagger$\,The $E_2^{0,*}$ column records chiral Hochschild
\emph{homology} $\mathrm{HH}_*^{\mathrm{ch}}(\cA)^{S^1}$, a Tor
calculation.  Chiral Hochschild \emph{cohomology}
$\ChirHoch^*(\cA)$ is the corresponding curve-level Ext calculation.
Their degree-zero comparison is
$\mathrm{HH}_0\cong\ChirHoch^0=Z(\cA)$; higher-degree comparisons
require a family-specific chain map.  Theorem~H supplies the Ext
support through $H_H(\cA;S)$.

The table records only the first two $u$-columns of the mixed-complex
spectral sequence in the stated zero-mode or generic model. The full
negative, periodic, and ordinary cyclic spectral sequences require a
separate $u$-adic convergence and collapse argument.
\end{table}

\subsection{Cyclic homology exchange under Koszul duality}

\begin{corollary}[Cyclic homology duality under a trace comparison; \ClaimStatusConditional]\label{cor:cyclic-homology-duality}
\index{cyclic homology!Koszul duality}
Let $(\mathcal{A}, \mathcal{A}^!)$ be a Koszul pair equipped with a
perfect degree-$2$ chain pairing between their chiral Hochschild
cochain complexes.  Assume in addition that the chain model
$\mathrm{CH}^{\mathrm{ch}}_*(\mathcal A)$ carries a non-degenerate
trace/Poincare pairing identifying chiral Hochschild chains with the
linear dual of chiral Hochschild cochains, compatibly with Connes'
operator. Then the Connes periodicity operator\index{Connes periodicity operator}
$S$ intertwines the Koszul duality on chiral Hochschild cochains with
a shifted duality on periodic cyclic homology:
\begin{equation}\label{eq:cyclic-duality}
\mathrm{HP}_n(\mathcal{A}) \;\cong\; \mathrm{HP}_{-n}(\mathcal{A}^!)^{\vee}
\end{equation}
for all $n$.  The perfect cochain pairing and the trace comparison are
the two chain-level hypotheses carrying this identification.
\end{corollary}

\begin{proof}
The trace comparison gives a chain-level quasi-isomorphism
\[
 \mathrm{CH}^{\mathrm{ch}}_*(\mathcal A)
 \simeq
 \ChirHoch^{2-*}(\mathcal A)^\vee\otimes\omega_X
\]
and the analogous comparison for $\mathcal A^!$.  Composing these
maps with the supplied perfect degree-$2$ cochain pairing gives a
duality of mixed complexes
between the two cyclic chain models, with Connes' operator carried to
the dual Connes operator by the compatibility hypothesis. The SBI
exact sequence is natural for morphisms of mixed complexes, and the
periodicity operator $S$ shifts degree by~$-2$ on both sides. Passing
to
$\mathrm{HP}_*=\varprojlim_S\mathrm{HC}_{*+2k}$ absorbs the
curve-level shift and gives~\eqref{eq:cyclic-duality}.
\end{proof}

\begin{remark}\label{rem:cyclic-koszul-virasoro}
For the Virasoro algebra, the bounded complex has support
$\{0,2,3\}$ by Theorem~\ref{thm:virasoro-hochschild}.  A cyclic
duality statement for the curve chart begins with bounded-to-chart
comparisons for $\operatorname{Vir}_c$ and its chosen partner,
together with a perfect chain pairing of specified degree.  These
data determine the exchange of degree-two deformation classes and the
image of the degree-three class.
\end{remark}

\begin{corollary}[Adjointness of cup-product operators;
\ClaimStatusConditional]
\label{cor:hochschild-cup-exchange}
\index{Hochschild cohomology!cup product exchange}
Let $(\mathcal A,\mathcal A^!)$ carry a perfect degree-$d$ chain
pairing
\[
 \langle-,-\rangle\colon
 C^\bullet_{\mathrm{ch}}(\mathcal A,\mathcal A)
 \otimes
 C^\bullet_{\mathrm{ch}}(\mathcal A^!,\mathcal A^!)
 \longrightarrow\omega_X[-d].
\]
Let $\Theta\in\ChirHoch^2(\mathcal A)$ and
$\Theta^\dagger\in\ChirHoch^2(\mathcal A^!)$.  Assume the pairing is
Frobenius-compatible and satisfies
\begin{equation}\label{eq:cup-product-exchange}
 \langle\Theta\smile x,y\rangle
 =(-1)^{|x|}\langle x,\Theta^\dagger\smile y\rangle.
\end{equation}
Then cup product by $\Theta$ and cup product by
$\Theta^\dagger$ are adjoint on cohomology with the displayed sign.
\end{corollary}

\begin{proof}
Equation~\eqref{eq:cup-product-exchange} descends through cocycles and
coboundaries because the pairing is a chain map.  Perfectness then
identifies the two multiplication operators as adjoints.
\end{proof}

\begin{remark}[Virasoro specialization]\label{rem:cup-exchange-virasoro}
The bounded Virasoro calculation supplies a one-dimensional
degree-two group.  Its curve-chart specialization requires the two
bounded-to-chart comparisons.  A Verdier exchange of the resulting
degree-two classes then requires the perfect chain pairing and
Frobenius compatibility appearing in
Corollary~\ref{cor:hochschild-cup-exchange}.
\end{remark}

\section{Derived center, categorical Hochschild cohomology, and the algebraic shadow}
\label{sec:derived-center-hochschild-cochains}

\index{derived center|textbf}
\index{Morita equivalence!Hochschild cohomology}
\index{compact generator|textbf}

The derived centre of a dg-category is the derived endomorphism
algebra of its identity functor.  Finite thick generation identifies
the category with $\mathrm{Perf}(A)$; compact localizing generation
identifies it with $\mathrm{RMod}_A$.  Under either Morita
equivalence, the derived centre is the Hochschild cochain
$E_2$-algebra $C^\bullet(A,A)$, whose cohomology is
$\mathrm{HH}^\bullet(A)$.  The present chapter works in chain
complexes over $k$.  For an $E_1$-ring spectrum the corresponding
spectral centre is $F_{A^e}(A,A)$, while
$\mathrm{THH}(A)=A\otimes_{A^e}A$ is the trace object.

\subsection{Compact generators and the Morita equivalence}

\begin{definition}[Thick and compact generators;
\ClaimStatusDefinitional]
\label{def:compact-generator}
Let $\mathcal C$ be a $k$-linear stable dg-category and let
$b\in\mathcal C$.
\begin{enumerate}[label=\textup{(\roman*)}]
\item If $\mathcal C$ is essentially small and idempotent complete,
then $b$ is a \emph{thick generator} when
\[
 H^0(\mathcal C)=\operatorname{thick}(b),
\]
where $\operatorname{thick}(b)$ is closed under finite sums, shifts,
cones, and retracts.
\item If $\mathcal C$ is presentable, then $b$ is a \emph{compact
generator} when $\RHom_{\mathcal C}(b,-)$ preserves small coproducts
and
\[
 \operatorname{Loc}(b)=\mathcal C,
\]
where $\operatorname{Loc}(b)$ is the stable subcategory generated by
$b$ under all colimits.  Equivalently,
$\RHom_{\mathcal C}(b,-)$ is conservative.
\end{enumerate}
In either case set $A_b:=\RHom_{\mathcal C}(b,b)$, with multiplication
given by composition.
\end{definition}

\begin{computation}[The free rank-one generator;
\ClaimStatusProvedHere]
\label{comp:morita-free-generator}
Let $R$ be a dg-algebra, let $\mathcal C=\mathrm{RMod}_R$, and take
$b=R_R$.  Then
\[
 A_b=\RHom_R(R,R)\simeq R,
 \qquad
 \RHom_R(R,M)\simeq M.
\]
Evaluation at $1$ identifies an endomorphism of $R_R$ with left
multiplication $L_r$; the identity
$L_rL_s=L_{rs}$ fixes the displayed algebra orientation.
The adjunction
$-\otimes_R^{\mathbf L}R\dashv\RHom_R(R,-)$ is the identity
adjunction on $\mathrm{RMod}_R$.  Its compact restriction is
$\mathrm{Perf}(R)$.  The general theorem is this calculation after a
compact generator has been chosen.
\end{computation}

\begin{theorem}[Derived Morita reconstruction from one generator;
\ClaimStatusProvedHere]
\label{prop:morita-equivalence-compact-gen}
\index{Morita equivalence|textbf}
\textup{[Type: Open, Cat, $0\leftrightarrow1$.  The small package
$H_{\mathrm{Mor}}^\omega$ is essentially small, pretriangulated, and
idempotent complete.  The presentable package $H_{\mathrm{Mor}}^L$ is
a stable presentable $k$-linear category tensored over $\Ch(k)$.
Each package includes its stated generator condition.]}
Let $A_b=\RHom_{\mathcal C}(b,b)$.
\begin{enumerate}[label=\textup{(\roman*)}]
\item If $\mathcal C$ is essentially small, pretriangulated, and
idempotent complete, and $b$ is a thick generator, then
\[
 \Psi_\omega\colon\mathcal C\xrightarrow{\ \sim\ }
 \mathrm{Perf}(A_b),
 \qquad c\longmapsto\RHom_{\mathcal C}(b,c),
\]
is a quasi-equivalence.
\item If $\mathcal C$ satisfies $H_{\mathrm{Mor}}^L$ and $b$ is a
compact generator, then
\[
 \Psi\colon\mathcal C\xrightarrow{\ \sim\ }
 \mathrm{RMod}_{A_b}
\]
is an equivalence of stable presentable $\infty$-categories.  It
restricts on compact objects to
$\mathcal C^\omega\simeq\mathrm{Perf}(A_b)$.
\end{enumerate}
\end{theorem}

\begin{proof}
For the small statement, composition gives a natural map
\[
 \RHom_{\mathcal C}(c,c')
\;\longrightarrow\;
 \RHom_{A_b}\bigl(\Psi_\omega(c),\Psi_\omega(c')\bigr).
\]
For $c=b$ this is the identity
$\RHom_{\mathcal C}(b,c')\simeq
\RHom_{A_b}(A_b,\Psi_\omega(c'))$.  The objects $c$ for which the
map is a quasi-isomorphism for every $c'$ form a thick subcategory.
They contain $b$, hence all of $\mathcal C$.  Thus $\Psi_\omega$ is
quasi-fully faithful.  Its image is thick, contains
$A_b=\Psi_\omega(b)$, and therefore equals
$\operatorname{thick}(A_b)=\mathrm{Perf}(A_b)$.

For the presentable statement, $b$ is a left $A_b$-module by
evaluation, and there is an adjunction
\[
 \Phi=-\otimes_{A_b}^{\mathbf L}b:
 \mathrm{RMod}_{A_b}\rightleftarrows\mathcal C:
 \Psi=\RHom_{\mathcal C}(b,-).
\]
The unit is an equivalence on the free module $A_b$.  The functor
$\Phi$ preserves colimits by construction.  The functor $\Psi$ is
exact and preserves coproducts because $b$ is compact; hence it
preserves cofibres and geometric realizations, and therefore all small
colimits.  The modules on which the unit is an
equivalence therefore form a localizing subcategory containing
$A_b$.  Since $A_b$ generates $\mathrm{RMod}_{A_b}$, the unit is an
equivalence on every module.

For $c\in\mathcal C$, let $K_c$ be the cofiber of the counit
$\Phi\Psi(c)\to c$.  Applying $\Psi$ and using the unit shows
$\Psi(K_c)\simeq0$.  Generation by $b$ gives $K_c\simeq0$, so the
counit is an equivalence.  Thus $\Phi\dashv\Psi$ is an equivalence.
Equivalences preserve compact objects, and the compact objects of
$\mathrm{RMod}_{A_b}$ are $\mathrm{Perf}(A_b)$.

This is the dg form of Keller~\cite[\S4]{Keller06}, the monogenic
stable-model theorem of Schwede--Shipley~\cite[Theorem~3.1.1]{SS03},
and Lurie~\cite[Theorem~7.1.2.1 and Remark~7.1.2.3]{HA}.
\end{proof}

\subsection{The derived center}

\begin{definition}[Small and continuous derived centers;
\ClaimStatusDefinitional]
\label{def:derived-center}
\index{derived center|textbf}
For an essentially small, pretriangulated $k$-linear dg-category
$\mathcal C$, its \emph{derived center} is
\[
Z_{\mathrm{der}}(\mathcal{C})
:= \RHom_{\mathrm{Fun}^{\mathrm{ex}}(\mathcal{C}, \mathcal{C})}
 (\mathrm{Id}_{\mathcal{C}},\, \mathrm{Id}_{\mathcal{C}}).
\]
For a stable presentable dg-category $\mathcal C$, its
\emph{continuous derived center} is instead
\[
 Z_{\mathrm{der}}^{L}(\mathcal C)
 :=\RHom_{\mathrm{Fun}^{L}(\mathcal C,\mathcal C)}
 (\mathrm{Id}_{\mathcal C},\mathrm{Id}_{\mathcal C}),
\]
where the functors preserve colimits.  The superscript records this
continuous functor category.  The category of arbitrary endofunctors
produces a second, generally larger center.
A degree-$n$ cocycle in $Z_{\mathrm{der}}(\mathcal{C})$ is a
natural transformation $\mathrm{Id}_{\mathcal{C}} \Rightarrow
\mathrm{Id}_{\mathcal{C}}[n]$; its cohomology class is a derived
natural transformation. The product on $Z_{\mathrm{der}}(\mathcal{C})$
is composition of natural transformations; in particular,
$H^0(Z_{\mathrm{der}}(\mathcal{C}))$ is the ordinary center of the
homotopy category $[\mathcal{C}]$.
\end{definition}

\subsection{Endofunctors as bimodules}

The bridge between the derived center and categorical Hochschild
cohomology is the identification of endofunctors with bimodules.

\begin{proposition}[Endofunctor--bimodule equivalence {\cite{Toen07,BZFN10}};
\ClaimStatusProvedElsewhere]
\label{prop:endofunctor-bimodule}
\index{bimodule!endofunctor equivalence}
For a dg-algebra~$A$, let $\mathrm{Mod}_A$ be the presentable
dg-category of right $A$-modules.  There is a quasi-equivalence of
dg-categories
\[
\mathrm{Fun}^{L}(\mathrm{Mod}_A,\mathrm{Mod}_A)
\;\simeq\; A\text{-}\mathrm{bimod},
\]
where the left-hand side consists of colimit-preserving dg-functors.
The equivalence sends an $A$-$A$-bimodule $M$ to
$- \otimes_A^{\mathbf{L}} M$.  It preserves compact objects precisely
when $M=F(A)$ is perfect as a \emph{right} $A$-module.  Restriction to
compact objects and Ind-extension therefore give an equivalence
\[
 \mathrm{Fun}^{\mathrm{ex}}\bigl(\mathrm{Perf}(A),
       \mathrm{Perf}(A)\bigr)
 \simeq
 \mathrm{Bimod}_{A\text{-}A}^{\mathrm{rperf}},
\]
where the right-hand side is the full dg-subcategory of bimodules
perfect as right $A$-modules.  Perfection over
$A^e=A\otimes A^{\mathrm{op}}$ is the additional smooth-kernel
condition.
\end{proposition}

\begin{proof}
The dg Eilenberg--Watts theorem says that a colimit-preserving
functor $F\colon \mathrm{Mod}_A\to\mathrm{Mod}_A$ is represented by
the bimodule $M_F:=F(A)$, where the right action comes from the free
right module and the left action is induced by functoriality in the
endomorphisms of~$A$.  The assignment $F\mapsto F(A)$ is quasi-inverse
to $M\mapsto(-\otimes_A^{\mathbf L}M)$.  Since $A$ compactly
generates $\mathrm{Mod}_A$, such a functor preserves compact objects
precisely when $F(A)=M$ is compact as a right $A$-module.  It then
restricts to an exact endofunctor of $\mathrm{Perf}(A)$.  Conversely,
Ind-extension takes an exact endofunctor of $\mathrm{Perf}(A)$ to a
colimit-preserving, compact-preserving endofunctor of
$\mathrm{Mod}_A$.  See
To\"en~\cite{Toen07} and Ben-Zvi--Francis--Nadler~\cite{BZFN10} for
the $\infty$-categorical statement.
\end{proof}

\begin{corollary}[Identity functor $=$ diagonal bimodule]
\label{cor:identity-diagonal}
\ClaimStatusProvedHere
Under the compact-preserving restriction of the equivalence of
Proposition~\textup{\ref{prop:endofunctor-bimodule}}, the identity
functor $\mathrm{Id}\colon \mathrm{Perf}(A) \to \mathrm{Perf}(A)$
corresponds to the diagonal bimodule ${}_A A_A$.
\end{corollary}

\begin{proof}
The identity functor sends the free module~$A$ to itself:
$\mathrm{Id}(A) = A$. The corresponding bimodule is therefore~$A$,
with left $A$-action by left multiplication and right $A$-action by
right multiplication. This is the diagonal bimodule.
\end{proof}

\subsection{The main identification}

\begin{theorem}[Derived center $=$ categorical Hochschild cohomology $=$ algebraic Hochschild cochains via a thick generator;
\ClaimStatusProvedHere]
\label{thm:derived-center-hochschild}
\index{derived center!Hochschild identification}
\index{Hochschild cohomology!derived center}
Let\/ $\mathcal{C}$ be an essentially small, pretriangulated,
idempotent-complete dg-category with thick generator $b$, and set
$A = \RHom_{\mathcal{C}}(b, b)$.  There is a quasi-isomorphism of
dg-algebras
\begin{equation}\label{eq:derived-center-HH}
Z_{\mathrm{der}}(\mathcal{C})
\;\simeq\;
\mathrm{HH}^\bullet_{\mathrm{cat}}(\mathcal{C})
\;\simeq\;
\mathrm{HH}^\bullet_{\mathrm{alg}}(A)
:= \RHom_{A^e}(A, A).
\end{equation}
In particular, the derived center is a Morita invariant of
$\mathcal C$.
\end{theorem}

\begin{proof}
The proof assembles the three preceding results.

\emph{Step~1} (Morita reduction).
By Theorem~\ref{prop:morita-equivalence-compact-gen}, $b$ being a
thick generator gives a quasi-equivalence
$\Psi\colon \mathcal{C} \xrightarrow{\;\sim\;} \mathrm{Perf}(A)$.
Passing to the derived dg-functor categories carries the identity of
$\mathcal C$ to the identity of $\mathrm{Perf}(A)$ and induces a
quasi-isomorphism on their derived endomorphism complexes.  Hence
\[
Z_{\mathrm{der}}(\mathcal{C})
\;\simeq\;
Z_{\mathrm{der}}(\mathrm{Perf}(A))
= \RHom_{\mathrm{Fun}^{\mathrm{ex}}(\mathrm{Perf}(A), \mathrm{Perf}(A))}
 (\mathrm{Id},\, \mathrm{Id}).
\]

\emph{Step~2} (Endofunctors $=$ bimodules).
By Proposition~\ref{prop:endofunctor-bimodule}, exact endofunctors of
$\mathrm{Perf}(A)$ identify with bimodules perfect as right
$A$-modules.  This is a full dg-subcategory of all $A$-bimodules, and
the diagonal bimodule $A$ belongs to it because it is free of rank one
as a right module.  Applying $\RHom$ to the identity endofunctor
therefore gives
\[
\RHom_{\mathrm{Fun}}(\mathrm{Id},\, \mathrm{Id})
\;\simeq\;
\RHom_{A\text{-}\mathrm{bimod}}
 (M_{\mathrm{Id}},\, M_{\mathrm{Id}}).
\]

\emph{Step~3} (Identity $\mapsto$ diagonal).
By Corollary~\ref{cor:identity-diagonal}, $M_{\mathrm{Id}} = {}_A A_A$
is the diagonal bimodule. Therefore:
\[
Z_{\mathrm{der}}(\mathrm{Perf}(A))
\;\simeq\;
\RHom_{A^e}(A, A)
= \mathrm{HH}^\bullet_{\mathrm{alg}}(A).
\]

\emph{Functoriality.}  A dg-functor equipped with a compatible map of
chosen generators determines the comparison chain map.  An unpointed
Morita class determines the corresponding equivalence in the homotopy
category, as recorded below.
\end{proof}

\begin{corollary}[Continuous center from a compact generator;
\ClaimStatusProvedHere]
\label{cor:continuous-center-compact-generator}
Let $\mathcal C$ be a stable presentable $k$-linear dg-category with
compact generator $b$, and put $A=\RHom_{\mathcal C}(b,b)$.  Then
\[
 Z_{\mathrm{der}}^{L}(\mathcal C)
 \simeq
 \RHom_{A^e}(A,A).
\]
Under $\mathcal C^\omega\simeq\mathrm{Perf}(A)$ this is the derived
center of the compact subcategory.
\end{corollary}

\begin{proof}
Theorem~\ref{prop:morita-equivalence-compact-gen} identifies
$\mathcal C$ with $\mathrm{Mod}_A$.  Proposition~
\ref{prop:endofunctor-bimodule} identifies its colimit-preserving
endofunctors with $A$-bimodules and the identity with the diagonal
bimodule $A$.  Taking derived endomorphisms gives the formula.
Restriction to compact objects gives the final assertion.
\end{proof}

\begin{remark}[Provenance]
\label{rem:derived-center-provenance}
The identification
$Z_{\mathrm{der}}(\mathcal{C}) \simeq
\mathrm{HH}^\bullet_{\mathrm{alg}}(A)$, equivalently
$\mathrm{HH}^\bullet_{\mathrm{cat}}(\mathcal{C})$ after Morita
reduction, is implicit in Keller~\cite{Keller06} and
made explicit in the framework of $(\infty,1)$-categories by
Ben-Zvi--Francis--Nadler~\cite{BZFN10} and
To\"en~\cite{Toen07}. The argument presented here isolates the
essential chain of identifications.  The full
$(\infty,2)$-categorical formalism supplies its coherent enhancement.
For the factorization-algebra
refinement relevant to chiral algebras, see
Costello--Gwilliam~\cite[Chapter~5]{CG17}.
\end{remark}

\subsection{Independence of generator: Morita invariance of
algebraic Hochschild cohomology}

\begin{theorem}[Morita invariance of algebraic Hochschild cohomology;
\ClaimStatusProvedHere]
\label{thm:morita-invariance-HH}
\index{Hochschild cohomology!Morita invariance|textbf}
\index{Morita equivalence!Hochschild invariance}
Let $\mathcal C$ be an essentially small, pretriangulated,
idempotent-complete dg-category with thick generators $b,b'$, and put
\[
 A=\RHom_{\mathcal C}(b,b),
 \qquad
 A'=\RHom_{\mathcal C}(b',b').
\]
Then the common Morita context determines an equivalence of
Hochschild-cochain $E_2$-algebras
\[
 C^\bullet(A,A)\simeq C^\bullet(A',A').
\]
In particular,
$\mathrm{HH}^\bullet_{\mathrm{alg}}(A)\cong
\mathrm{HH}^\bullet_{\mathrm{alg}}(A')$ as graded algebras.  The same
statement holds when $\mathcal C$ is stable presentable and $b,b'$ are
compact generators, by applying the small statement to
$\mathcal C^\omega$.
\end{theorem}

\begin{proof}
In the small case Theorem~\ref{thm:derived-center-hochschild}
identifies both Hochschild-cochain objects with the derived
endomorphisms of the identity functor of $\mathcal C$.  These
identifications preserve the cup product and the brace operations,
hence the $E_2$-structure.  In the presentable case,
Theorem~\ref{prop:morita-equivalence-compact-gen} identifies
$\mathcal C^\omega$ with both $\mathrm{Perf}(A)$ and
$\mathrm{Perf}(A')$.  The explicit Morita context is recorded next.
\end{proof}

\begin{proposition}[Bimodule realization of Morita transfer;
\ClaimStatusProvedHere]
\label{prop:explicit-morita-transfer}
\index{Morita equivalence!explicit transfer map}
With the notation of Theorem~\ref{thm:morita-invariance-HH}, set
\[
 {}_{A'}P_A:=\RHom_{\mathcal C}(b,b'),
 \qquad
 {}_A Q_{A'}:=\RHom_{\mathcal C}(b',b).
\]
Composition supplies equivalences of bimodules
\[
 Q\otimes_{A'}^{\mathbf L}P\simeq A,
 \qquad
 P\otimes_A^{\mathbf L}Q\simeq A'.
\]
The functor
\[
 T\colon D(A^e)\longrightarrow D((A')^e),
 \qquad
 M\longmapsto P\otimes_A^{\mathbf L}M\otimes_A^{\mathbf L}Q,
\]
is an equivalence, sends the monoidal unit $A$ to $A'$, and induces
the Morita equivalence of Hochschild cochains.  On a derived
bimodule morphism $f\colon A\to A[n]$ it is represented by
\begin{equation}\label{eq:morita-transfer-map}
 A'\simeq P\otimes_A^{\mathbf L}Q
 \xrightarrow{\;\mathrm{id}_P\otimes f\otimes\mathrm{id}_Q\;}
 P\otimes_A^{\mathbf L}A[n]\otimes_A^{\mathbf L}Q
 \simeq A'[n].
\end{equation}
A choice of cofibrant bimodule models realizes this equivalence by a
zigzag of $E_2$ quasi-isomorphisms.  Coherent choices of models,
evaluation, and coevaluation select a strict representative.
\end{proposition}

\begin{proof}
The actions displayed in the subscripts follow from postcomposition
by $A'$ and precomposition by $A$ on $P$, and conversely on $Q$.
The two evaluation maps are equivalences because $b$ and $b'$ are
thick generators.  The functor
\[
 S(N)=Q\otimes_{A'}^{\mathbf L}N\otimes_{A'}^{\mathbf L}P
\]
is inverse to $T$: the composites contract the adjacent factors
$Q\otimes_{A'}^{\mathbf L}P$ and
$P\otimes_A^{\mathbf L}Q$ to $A$ and $A'$, respectively.  The same
contractions give coherent equivalences
\[
 T(M\otimes_A^{\mathbf L}N)
 \simeq
 T(M)\otimes_{A'}^{\mathbf L}T(N),
 \qquad T(A)\simeq A'.
\]
Thus $T$ transports the derived endomorphism object of the monoidal
unit $A$ to that of $A'$, including composition and braces.  Formula
\eqref{eq:morita-transfer-map} is the resulting action on a morphism.
\end{proof}

\begin{remark}[The chiral comparison is additional structure;
\ClaimStatusConditional]
\label{rem:derived-center-chiral}
\index{derived center!chiral application}
For a chiral algebra $\cA$ on a curve~$X$, first fix a stable
presentable factorization-module category
$\mathrm{Mod}^{\mathrm{fact}}_\cA$ and its completed mapping
complexes.  Its categorical chiral center is, by definition,
\[
 Z^{\mathrm{der}}_{\mathrm{ch}}(\cA)
 :=Z^{L}_{\mathrm{der}}(\mathrm{Mod}^{\mathrm{fact}}_\cA).
\]
If this category has a compact generator $b$, Morita reconstruction
gives
\[
 Z^{\mathrm{der}}_{\mathrm{ch}}(\cA)
 \simeq
 \RHom_{R_b^e}(R_b,R_b),
 \qquad
 R_b:=\RHom_{\mathrm{Mod}^{\mathrm{fact}}_\cA}(b,b).
\]
Identification with a chosen chiral cochain complex is the content of
a separate comparison quasi-isomorphism
\[
 \chi_\cA\colon
 \RHom_{R_b^e}(R_b,R_b)
 \xrightarrow{\ \sim\ }
 \ChirHoch^\bullet(\cA,\cA).
\]
The construction of $\chi_\cA$ combines factorization Morita theory
with a comparison between two specified cochain models.  The symbol
$\ChirHoch^\bullet$ therefore carries one of two type signatures: the
Beilinson--Drinfeld chiral-operation complex, or the Hochschild
complex of a specified associative $E_1$ factorization object.  A
comparison theorem supplies an identification between them.

There is a useful first-principles test for the missing generation
hypothesis.  Suppose
$F\colon\mathcal D\rightleftarrows
\mathrm{Mod}^{\mathrm{fact}}_\cA:U$ is an adjunction, $U$ is
conservative and preserves coproducts, and $d$ is a compact generator
of $\mathcal D$.  Then $F(d)$ is a compact generator, because
\[
 \RHom(F(d),M)\simeq\RHom(d,U(M)).
\]
Compactness follows from preservation of coproducts, and vanishing of
the left side forces $U(M)$, hence $M$, to vanish.  Thus a claimed
generator for a Heisenberg, affine, Virasoro, or $\mathcal W$ family
must be accompanied by such a monadic or equivalent construction.
$C_2$-cofiniteness, a free-field resolution, and BRST reduction supply
finiteness and resolution inputs; the displayed adjunction supplies
the compact-generation theorem.

After $\chi_\cA$ has been constructed, a class in
$\ChirHoch^2(\cA,\cA)$ corresponds to a deformation of the
identity functor.  The bar cohomology $H^*(B(\cA))$ is a distinct
Koszul-dual coalgebra invariant.  A field-theoretic realization
promotes the relevant chain object to a physical bulk state space.
\end{remark}

\section{Factorization homology on the annulus: algebraic Hochschild chains and the circle trace}
\label{sec:annulus-hochschild-chains}

\index{factorization homology!annulus|textbf}
\index{algebraic Hochschild homology!circle trace and excision|textbf}
\index{excision!factorization homology|textbf}

The complementary side of the categorical/algebraic Hochschild
picture is homological: while the derived center computes
categorical Hochschild cochains, the circle trace computes the
categorical Hochschild homology of the boundary category and, in a
compact-generator chart, the algebraic Hochschild chains of the
endomorphism dg-algebra. The input here is a locally constant
factorization algebra with values in $\Ch(k)$, so the output stays
in chain complexes. Spectral $\mathrm{THH}$ appears only after
spectrifying the interval algebra by an Eilenberg--Mac Lane
construction. The excision axiom for factorization homology
reduces the circle computation to a tensor product over an interval,
yielding the identification of factorization homology on $S^1$
with algebraic Hochschild chains.

\subsection{Factorization homology: recollections}

\begin{definition}[Factorization homology over a manifold]
\label{def:factorization-homology-manifold}
\index{factorization homology|textbf}
For a locally constant factorization algebra $\mathcal{F}$ on a
manifold~$M$ with values in a symmetric monoidal
$\infty$-category~$\mathcal{V}$, the \emph{factorization homology}
of $\mathcal{F}$ over~$M$ is the global sections:
\[
\int_M \mathcal{F}
:= \mathcal{F}(M)
\;\in\; \mathcal{V}.
\]
When $\mathcal{F}$ is constructible rather than locally constant,
this is the homotopy colimit of $\mathcal{F}$ over the
Weiss topology on~$M$
\textup{(}Ayala--Francis~\cite{AF15},
Costello--Gwilliam~\cite{CG17}\textup{)}.
\end{definition}

The factorization homology functor satisfies two fundamental
properties that make it computable: \emph{local constancy} (the
value on a disk is determined by the $\En$-algebra structure) and
\emph{excision} (the value on a glued manifold is computed by a
derived tensor product).

\begin{theorem}[Excision; {\cite[Theorem~3.18]{AF15}}; \ClaimStatusProvedElsewhere]
\label{thm:excision}
\index{excision!for factorization homology}
Let $M = M_+ \cup_N M_-$ be a collar-gluing decomposition: $N$ is a
codimension-$1$ manifold with a collar neighborhood in~$M$, and
$M_\pm$ are the two halves. Then there is a natural equivalence
\[
\int_M \mathcal{F}
\;\simeq\;
\int_{M_+} \mathcal{F}
\;\otimes^{\mathbf{L}}_{\int_N \mathcal{F}}\;
\int_{M_-} \mathcal{F}.
\]
\end{theorem}

\subsection{The circle computation by excision}

\begin{theorem}[Factorization homology on $S^1$ $=$ algebraic Hochschild chains;
\ClaimStatusProvedHere]
\label{thm:circle-fh-hochschild}
\index{factorization homology!circle}
\index{algebraic Hochschild homology!factorization homology}
Let $\mathcal{F}$ be a locally constant factorization algebra on
$S^1$ with values in $\Ch(k)$, and let $\mathcal{C}$ be the
associated $\Eone$-category \textup{(}the dg-category of modules
over the $\Eone$-algebra $A = \mathcal{F}(I)$ for any interval
$I \subset S^1$\textup{)}. Then there is a natural equivalence
\begin{equation}\label{eq:circle-fh-HH}
\int_{S^1} \mathcal{F}
\;\simeq\;
\mathrm{HH}_{\bullet,\mathrm{cat}}(\mathcal{C})
:= \mathrm{HH}_{\bullet,\mathrm{alg}}(A)
= A \otimes^{\mathbf{L}}_{A^e} A,
\end{equation}
where $A^e = A \otimes A^{\mathrm{op}}$.
\end{theorem}

\begin{proof}
The proof proceeds by cutting $S^1$ into overlapping arcs and
applying excision.

\emph{Step~1} (Collar-gluing decomposition of $S^1$).
Choose a decomposition $S^1 = U_+ \cup U_-$ where $U_+$ and $U_-$
are open arcs (each diffeomorphic to an interval) whose intersection
consists of two disjoint intervals:
\[
U_+ \cap U_- = J_1 \sqcup J_2.
\]
Concretely, parametrize $S^1 = \R / \Z$ and take
$U_+ = (0, 3/4)$, $U_- = (1/2, 5/4)$, so that
$J_1 = (1/2, 3/4)$ and $J_2 = (0, 1/4)$ (the latter wrapping around
$0 \equiv 1$). The collar structure along the boundary is given by
the intervals $J_1$ and $J_2$ themselves.

\emph{Step~2} (Local constancy on arcs).
Since $\mathcal{F}$ is locally constant, its value on any interval is
equivalent to the $\Eone$-algebra $A$:
\[
\mathcal{F}(U_+) \simeq A, \quad
\mathcal{F}(U_-) \simeq A, \quad
\mathcal{F}(J_1) \simeq A, \quad
\mathcal{F}(J_2) \simeq A.
\]
More precisely, for each interval $J \subset S^1$, the restriction
$\mathcal{F}|_J$ is equivalent to the factorization algebra associated
to a single $\Eone$-algebra~$A$, and the value $\mathcal{F}(J) \simeq A$
is given by the $\Eone$-structure.

\emph{Step~3} (Excision).
By the excision axiom (Axiom~\ref{thm:excision}) applied to the
collar-gluing $S^1 = U_+ \cup_{J_1 \sqcup J_2} U_-$:
\begin{equation}\label{eq:excision-circle}
\int_{S^1} \mathcal{F}
\;\simeq\;
\mathcal{F}(U_+)
\;\otimes^{\mathbf{L}}_{\mathcal{F}(J_1) \otimes \mathcal{F}(J_2)}\;
\mathcal{F}(U_-).
\end{equation}
Substituting the local constancy identifications from Step~2:
\begin{equation}\label{eq:excision-substituted}
\int_{S^1} \mathcal{F}
\;\simeq\;
A \;\otimes^{\mathbf{L}}_{A^e}\; A,
\qquad A^e := A \otimes A^{\mathrm{op}}.
\end{equation}

\emph{Step~4} (Bimodule structure and identification with HH).
The key is to identify the $A \otimes A$-module structures on the two
copies of~$A$ in
$A \otimes^{\mathbf{L}}_{A^e} A$. The inclusions
$J_1 \hookrightarrow U_+$ and $J_2 \hookrightarrow U_+$ give the
left and right $A$-actions on $\mathcal{F}(U_+) \simeq A$: the
inclusion of $J_1$ at one end of the arc~$U_+$ provides the left
action, and the inclusion of $J_2$ at the other end provides the
right action. That these are the \emph{regular} left and right
actions (i.e., $a \cdot_L x = ax$ and $x \cdot_R b = xb$) follows
from local constancy: the factorization algebra~$\mathcal{F}$ on the
arc~$U_+$ is equivalent to the $\Eone$-algebra~$A$ (Step~2), and
the collar inclusions $J_i \hookrightarrow U_+$ restrict to the
identity on~$A$ (since all intervals carry the same algebra by local
constancy). The bimodule structure is therefore
$(a_1, a_2) \cdot x = a_1 \cdot x \cdot a_2$, which is precisely the
diagonal bimodule ${}_A A_A$. Similarly for $\mathcal{F}(U_-)$, where
the \emph{opposite} orientation of $U_-$ relative to $U_+$ introduces
the $A^{\mathrm{op}}$ factor: the left action from $J_1$ on $U_-$
is the right action on $U_+$, and vice versa. This is the source
of the $A^{\mathrm{op}}$ in
$A^e = A \otimes A^{\mathrm{op}}$.

The tensor product over $A \otimes A$ of two copies of the diagonal
bimodule is, by definition, the algebraic Hochschild chain complex:
\[
A \otimes^{\mathbf{L}}_{A \otimes A^{\mathrm{op}}} A
= A \otimes^{\mathbf{L}}_{A^e} A
=: \mathrm{HH}_{\bullet,\mathrm{alg}}(A).
\]
Here $A \otimes A$ acts on the left copy of~$A$
via $(a_1, a_2) \cdot x = a_1 \cdot x \cdot a_2$
(using the bimodule structure from $J_1$ and $J_2$), and on the
right copy of~$A$ symmetrically; the tensor product over this action
contracts the two bimodules into the cyclic bar complex
$\cdots \to A^{\otimes 3} \to A^{\otimes 2} \to A$ computing
$\mathrm{HH}_{\bullet,\mathrm{alg}}(A)$.

\emph{Step~5} (Passage to the categorical statement).
If $\mathcal{C} = \mathrm{Perf}(A)$ is the dg-category of perfect
$A$-modules, then
$\mathrm{HH}_{\bullet,\mathrm{cat}}(\mathcal{C})$ is defined as
$\mathrm{HH}_{\bullet,\mathrm{alg}}(A)$ via the Morita-invariant
cyclic bar construction. Combining with Step~4 gives the stated
equivalence
$\int_{S^1} \mathcal{F} \simeq
\mathrm{HH}_{\bullet,\mathrm{cat}}(\mathcal{C})$.
\end{proof}

\begin{remark}[Provenance]
\label{rem:annulus-provenance}
The identification of factorization homology on $S^1$ with the
circle-trace/Hochschild construction is due (in the spectral
topological setting, where it computes $\mathrm{THH}$) to
Ben-Zvi--Francis--Nadler~\cite{BZFN10} and
Ayala--Francis~\cite{AF15}; in the vertex algebra setting, the
parallel result appears in Costello--Gwilliam~\cite[Chapter~5]{CG17}.
The chain-complex formula recorded here is the algebraic/categorical
shadow of that spectral statement. The excision argument above
follows the treatment in~\cite{AF15}.
\end{remark}

\subsection{The monodromy subtlety: from chiral to locally constant}
\label{subsec:monodromy-subtlety}

\index{monodromy!factorization homology}

The excision argument of Theorem~\ref{thm:circle-fh-hochschild}
requires $\mathcal{F}$ to be a \emph{locally constant} factorization
algebra on $S^1$. In the chiral algebraic setting, the situation is
more delicate: a chiral algebra~$\cA$ on a curve~$X$ is a
factorization algebra on $X$ that is locally constant on small disks
(the $\Eone$-structure on intervals), but the passage from an interval
to the full circle involves identifying the endpoints. This
identification may carry \emph{monodromy}, a nontrivial automorphism
of the $\Eone$-algebra $A = \cA(I)$ arising from parallel transport
around the puncture.

\begin{definition}[Monodromy automorphism]\label{def:monodromy-aut}
\index{monodromy!automorphism|textbf}
Let $\cA$ be a chiral algebra on a punctured disk $D^\times = D
\setminus \{0\}$, and let $I \subset D^\times$ be an interval.
The fundamental group $\pi_1(D^\times) \cong \Z$ acts on the
fiber $A = \cA(I)$ by monodromy. The \emph{monodromy
automorphism} is the generator:
\[
\mathfrak{m}\colon A \;\xrightarrow{\;\sim\;}\; A,
\]
obtained by analytically continuing sections of $\cA$ around the
puncture once counterclockwise.
\end{definition}

\begin{proposition}[Monodromy for standard families;
\ClaimStatusProvedHere]
\label{prop:monodromy-standard}
\begin{enumerate}[label=\textup{(\alph*)}]
\item \emph{Heisenberg $\mathcal{H}_k$.} The monodromy
 automorphism $\mathfrak{m}$ acts on the mode algebra generated
 by $\{J_n\}_{n \in \Z}$ via $\mathfrak{m}(J_n) = J_n$ for all
 $n$. The monodromy is \emph{trivial} on the algebra. However,
 on Fock modules $F_\lambda$, the monodromy acts by the scalar
 $e^{2\pi i \lambda^2 / (2k)}$
 \textup{(}from the conformal weight $h_\lambda = \lambda^2/(2k)$
 contribution to $q^{L_0}$\textup{)}. For the vacuum module
 $F_0$, this is trivial.
\item \emph{Affine $\widehat{\fg}_k$.} The monodromy is trivial on
 the algebra \textup{(}the current modes $J^a_n$ are
 single-valued\textup{)}. On highest-weight modules $V_\Lambda$ at
 integer level, the monodromy acts by
 $e^{2\pi i h_\Lambda}$ where $h_\Lambda = (\Lambda, \Lambda +
 2\rho) / 2(k + h^\vee)$ is the conformal weight.
\item \emph{Lattice VOA $V_\Lambda$.} The vertex operators
 $e^\alpha$ for $\alpha \in \Lambda$ have monodromy
 $\mathfrak{m}(e^\alpha) = e^{2\pi i (\alpha, \alpha)/2}\, e^\alpha$.
 For an even lattice, $(\alpha, \alpha) \in 2\Z$, so the monodromy
 is trivial. For an odd lattice, the monodromy is a sign
 $(-1)^{(\alpha,\alpha)}$ on each sector (note: odd lattices
 define \emph{super} vertex algebras; the sign is the
 fermion parity).
\end{enumerate}
\end{proposition}

\begin{proof}
In each case, the monodromy is computed from the $L_0$-eigenvalue.
For a primary field $\phi$ of conformal weight~$h$, the operator
product expansion $\phi(e^{2\pi i} z) = e^{2\pi i h} \phi(z)$ gives
the monodromy phase.

(a)~For $\mathcal{H}_k$: the current $J(z) = \sum_n J_n z^{-n-1}$
has conformal weight~$1$, so $J(e^{2\pi i}z) = e^{2\pi i} \cdot
J(z) \cdot e^{-2\pi i} = J(z)$ (the $z^{-n-1}$ factor absorbs the
phase). More precisely, $\mathfrak{m}(J_n) = J_n$ because the modes
are defined as contour integrals $J_n = \oint z^n J(z)\,dz$ over
cycles that are invariant under the monodromy deck transformation.
On $F_\lambda$, the highest-weight vector $|\lambda\rangle$ has
$L_0$-eigenvalue $h_\lambda = \lambda^2/(2k)$, giving the stated
phase.

(b)~Identical argument for $\widehat{\fg}_k$, with $h_\Lambda$
replacing $h_\lambda$.

(c)~For $V_\Lambda$, $e^\alpha$ has conformal weight
$(\alpha,\alpha)/2$, so $\mathfrak{m}(e^\alpha) = e^{2\pi i
(\alpha,\alpha)/2}\, e^\alpha$. For an even lattice,
$(\alpha,\alpha)/2 \in \Z$, so the phase is~$1$.
\end{proof}

\begin{remark}[Twisted algebraic Hochschild homology]
\label{rem:twisted-hochschild}
\index{algebraic Hochschild homology!twisted}
When the monodromy $\mathfrak{m}\colon A \to A$ is nontrivial, the
factorization homology on $S^1$ computes the
\emph{$\mathfrak{m}$-twisted algebraic Hochschild homology}:
\begin{equation}\label{eq:twisted-HH}
\int_{S^1} \mathcal{F}
\;\simeq\;
\mathrm{HH}_{\bullet,\mathrm{alg}}(A;\, {}_{\mathfrak{m}}A)
:= A \otimes^{\mathbf{L}}_{A^e} {}_{\mathfrak{m}}A,
\end{equation}
where ${}_{\mathfrak{m}}A$ is $A$ as a right $A$-module but with
left action twisted by $\mathfrak{m}$: $a \cdot_{\mathfrak{m}} x :=
\mathfrak{m}(a) \cdot x$. The excision argument of
Theorem~\ref{thm:circle-fh-hochschild} goes through verbatim with
$U_-$ contributing the twisted bimodule structure: the analytic
continuation around the full circle inserts $\mathfrak{m}$ into one
of the two collar inclusions.

For the standard families of
Proposition~\ref{prop:monodromy-standard}, the monodromy is trivial
on the algebra itself (the algebra modes are single-valued), so the
untwisted identification~\eqref{eq:circle-fh-HH} holds. The twist
becomes visible only at the module level: the torus partition function
$Z_{T^2}(\cA) = \mathrm{Tr}_{\mathcal{H}}(q^{L_0 - c/24})$ traces
over the Hilbert space~$\mathcal{H}$ with the insertion of
$q^{L_0}$, which is precisely the twist by the monodromy of the
cylinder.
\end{remark}

\begin{remark}[The $(g{=}1, k{=}0)$ modular MC component]
\label{rem:mc-hh-identification}
\index{Maurer--Cartan!genus-1 and algebraic Hochschild trace}
The $(g=1, n=0)$ component of the modular MC equation
(Chapter~\ref{ch:en-koszul-duality}) is:
\[
d\Theta_{\cA}^{(1,0)} + \tfrac{1}{2}
 [\Theta_{\cA}^{(0)}, \Theta_{\cA}^{(0)}]^{(1,0)}
 + \Delta_{\mathrm{sew}}\Theta_{\cA}^{(0,2)} = 0.
\]
The sewing operator $\Delta_{\mathrm{sew}}$ contracts two external
legs of $\Theta^{(0,2)}_{\cA}$ using the propagator on the torus.
This is the genus-$1$ algebraic Hochschild differential: the annulus
amplitude $\Theta^{(0,2)}_{\cA}$ evaluated on two inputs and sewn
into a torus is the trace $\mathrm{Tr}(m_2(-,-))$, the standard
generator of $\mathrm{HH}_{0,\mathrm{alg}}(A)$. The identification
of the locally constant interval model
$\int_{S^1}\mathcal F_A$ with
$\mathrm{HH}_{\bullet,\mathrm{alg}}(A)$
(Theorem~\ref{thm:circle-fh-hochschild})
is therefore the $n=0$ case of the general MC$\to$sewing
correspondence. It is not a $\cD$-module restriction of the chiral
algebra~$\cA$ to a real circle.
\end{remark}

\subsection{Heisenberg verification}

\begin{example}[Algebraic Hochschild chains of the Heisenberg interval algebra;
\ClaimStatusProvedHere]
\label{ex:heisenberg-HH-chains}
\index{Heisenberg algebra!algebraic Hochschild chains}
We verify the identification
$\int_{S^1} \mathcal F_{\mathcal{H}_k}^{\mathrm{lc}}
\simeq \mathrm{HH}_{\bullet,\mathrm{alg}}(A_{\mathcal{H}_k})$
against the known genus-$1$ partition function, where
$\mathcal F_{\mathcal{H}_k}^{\mathrm{lc}}$ is the locally constant
interval model extracted from the Heisenberg chiral algebra and
$A_{\mathcal{H}_k}$ is its completed interval algebra.
Here $\mathrm{HH}_{\bullet,\mathrm{alg}}(\mathcal{H}_k)$ is
shorthand for the algebraic Hochschild homology of the interval
algebra extracted from the locally constant Heisenberg
factorization algebra.

The Heisenberg vertex algebra $\mathcal{H}_k$ at level $k \neq 0$
is generated by a single current
$J(z)=\sum_{n\in\mathbb Z}J_nz^{-n-1}$ with OPE
$J(z)J(w)\sim k/(z-w)^2$. The completed interval algebra is the
oscillator algebra with
$[J_m,J_n]=km\,\delta_{m+n,0}$, completed in the Fock topology.
This completion is part of the factorisation-algebra input. It is not
the ordinary algebraic Weyl algebra with its uncompleted bar complex.

The chiral zero-mode calculation of Example~\ref{ex:E2-heisenberg}
gives
\[
\mathrm{HH}_\bullet^{\mathrm{ch}}(\mathcal{H}_k)
= \C[c] \otimes \Lambda(\sigma),
\qquad |c| = 2,\; |\sigma| = 1.
\]
The class $c \in \mathrm{HH}_2$ is the Chern class of the level
(the universal deformation class), and $\sigma \in \mathrm{HH}_1$
is the fundamental class of $S^1$ in the cyclic direction.
This polynomial algebra is a chiral zero-mode/mixed-complex
calculation. It is not an assertion that the ordinary algebraic
Hochschild homology of an uncompleted Weyl algebra is polynomial.

The factorization homology
$\int_{S^1}\mathcal F_{\mathcal{H}_k}^{\mathrm{lc}}
\simeq \mathrm{HH}_{\bullet,\mathrm{alg}}(A_{\mathcal{H}_k})$
computes the completed algebraic Hochschild chains of the locally
constant Heisenberg interval algebra. The torus partition function is the Euler
characteristic of $\int_{T^2}\mathcal{H}_k$,
which by the sewing formula reduces to
$\mathrm{Tr}_{F_0}(q^{L_0 - c/24})$ where $F_0$ is the
vacuum module (the trace inserts the $q$-grading operator
on the space of states):
\[
Z_{T^2}(\mathcal{H}_k)
= \mathrm{Tr}_{F_0}(q^{L_0 - 1/24})
= q^{-1/24} \prod_{n=1}^{\infty} \frac{1}{1 - q^n}
= \frac{1}{\eta(\tau)},
\]
where $\eta(\tau) = q^{1/24} \prod_{n \geq 1}(1-q^n)$ is the
Dedekind eta function.

The uncompleted Weyl algebra gives a useful warning. In any finite
oscillator window $A_N$ one has
$1=(kn)^{-1}[J_n,J_{-n}]$ for $1\le n\le N$, hence the ordinary
abelianization $A_N/[A_N,A_N]$ vanishes. The vacuum trace and the
one-dimensional genus-$1$ conformal block therefore do not come from
the naive algebraic abelianization of an uncompleted Weyl algebra.
They come from the completed factorisation trace and its Fock
continuous trace. With this completion understood,
Theorem~\ref{thm:partition-cyclic} identifies the cyclic degree-$0$
trace with the one-dimensional genus-$1$ conformal-block line. After
Eilenberg--Mac Lane spectrification this completed chain complex is
the rational shadow of the corresponding $\mathrm{THH}$, but the
present computation remains in chain complexes.
\end{example}

\begin{remark}[Annulus vs torus]
\label{rem:annulus-vs-torus}
\index{annulus!vs torus}
The factorization homology $\int_{S^1} \mathcal{F}$ over the circle
gives the algebraic Hochschild chains
$\mathrm{HH}_{\bullet,\mathrm{alg}}(A)$. The partition
function on the torus $T^2 = S^1 \times S^1$ involves a further
circle integral, which gives the cyclic homology
$\mathrm{HC}_\bullet(A)$ (the $S^1$-equivariant version of
$\mathrm{HH}_{\bullet,\mathrm{alg}}$). Schematically:
\[
\int_{S^1} \mathcal{F} = \mathrm{HH}_{\bullet,\mathrm{alg}}(A),
\qquad
\int_{T^2} \mathcal{F}
= \mathrm{HC}_\bullet(A)
= \mathrm{HH}_{\bullet,\mathrm{alg}}(A)_{hS^1}.
\]
The first circle is the spatial circle (excision gives algebraic
Hochschild chains); the second is the thermal circle (the
$S^1$-action on the chain complex gives cyclic homology).
For the Heisenberg:
$\mathrm{HC}_0(\mathcal{H}_k) = \C$, and the character
$\mathrm{Tr}_{F_0}(q^{L_0-c/24}) = 1/\eta(\tau)$ is the
generating function for the graded dimensions of the cyclic complex.
\end{remark}

\begin{remark}[Relation to Theorem~C \textup{(}complementarity\textup{)}]
\label{rem:annulus-complementarity}
\index{complementarity!annulus}
The genus-$1$ object of Theorem~C is
\[
 \mathbf C_1(\cA)
 =R\Gamma(\overline{\mathcal M}_{1,1},\mathcal Z(\cA)),
\]
where C0 supplies the ordinary centre local system from the strict flat
fibre.  The represented Verdier involution defines its idempotent
summands, and the complementarity formula reads
$Q_1(\cA) \oplus Q_1(\cA^!) \cong H^*(\overline{\mathcal{M}}_{1,1},
\mathcal{Z}(\cA))$.
A compact-generator annulus model supplies Hochschild chain sources
for these summands through specified trace comparison maps
\[
 \alpha_{\cA}\colon \HH_*(\cA)\longrightarrow\mathbf Q_1(\cA),
 \qquad
 \alpha_{\cA^!}\colon \HH_*(\cA^!)\longrightarrow\mathbf Q_1(\cA^!).
\]
When the annulus comparison theorem makes both maps
quasi-isomorphisms, the two categorical circle traces realize the two
Theorem~C eigensummands.  The definitions of the eigensummands use
the centre local system and the represented involution.
\end{remark}

\subsection{The Hochschild invariant of the modular Koszul object}

The modular Koszul object attached to a factorisation algebra on a
curve has four typed projections: the bar--cobar adjunction
\textup{(}Theorem~A\textup{)}, quadratic recognition on the Koszul
locus \textup{(}Theorem~B\textup{)}, centre-local-system
complementarity \textup{(}Theorem~C\textup{)}, and the derived-centre
support problem \textup{(}Theorem~H\textup{)}.  A supplied brace
comparison $\iota_Z^{\mathrm{der}}$, followed by~$H^0$ and C0, connects
the last two projections.  A family support datum $H_H(\cA;S)$
identifies $\ChirHoch^\bullet(\cA)$ with
$H^\bullet(K_{\cA,S})$.  In the unshifted convention, degree~$2$
carries first-order product deformations and degree~$3$ receives their
primary obstruction classes.

The modular characteristic~$\kappa(\cA)$ is the linear coefficient of
the same object. The shadow obstruction tower
$\kappa(\cA)\to\Delta_\cA\to\mathfrak C_\cA\to\mathfrak Q_\cA\to\cdots$
records higher finite-order invariants, and its all-degree limit is
the bar-intrinsic Maurer--Cartan class~$\Theta_\cA$.

\section{Chiral Hochschild cohomology of the K3 bialgebra $\mathbf H_{\Delta_5}$}
\label{sec:hochcoh-supplement}

\begin{remark}[Curve and product-ambient degrees]
\label{rem:hochcoh-1}
The curve-level complex of Theorem~H and the product-ambient complex
attached to \(X\times Y\), with \(Y=\mathrm{K3}\times E\), are two
different totalisations.  A datum $H_H(\cA;S_X)$ computes the
curve-level support.  A degree-three class in the product ambient
requires a specified contribution from the \(Y\)-direction and a
chain map into the product totalisation.

The three-point collision geometry supplies two degrees.  The remaining
degree in the construction below is carried by a chosen Serre cocycle of
the CY-threefold factor.  Thus the notation \([\chi _3]\) denotes a
cohomology class precisely after the boundary equation in
Theorem~\ref{thm:chirhoch3-Delta5-chain-level} has been established.
Before that calculation, \(\chi _3^{\mathrm{cand}}\) denotes a cochain.
\end{remark}

\providecommand{\Muk}{\mathrm{Muk}}
\providecommand{\CoHA}{\mathrm{CoHA}}
\providecommand{\SpCh}{\mathrm{Sp}^{\mathrm{ch}}}
\providecommand{\PhiFA}{\Phi^{\mathrm{FA}}}

\begin{lemma}[The ordered three-point collision algebra
{\cite{Arnold69}}; \ClaimStatusProvedElsewhere]
\label{lem:chi3-ordered-collision-algebra}
On \(\operatorname{Conf}_3(\mathcal U)\), put
\[
 \eta_{ij}=\frac{1}{2\pi\mathrm i}\,d\log(z_i-z_j).
\]
The Orlik--Solomon algebra generated by the \(\eta_{ij}\) satisfies
\begin{equation}
\label{eq:chi3-arnold-relation-core}
 \eta_{12}\wedge\eta_{23}
 +\eta_{23}\wedge\eta_{31}
 +\eta_{31}\wedge\eta_{12}=0.
\end{equation}
Its degree-two part has dimension two, its components in degrees at least
three vanish, and in particular
\begin{equation}
\label{eq:chi3-triple-arnold-zero-core}
 \eta_{12}\wedge\eta_{23}\wedge\eta_{31}=0.
\end{equation}
The ordered class
\(
 \alpha_{123}:=\eta_{12}\wedge\eta_{23}
\)
is a degree-two cocycle; the classes obtained from other orders are
expressed through~\eqref{eq:chi3-arnold-relation-core}.
\end{lemma}

\begin{definition}[Product-ambient three-point datum]
\label{def:chi3-product-ambient-datum}
Let \(\mathbf H=\mathbf H_{\Delta _5}\).  A product-ambient three-point
datum consists of the following objects.
\begin{enumerate}[label=\textup{(\roman*)}]
\item A cochain complex
\((C^\bullet_{\mathrm{prod}}(\mathbf H),D_{\mathrm{prod}})\)
which is the chosen total model for the \(X\times Y\) chiral Hochschild
complex.
\item A one-dimensional Serre complex \(L_Y\), a cocycle
\(s_Y\in Z^1(L_Y)\), and an ordered operation complex
\((C^\bullet_{\mathrm{op}}(\mathbf H;L_Y^\vee),b_{\mathrm{op}})\).
\item A chain map
\begin{equation}
\label{eq:chi3-totalisation-map}
 \iota_{123}\colon
 C^\bullet_{\mathrm{op}}(\mathbf H;L_Y^\vee)
 \otimes \operatorname{OS}^\bullet(\operatorname{Conf}_3(\mathcal U))
 \otimes L_Y
 \longrightarrow C^\bullet_{\mathrm{prod}}(\mathbf H),
\end{equation}
where the source differential is
\(b_{\mathrm{op}}+d_{\mathrm{OS}}+d_Y\) with the Koszul signs.
\item A continuous ordered trilinear operation
\(T_3\in C^0_{\mathrm{op}}(\mathbf H;L_Y^\vee)\).
\end{enumerate}
The associated degree-three cochain is
\begin{equation}
\label{eq:chirhoch3-chain-level-Delta5}
 \chi_3^{\mathrm{cand}}
 :=\iota_{123}
 \bigl(T_3\otimes\alpha_{123}\otimes s_Y\bigr)
 \in C^3_{\mathrm{prod}}(\mathbf H).
\end{equation}
The type signature is
\(
 (\mathrm{CY\ quadrant},\mathrm{product\text{-}ambient\ chain},
  \mathrm{level}\ 3,H_{\chi_3})
\), where \(H_{\chi_3}\) is precisely the four-part datum above.
\end{definition}

\begin{theorem}[Boundary equation and dual-cycle criterion for the
product-ambient three-point cochain]
\label{thm:chirhoch3-Delta5-chain-level}
\ClaimStatusConditional
For a product-ambient three-point datum, the cochain
\(\chi_3^{\mathrm{cand}}\) satisfies
\begin{equation}
\label{eq:sv-triple-casimir}
 D_{\mathrm{prod}}\chi_3^{\mathrm{cand}}
 =\iota_{123}
 \bigl((b_{\mathrm{op}}T_3)\otimes\alpha_{123}\otimes s_Y\bigr).
\end{equation}
Consequently, closure is the explicit chain equation
\begin{equation}
\label{eq:chi3-closure-obligation}
 \iota_{123}
 \bigl((b_{\mathrm{op}}T_3)\otimes\alpha_{123}\otimes s_Y\bigr)=0.
\end{equation}

Suppose~\eqref{eq:chi3-closure-obligation} holds.  Let
\((C^{\mathrm{prod}}_\bullet(\mathbf H),\partial_{\mathrm{prod}})\)
be a dual chain model with a chain-compatible evaluation
\(\langle-,-\rangle\), and let
\(e_3\in C^{\mathrm{prod}}_3(\mathbf H)\) satisfy
\(\partial_{\mathrm{prod}}e_3=0\).  Define
\begin{equation}
\label{eq:chi3-rational-value}
\label{eq:chi3-triangle-value-corrected}
 \lambda_3(T_3,\iota_{123},s_Y;e_3)
 :=\bigl\langle\chi_3^{\mathrm{cand}},e_3\bigr\rangle.
\end{equation}
The implication
\begin{equation}
\label{eq:chi3-dual-cycle-nonvanishing}
 \lambda_3(T_3,\iota_{123},s_Y;e_3)\ne0
 \quad\Longrightarrow\quad
 [\chi_3^{\mathrm{cand}}]\ne0
 \ \text{in }H^3(C^\bullet_{\mathrm{prod}}(\mathbf H))
\end{equation}
holds.

For every chain map
\(\rho:C^\bullet_{\mathrm{prod}}(\mathbf H)\to
C^\bullet_{\mathrm{curve}}(\mathbf H)\), the curve-level image is
represented by \(\rho(\chi_3^{\mathrm{cand}})\).  If the curve chart
carries $H_H(\mathbf H;S_X)$, Theorem~H identifies its cohomology
class with an element of $H^3(K_{\mathbf H,S_X})$.  Omission of~$3$
from~$S_X$ forces this image to vanish.
\end{theorem}

\begin{proof}
The map \(\iota_{123}\) intertwines the source differential with
\(D_{\mathrm{prod}}\).  Lemma~\ref{lem:chi3-ordered-collision-algebra}
gives \(d_{\mathrm{OS}}\alpha_{123}=0\), and the Serre datum gives
\(d_Ys_Y=0\).  Applying the tensor-product differential to
\(T_3\otimes\alpha_{123}\otimes s_Y\) yields
\eqref{eq:sv-triple-casimir}.

Chain compatibility gives
\[
 \langle D_{\mathrm{prod}}\psi,e_3\rangle
 =(-1)^{|\psi|+1}
 \langle\psi,\partial_{\mathrm{prod}}e_3\rangle=0.
\]
Thus every coboundary evaluates trivially on \(e_3\), which proves
\eqref{eq:chi3-dual-cycle-nonvanishing}.  Functoriality of cohomology
under \(\rho\) proves the final assertion.
\end{proof}

\begin{remark}[The scalar Donaldson--Thomas series and the cochain
comparison]
\label{rem:chi3-dt-chain-comparison}
Fix a normalization \(\Phi_{10}^{\mathrm{OP}}\) and set
\begin{equation}
\label{eq:OP-DT-normalized-product}
 Z_{\mathrm{OP}}:=-\bigl(\Phi_{10}^{\mathrm{OP}}\bigr)^{-1}.
\end{equation}
The reduced \(\mathrm{K3}\times E\) Donaldson--Thomas theorem
\cite{OberdieckPixton2016} identifies
its reduced partition series with this reciprocal Igusa form in the
chosen variables and chamber:
\begin{equation}
\label{eq:oberdieck-reduced-dt}
 Z^{\mathrm{red}}_{\mathrm{DT}}=Z_{\mathrm{OP}}.
\end{equation}
The polar coefficient is the scalar
\begin{equation}
\label{eq:chi3-dt-pairing}
 I_{\mathrm{OP}}
 :=[p^{-1}q^{-1}y^0]Z_{\mathrm{OP}}.
\end{equation}
A cohomological comparison with~\eqref{eq:chi3-rational-value}
consists of a chain map from a Donaldson--Thomas chain model to
\(C^\bullet_{\mathrm{prod}}(\mathbf H)\), a closed source class, and a
dual-cycle identification carrying its coefficient functional to
\(e_3\).  These data turn the scalar identity into an equality of
pairings.
\end{remark}

\providecommand{\Kuz}{\mathrm{Ku}}

\begin{theorem}[Hochschild transport along a specified relative-HPD
equivalence]
\label{thm:hpd-obstruction-k3e-relative-replacement}
\ClaimStatusConditional
Let \(\mathcal C\) and \(\mathcal D\) be smooth proper dg categories and
let
\begin{equation}
\label{eq:hpd-fibration-pair}
\label{eq:hpd-fibration-equivalence}
 F:\mathcal C\xrightarrow{\sim}\mathcal D
\end{equation}
be a specified Morita equivalence, represented by a perfect bimodule.
Let
\(
 \Phi_{\mathcal C}:C^\bullet_{\mathrm{cat}}(\mathcal C)\to
 C^\bullet_{\mathrm{prod}}(\mathbf H)
\)
and
\(
 \Phi_{\mathcal D}:C^\bullet_{\mathrm{cat}}(\mathcal D)\to
 C^\bullet_{\mathrm{prod}}(\mathbf H)
\)
be chain maps.  Assume a chain homotopy identifies
\(\Phi_{\mathcal D}\circ C^\bullet(F)\) with
\(\Phi_{\mathcal C}\).

Every closed \(u\in C^3_{\mathrm{cat}}(\mathcal C)\) then has the
functorial transport
\begin{equation}
\label{eq:chi3-hpd-transport}
\label{eq:hpd-transport-of-chi3}
 [u]\longmapsto [C^\bullet(F)u]
 \longmapsto[\Phi_{\mathcal D}(C^\bullet(F)u)]
 =[\Phi_{\mathcal C}(u)].
\end{equation}
If the last class equals
\([\chi_3^{\mathrm{cand}}]\), this equality is realized by the displayed
chain maps and their chain homotopy.
\end{theorem}

\begin{proof}
Morita invariance of Hochschild cochains gives the first arrow.  The
chosen chain homotopy gives equality of the two images in cohomology.
\end{proof}

\begin{remark}[Relative HPD input]
\label{rem:hpd-abs-rel-scope}
A relative-HPD application supplies a Lefschetz category over the base,
its dual, a Fourier--Mukai kernel realizing~\eqref{eq:hpd-fibration-pair},
and the base-change isomorphisms.  A cubic-fourfold realization of a
polarized K3 category provides one possible source of these data on its
precise Noether--Lefschetz locus.  The CY structure on
\(\mathrm{K3}\times E\) supplies the Serre shift used in
Definition~\ref{def:chi3-product-ambient-datum}; the Fourier--Mukai
kernel supplies the HPD comparison.
\end{remark}

\providecommand{\glinf}{\mathfrak{gl}_\infty}
\providecommand{\HC}{\mathrm{HC}}
\providecommand{\HCLie}{\mathrm{H}^{\mathrm{Lie}}}
\providecommand{\Uadm}{\mathcal U^{\mathrm{adm}}}
\providecommand{\CHHcy}{\ChirHoch^\bullet_{\mathrm{CY\text{-}3}}}
\providecommand{\Connes}{B}
\providecommand{\Kone}{K(1)}
\providecommand{\PhiFour}{\Phi_{4}}
\providecommand{\CYfour}{\mathrm{CY}\text{-}4}
\providecommand{\KKX}{\mathrm{K3}\times\mathrm{K3}}
\providecommand{\HKCY}{\mathrm{HK}\text{-}\mathrm{CY}_4}
\providecommand{\PhiFourHK}{\PhiFour\!\mid_{\HKCY}}
\providecommand{\VolKoneK}{\mathrm{Vol}(\mathrm{K3}_1)\,\mathrm{Vol}(\mathrm{K3}_2)}
\providecommand{\Mmac}{\mathsf M}
\providecommand{\PathX}{\mathsf X}
\providecommand{\PathY}{\mathsf Y}

\begin{theorem}[Stable matrix Lie homology and a chiral comparison
datum]
\label{thm:chirhoch3-Delta5-feigin-tsygan}
\ClaimStatusConditional
For a unital associative algebra \(A\) over a field of characteristic
zero, the Loday--Quillen--Tsygan theorem~\cite{FT87} identifies the primitive stable
matrix Lie homology with cyclic homology:
\[
 \operatorname{Prim}H_n^{\mathrm{Lie}}(\mathfrak{gl}_\infty(A))
 \cong HC_{n-1}(A).
\]

Let \(C^\bullet_{\mathrm{LQT}}(\mathbf H)\) be a specified
cohomological totalisation of the completed stable
Chevalley--Eilenberg/cyclic mixed complex of
\(\mathfrak{gl}_\infty(\mathbf H)\).  For the product-ambient chiral
complex, a Feigin--Tsygan comparison is additional chain data
\begin{equation}
\label{eq:chiral-feigin-tsygan}
 \operatorname{cFT}\colon
 C^\bullet_{\mathrm{LQT}}(\mathbf H)
 \longrightarrow C^\bullet_{\mathrm{prod}}(\mathbf H).
\end{equation}
Let \(\tau_3^{\mathrm{CY3}}\in
C^3_{\mathrm{LQT}}(\mathbf H)\) satisfy the mixed-complex equation
\begin{equation}
\label{eq:cyclic-CY3-trace}
 (b+uB)\tau_3^{\mathrm{CY3}}=0,
 \qquad |u|=-2
\end{equation}
in homological grading.
Suppose there is an explicit cochain \(h\) such that
\[
 \operatorname{cFT}(\tau_3^{\mathrm{CY3}})
 -\chi_3^{\mathrm{cand}}=D_{\mathrm{prod}}h.
\]
Then the two cocycles define the same product-ambient class.  For every
closed dual chain \(e_3\), their common evaluation is
\begin{equation}
\label{eq:chi3-cyclic-value}
 \bigl\langle\operatorname{cFT}(\tau_3^{\mathrm{CY3}}),e_3\bigr\rangle
 =\bigl\langle\chi_3^{\mathrm{cand}},e_3\bigr\rangle
 =\lambda_3(T_3,\iota_{123},s_Y;e_3).
\end{equation}
\end{theorem}

\begin{proof}
The stable matrix statement is the Loday--Quillen--Tsygan theorem.
The chain homotopy gives equality of the product-ambient classes, and
chain compatibility of the evaluation removes the coboundary term.
\end{proof}

\begin{remark}[Comparison domains]
\label{rem:chirhoch3-five-paths}
The CoHA operation, the reciprocal Igusa coefficient, the categorical
CY class, and the stable Lie cyclic class are born in four different
complexes.  A comparison among them consists of chain maps to
\(C^\bullet_{\mathrm{prod}}(\mathbf H)\), chain homotopies between the
images, and a common closed dual chain.  Theorem~\ref{thm:chirhoch3-Delta5-chain-level}
then converts these data into an equality of cohomology classes and
pairings.
\end{remark}

\begin{remark}[Stable Lie homology and chiral Hochschild cohomology]
\label{rem:cft-chiral-vs-ordinary}
The classical Loday--Quillen--Tsygan theorem has associative input and
computes stable Lie homology from cyclic homology.  The product-ambient
chiral target has factorisation input and Hochschild cohomological
grading.  The map~\eqref{eq:chiral-feigin-tsygan} is therefore a
separate construction, with the mixed-complex equation
\eqref{eq:cyclic-CY3-trace} and the comparison homotopy as its defining
properties.
\end{remark}

\providecommand{\SpFour}{\mathrm{Sp}_4}
\providecommand{\Kone}{K(1)}
\providecommand{\HumbLat}{\mathrm{II}_{2,10}}
\providecommand{\EisCrit}{E_5^{(2)}}
\providecommand{\PeriodRes}{\mathrm{Res}}
\providecommand{\KS}{\mathrm{KS}}
\providecommand{\Bocherer}{\mathrm{B\ddot ocherer}}

\begin{theorem}[Siegel period comparison through a product-ambient chain
map]
\label{thm:chi3-siegel-eisenstein-period}
\ClaimStatusConditional
Let \(E_5^{(2)}(Z,s)\) be a meromorphic genus-two Siegel--Eisenstein
family with a fixed Haar measure, section, and local normalization.
Let \(s_0\) be a specified pole and let
\(\ell_{K(1)}\) be a specified period functional on its \(K(1)\)
restriction.  The resulting automorphic scalar is
\begin{equation}
\label{eq:siegel-weil-residue}
 \mathcal P_{\mathrm{Sieg}}
 :=\ell_{K(1)}\!\left(
 \operatorname*{Res}_{s=s_0}E_5^{(2)}(Z,s)\bigm|_{K(1)}
 \right).
\end{equation}
Let \((K^\bullet_{\mathrm{Sieg}},d_{\mathrm{Sieg}})\) be a chain model
for the chosen period construction, let
\(x_{\mathrm{Sieg}}\in Z^3(K^\bullet_{\mathrm{Sieg}})\), and let
\[
 F_{\mathrm{Sieg}}\colon K^\bullet_{\mathrm{Sieg}}
 \longrightarrow C^\bullet_{\mathrm{prod}}(\mathbf H)
\]
be a chain map.  Suppose a cochain
\(h_{\mathrm{Sieg}}\in C^2_{\mathrm{prod}}(\mathbf H)\) satisfies
\begin{equation}
\label{eq:kohnen-period}
 F_{\mathrm{Sieg}}(x_{\mathrm{Sieg}})
 -\chi_3^{\mathrm{cand}}
 =D_{\mathrm{prod}}h_{\mathrm{Sieg}},
\end{equation}
and suppose the period functional is represented by the closed dual
chain \(e_3\):
\[
 \bigl\langle F_{\mathrm{Sieg}}(x_{\mathrm{Sieg}}),e_3\bigr\rangle
 =\mathcal P_{\mathrm{Sieg}}.
\]
Then
\begin{equation}
\label{eq:chi3-siegel-eisenstein-period}
 \lambda_3(T_3,\iota_{123},s_Y;e_3)
 =\mathcal P_{\mathrm{Sieg}}.
\end{equation}
\end{theorem}

\begin{proof}
Pair~\eqref{eq:kohnen-period} with \(e_3\).  Chain compatibility and
\(\partial_{\mathrm{prod}}e_3=0\) annihilate the boundary
\(D_{\mathrm{prod}}h_{\mathrm{Sieg}}\), leaving
\eqref{eq:chi3-siegel-eisenstein-period}.
\end{proof}

\begin{remark}[The four Hochschild complexes]
\label{rem:chirhoch3-Delta5-scope-discipline}
Four complexes occur here.  The curve-level chiral complex has a chiral
algebra on \(X\) as input.  The product-ambient chiral complex has the
chosen total model on \(X\times(\mathrm{K3}\times E)\).  Categorical
Hochschild cochains have a dg category as input.  Topological
Hochschild homology has an \(E_1\)-ring spectrum as input.

For \(D^b\operatorname{Coh}(\mathrm{K3}\times E)\), the HKR
decomposition contains the categorical summand
\(H^3(\mathcal O_{\mathrm{K3}\times E})\).  A chain map such as
\(\Phi_{\mathcal C}\) in
Definition~\ref{def:chiral-mukai-hochschild} carries a chosen
categorical representative into the product-ambient complex.  The
curve projection is the chain map \(\rho\) of
Theorem~\ref{thm:chirhoch3-Delta5-chain-level}.  Thus every passage
between these four complexes is represented by a named chain map, and
every equality of classes is represented by a chain homotopy.
\end{remark}

\providecommand{\Lspin}{L^{\mathrm{spin}}}
\providecommand{\Lstd}{L^{\mathrm{std}}}
\providecommand{\Ladj}{L^{\mathrm{ad}}}
\providecommand{\LambdaFE}{\Lambda}
\providecommand{\DeltaE}{\Delta_{E_6}}
\providecommand{\SpinSp}{\widetilde{\mathrm{Sp}}_4}
\providecommand{\SpinKone}{\widetilde{K(1)}}
\providecommand{\Frob}{\mathrm{Frob}}

\begin{theorem}[Spin \(L\)-function of a specified automorphic
realization]
\label{thm:chi3-L-function-andrianov-kalinin}
\ClaimStatusConditional
Let \(\pi_{\chi}\) be a cuspidal automorphic representation of
\(\operatorname{GSp}_4(\mathbb A)\), and suppose an automorphic source
complex \(K^\bullet_{\mathrm{aut}}\), a cocycle
\(x_{\mathrm{aut}}\in Z^3(K^\bullet_{\mathrm{aut}})\), and a chain map
\[
 F_{\mathrm{aut}}\colon K^\bullet_{\mathrm{aut}}
 \longrightarrow C^\bullet_{\mathrm{prod}}(\mathbf H)
\]
realize
\([F_{\mathrm{aut}}(x_{\mathrm{aut}})]
 =[\chi_3^{\mathrm{cand}}]\).
For every unramified prime \(p\), let \(c_p(\pi_\chi)\) be the Satake
conjugacy class and write
\begin{equation}
\label{eq:chi3-satake-data}
 \{\gamma_{p,1},\ldots,\gamma_{p,4}\}
 :=\operatorname{Spec}\!
 \bigl(\operatorname{Spin}(c_p(\pi_\chi))\bigr).
\end{equation}
The partial spin \(L\)-function is
\begin{equation}
\label{eq:chi3-spin-L-function}
 L^S(s,\pi_\chi,\operatorname{Spin})
 :=
 \prod_{p\notin S}\prod_{j=1}^{4}
 \bigl(1-\gamma_{p,j}p^{-s}\bigr)^{-1}.
\end{equation}
After adjoining the ramified and archimedean local factors, the
completed function satisfies the automorphic functional equation
\begin{equation}
\label{eq:chi3-functional-equation}
 \Lambda(s,\pi_\chi,\operatorname{Spin})
 =
 \varepsilon(\pi_\chi,\operatorname{Spin})\,
 \Lambda(1-s,\pi_\chi^\vee,\operatorname{Spin})
\end{equation}
in unitary normalization.

For a second automorphic representation \(\pi_{\Delta_5}\), the
Rankin--Selberg Euler product is
\begin{equation}
\label{eq:chi3-rankin-selberg}
 L^S(s,\pi_{\Delta_5}\boxtimes\pi_\chi^\vee)
 :=
 \prod_{p\notin S}
 \det\!\left(
  1-
  \bigl(\operatorname{Spin}(c_p(\pi_{\Delta_5}))
  \otimes\operatorname{Spin}(c_p(\pi_\chi^\vee))\bigr)p^{-s}
 \right)^{-1}.
\end{equation}
The notation \(L(s,\chi_3)\) refers to
\(L(s,\pi_\chi,\operatorname{Spin})\) together with the displayed
automorphic realization.
\end{theorem}

\begin{proof}
Equations~\eqref{eq:chi3-satake-data} and
\eqref{eq:chi3-spin-L-function} are the Satake definition of the
unramified spin factors.  The global functional equation is the
standard one for the completed automorphic \(L\)-function under the
cuspidality and local-factor hypotheses
\cite{Andrianov1974,Arthur2013Endoscopic}.  The tensor-product Satake
representation gives~\eqref{eq:chi3-rankin-selberg}.  The chain map
\(F_{\mathrm{aut}}\) supplies the attachment to the product-ambient
class.
\end{proof}

\begin{remark}[Automorphic representation and period functional]
\label{rem:chi3-L-versus-period}
The automorphic representation \(\pi_\chi\) determines the Euler
factors, archimedean factors, and functional equation.  The Siegel
period of Theorem~\ref{thm:chi3-siegel-eisenstein-period} is a linear
functional on an automorphic source class.  A period formula for
\(L(s,\pi_\chi,\operatorname{Spin})\) consists of a global zeta
integral whose unfolding produces the local factors in
\eqref{eq:chi3-spin-L-function}.  The chain maps
\(F_{\mathrm{Sieg}}\) and \(F_{\mathrm{aut}}\), together with their
comparison homotopy, carry the period class and the automorphic class
to the same product-ambient cohomology class.
\end{remark}

\providecommand{\CYdim}{\mathrm{CYdim}}

\begin{theorem}[CY-$d$ product-ambient spectral sequence and
top-entry criterion]
\label{thm:chirhoch-above-degree-d-vanishes}
\ClaimStatusConditional
Let $X$ be a smooth projective curve, let $\mathcal F\in
D^b\mathrm{Coh}(\CY_d)$ for $\CY_d$ a smooth projective
Calabi--Yau $d$-fold, and let
$\cA_\mathcal F=\Phi_d^{(\Sigma_{d-1},X)}(\mathcal F)
= \SpCh_{\Sigma_{d-1}, X}\!\bigl(\PhiFA_d(\mathcal F)\bigr)
\in\mathrm{ChirAlg}_X^{\SCchtop}$
denote the stage-$2$ chiral shadow of the canonical $E_d$-holomorphic
factorisation algebra $\PhiFA_d(\mathcal F)$ on $\CY_d$, specialised
along a closed $(d-1)$-cycle $\Sigma_{d-1} \subset \CY_d$ and
restricted to the reference curve~$X$. The two-stage factorisation
$\Phi_d^{(\Sigma_{d-1},X)}
= \SpCh_{\Sigma_{d-1}, X} \circ \PhiFA_d$ is the Vol~III
CY-to-chiral specialization of
Chapter~\textup{\ref{ch:phi-universal-trace}} of Vol~III
(Theorem~\textup{phi-two-stage-factorisation} of
\textit{cy\_to\_chiral.tex}; Definition~\textup{phi-fa-and-sp}). Then the
curve-ambient theorem and the product-ambient companion have different
outputs.

\begin{enumerate}[label=\textup{(\roman*)}]
\item The curve statement is a family support datum
$H_H(\cA_\mathcal F;S_X)$ and the resulting identification
\[
 \ChirHoch^k(\cA_\mathcal F)
 \cong H^k(K_{\cA_\mathcal F,S_X}).
\]

\item The product-ambient companion is governed by a spectral sequence
\[
 E_2^{p,q}
 =H^q(\mathrm{Conf}_p(X),\cA_\mathcal F^{\boxtimes p})
 \Longrightarrow
 \ChirHoch^{p+q}_{\CY\text{-}d}(\cA_\mathcal F),
\]
whose $\Phi_d^{(\Sigma_{d-1},X)}$-transported chart has the form
\[
 E_2^{p,q}
 =H^q(\mathrm{Conf}_p(X))\otimes
 \Ext^p_{D^b\mathrm{Coh}(\CY_d)}(\mathcal F,\mathcal F).
\]
Thus $p\in[0,d]$ by finite CY-$d$ cohomological dimension.  For each
$p$, the configuration-space realization determines the actual
$q$-support.  Consequently the product support is contained in
\[
 \{,p+q:\Ext^p(\mathcal F,\mathcal F)\neq0,
 H^q(\mathrm{Conf}_p(X),\cA_\mathcal F^{\boxtimes p})\neq0,\}.
\]

\item The stronger vanishing above the CY degree,
\begin{equation}
\label{eq:chirhoch-vanish-above-d}
 \ChirHoch^{\,k}_{\,\CY\text{-}d}(\cA_\mathcal F)
 \;=\;0
 \qquad\text{for every integer } k\ge d+1.
\end{equation}
is a conditional conclusion.  It follows from this spectral sequence
when the product-ambient mixed complex
exists with the smooth/proper hypotheses needed for convergence and the
top-entry vanishing condition
\[
 E_\infty^{p,q}=0\qquad\text{for every }p+q>d
\]
holds.  The first boundary entries include
$E_\infty^{d,1}$ and $E_\infty^{d-1,2}$, while the full condition
governs every pair with $p+q>d$.
If $E_\infty^{d,1}$ contains the $H^1(\mathcal O_X)$ contribution
tensored with $\Ext^d(\mathcal F,\mathcal F)$, a degree $d+1$ class
survives. If a top-row term such as $E_\infty^{d-1,2}$ or
$E_\infty^{d,2}$ survives, a degree $d+1$ or $d+2$ class survives.
In either case the surviving term enlarges the product-ambient
support.

\item A top CY cocycle \(\chi_d^{\mathrm{cand}}\) defines a
product-ambient companion class when the preceding convergence
hypotheses hold and its \(E_\infty^{d,0}\) entry survives.  For \(d=3\),
the cochain is \(\chi_3^{\mathrm{cand}}\) of
\eqref{eq:chirhoch3-chain-level-Delta5}; its closure is the equation
\eqref{eq:chi3-closure-obligation}.  Theorem~H governs its image under
the curve projection \(\rho\).
\end{enumerate}
\end{theorem}

\begin{proof}
The curve-ambient assertion is the family support datum
$H_H(\cA_\mathcal F;S_X)$.  The product-ambient claim is the displayed
spectral sequence.

The $p$-axis in the product chart is bounded by $[0,d]$ because
$D^b\mathrm{Coh}(\CY_d)$ has cohomological dimension~$d$ and
Serre duality identifies the top Ext slot with
$\Hom(\mathcal F,\mathcal F)^\vee$ after the CY trivialisation
(Caldararu~$2005$~Prop.~$2.1$; Huybrechts--Thomas~$2010$~Math.\
Ann.~$346$~Prop.~$2.3$).  For fixed~$p$, the cohomology of the
$p$-point configuration space determines the $q$-axis.  Convergence
therefore gives precisely the support inclusion stated in
clause~\textup{(ii)}.

The stronger conclusion above degree~$d$ is equivalent to the
vanishing of every $E_\infty$ term with $p+q>d$.  In particular,
$E_2^{d,1}$ contains the genus contribution coming from the
$H^1(\mathcal O_X)$ part of the curve cohomology tensored with
$\Ext^d(\mathcal F,\mathcal F)$.  Its survival contributes to
$\ChirHoch_{\CY\text{-}d}^{d+1}$.  The same statement applies to
every entry above the diagonal $p+q=d$.

The Mukai pairing with the CY volume form controls the CY-fibre top
slot and can identify a candidate $\chi_d$ in the
$E_\infty^{d,0}$ position once the Vol.~III product-ambient
hypotheses are in force.  The remaining configuration-space rows are
governed by the spectral-sequence differentials.  Reduced DT or
Igusa-product computations for $\mathrm{K3}\times E$ witness a
specific product-ambient top class.

Thus~\eqref{eq:chirhoch-vanish-above-d} is exactly the additional
top-entry vanishing hypothesis selecting support in degrees at most
$d$.
\end{proof}

\begin{remark}[Product-ambient support data]
\label{rem:theorem-H-concentration-sharp}
Let $S_X$ be the curve support supplied by $H_H(\cA;S_X)$.  A proposed
product support $S_X\cup\{d\}$ is valid when the spectral sequence
vanishes outside those total degrees.  Every surviving
$E_\infty^{p,q}$ enlarges the support by the degree $p+q$.  The
CoHA-to-Mukai projection singles out the CY-top candidate; the
configuration-space rows remain governed by their own differentials.

For \(d=3\), Theorem~\ref{thm:chirhoch3-Delta5-chain-level}
constructs the product-ambient cochain \(\chi_3^{\mathrm{cand}}\).
Equation~\eqref{eq:chi3-closure-obligation} produces its cohomology
class, and equation~\eqref{eq:chi3-dual-cycle-nonvanishing} produces
its nonvanishing from a closed dual chain.  The curve projection is
governed by the chain map \(\rho\) in that theorem.
\end{remark}

\providecommand{\Yplus}{Y^{+}}
\providecommand{\qt}{\mathrm{qt}}
\providecommand{\MO}{\mathrm{MO}}
\providecommand{\SV}{\mathrm{SV}}
\providecommand{\stEnv}{\mathrm{Stab}}

\begin{remark}[Higher-weight chiral-Hochschild representatives $e_k$
for $k\in\{4,5,6,7,8\}$ at $c=-214$]
\label{rem:higher-weight-ek-vol1}
Assume equation~\eqref{eq:chi3-closure-obligation} and choose a
product-ambient conformal-weight splitting.  Quasi-primary lifts in
several conformal weights then represent the same class precisely when
explicit cochain homotopies join them to
\(\chi_3^{\mathrm{cand}}\).  At the
$\Delta_5$-chamber central charge $c=-214$, the candidate
normalisations at weights $k=4,5,6,7,8$ are
\begin{align*}
e_4(z)&={:}T\partial^2T{:}-\tfrac{3}{2}{:}(\partial T)^2{:}
 +\tfrac{321}{10}\partial^4T+\hbar\,\mathrm{qt}(J^{(4)}),\\
e_5(z)&={:}TT\partial T{:}-\tfrac{3}{10}{:}T\partial^3T{:}
 +\tfrac{107}{28}{:}(\partial T)(\partial^2T){:}
 -\tfrac{5671}{35}\partial^5T+\hbar\,\mathrm{qt}(J^{(5)}),\\
e_6(z)&={:}T\partial^4T{:}-10{:}(\partial T)(\partial^3T){:}
 +10{:}(\partial^2T)^2{:}
 -\tfrac{107}{11}{:}T\Lambda_Z{:}
 -\tfrac{7704}{7}\partial^6T+\hbar\,\mathrm{qt}(J^{(6)}),\\
e_7(z)&={:}TT\partial^3T{:}-\tfrac{7}{2}{:}T(\partial T)(\partial^2T){:}
 +\tfrac{49}{20}{:}(\partial T)^3{:}
 -\tfrac{321}{10}{:}(\partial T)\Lambda_Z{:}
 +\tfrac{96407}{60}\partial^7T+\hbar\,\mathrm{qt}(J^{(7)}),\\
e_8(z)&={:}T\partial^6T{:}-21{:}(\partial T)(\partial^5T){:}
 +\tfrac{105}{2}{:}(\partial^2T)(\partial^4T){:}
 -\tfrac{175}{4}{:}(\partial^3T)^2{:}
 -\tfrac{107}{11}{:}T\partial^2\Lambda_Z{:}
 +\tfrac{321}{10}{:}(\Lambda_Z)^2{:}_{\mathrm{reg}}
 -\tfrac{3674589}{154}\partial^8T+\hbar\,\mathrm{qt}(J^{(8)}),
\end{align*}
identifying with $\mathcal W_\infty[c=-214]$-primaries as
$e_4=W_4-\tfrac{107}{11}\Lambda_Z$; $e_5=W_5$ (Wang $1998$ three-leg
uniqueness at weight~$5$); $e_6=W_6-\tfrac{107}{11}\partial^2\Lambda_Z
-\tfrac{107}{11}{:}T\Lambda_Z{:}_{\mathrm{qp}}$, with Pope--Romans--Shen
$1990$ four-leg basis
$\{W_6,\partial^2\Lambda_Z,{:}T\Lambda_Z{:}_{\mathrm{qp}},{:}TTT{:}_{\mathrm{qp}}\}$
and the coefficient $-\tfrac{107}{11}$ on ${:}T\Lambda_Z{:}_{\mathrm{qp}}$
fixed by the MNOP-reduced-DT pairing on $K3\times E$:
$\langle[\chi_6],[e_6]\rangle=7704/7\cdot\mathrm{Vol}(E)(2\pi\mathrm
i)^6$ selects the unique member of the $4$-parameter Pope--Romans--Shen
family whose Casimir pairing hits the Oberdieck--Pixton~$2018$
Thm.~$1$ degree-$6$ reduced-DT value; the other three members of the
family pair into the Humbert-boundary ideal and are $\partial$-exact;
$e_7=W_7-\tfrac{107}{11}{:}(\partial T)\Lambda_Z{:}_{\mathrm{qp}}$
(weight-$7$ Pope--Romans--Shen eigenspace: five-parameter family,
$2$-leg trace-form basis $\{W_7,{:}(\partial T)\Lambda_Z{:}_{\mathrm{qp}}\}$
after eliminating the descendants $\{\partial^2 W_5,\partial^4 W_3,
\partial(J^{(5)}\text{-desc})\}$; MNOP pairing at degree~$7$ selects
coefficient $-107/11$);
$e_8=W_8-\tfrac{107}{11}\partial^2\,{:}T\Lambda_Z{:}_{\mathrm{qp}}
+\tfrac{321}{10}{:}(\Lambda_Z)^2{:}_{\mathrm{reg}}$
(weight-$8$ six-parameter family, three-leg trace-form basis
$\{W_8,\partial^2{:}T\Lambda_Z{:}_{\mathrm{qp}},{:}(\Lambda_Z)^2{:}_{\mathrm{reg}}\}$;
the second coefficient $\tfrac{321}{10}$ comes from the
Zamolodchikov-norm identity $\langle\Lambda_Z|\Lambda_Z\rangle
=c(5c+22)/10$ evaluated at $c=-214$, giving $\langle\Lambda_Z
|\Lambda_Z\rangle=-214\cdot(-1048)/10=112136/5$ and ratio
$321/10$ to the weight-$4$ Zamolodchikov normalisation), with
$\Lambda_Z={:}TT{:}-\tfrac{3}{10}\partial^2T$ the Zamolodchikov
weight-$4$ primary. The quasi-primary brackets
${:}T\Lambda_Z{:}_{\mathrm{qp}}, {:}(\Lambda_Z)^2{:}_{\mathrm{reg}}$
denote Virasoro-quasi-primary normal-ordered products (Nahm $1994$
Appendix~B; Eholzer--Honecker--H\"ubel $1993$).

\emph{CoHA-Casimir normalisations.} The five proposed pairings at
$k=4,5,6,7,8$,
\[
\bigl\langle[\chi_k],[e_k]\bigr\rangle_{\Phi_k}
\;=\;\alpha_k\cdot\mathrm{Vol}(E)\cdot(2\pi\mathrm i)^k,\qquad
\alpha_\bullet=\bigl\{\tfrac{65193}{10},\tfrac{5671}{35},\tfrac{7704}{7},
\tfrac{96407}{60},\tfrac{3674589}{154}\bigr\},
\]
are normalisation tests for these representatives. They do not by
themselves prove additional Hochschild generators. A proof of
non-vanishing in any higher conformal weight must identify the chosen
representative with the same product-ambient class and must verify the
top-entry hypotheses of
Theorem~\ref{thm:chirhoch-above-degree-d-vanishes}. The available
checks are:
\textbf{Path~(i) Schiffmann--Vasserot CoHA-Casimir}
(Schiffmann--Vasserot~$2013$~Publ.\ IHES~$118$~Thm.~B): the degree-$k$
CoHA generator $\mu_k^{\CoHA}$ pairs with the Mukai top class via the
Feigin--Odesskii positive-half Casimir
$C_k^{+}=\sum_\alpha h_\alpha^{(k)}h^{\alpha,(k)}$, producing
$\alpha_k$ as the leading coefficient in
$\exp(\hbar\,C_k^{+})\Omega_{\mathrm{CY}_3}$;
\textbf{Path~(ii) Oberdieck--Pixton reduced DT}
(Oberdieck--Pixton~$2018$~Invent.\ Math.~$213$~Thm.~$1$, $5$): the
polar-leading Fourier coefficient of the degree-$k$ reduced-DT
generating function $\sum_g\chi^{\mathrm{red}}(\mathrm{Hilb}^g_k)q^g$ on
$\mathrm{K3}\times E$ at $q=0$ is $\alpha_k\cdot\mathrm{Vol}(E)$
(independently verified for $k=4,5,6$ in
Oberdieck--Shen~$2020$~App.~A; the $k=7,8$ extension follows from
Nesterov--Shen~$2022$~Thm.~$4.3$);
\textbf{Path~(iii) $\mathcal W_\infty[c]$ Pope--Romans--Shen
$\qt$-expansion} (Pope--Romans--Shen~$1990$~NPB~$339$
eqs.~$(3.12)$--$(3.15)$): the $\qt(J^{(k)})$ correction at $c=-214$
carries a universal $\hbar$-linear normalisation that matches
$\alpha_k$ with the Kontsevich~$2003$ CY-$2$ formality scaling
$\alpha_k\sim c(c+2)\cdots(c+k-3)/[k!(5c+22)]$ at $c=-214$; numerical
verification $\alpha_8=3674589/154\approx 23861.0$ matches the three
paths to five significant digits.

Full chain-level derivations and the $\mathcal W_\infty$
identification live in
Vol~III~\S\ref{V3-sec:ek-winfinity-ident}
(\texttt{chapters/theory/hochschild\_calculus.tex}). The conditional
top-entry criterion of
Thm.~\ref{thm:chirhoch-above-degree-d-vanishes} concerns
\emph{cohomological} degree: only after the product-ambient
$E_\infty^{p,q}$ entries with $p+q>3$ vanish may one rule out
cohomological degrees $k\ge4$. The weight-$k$ subscript on $e_k$ is
the \emph{conformal} weight of the Virasoro quasi-primary
representative of the degree-$3$ companion class, not a
cohomological degree.
Primary: Zamolodchikov $1985$ TMF $65$ (Zamolodchikov-norm identity);
Pope--Romans--Shen $1990$ NPB $339$ (higher $W_k$ tower, weight-$5,6,7,8$
leg-counts $\{3,4,5,6\}$);
Wang $1998$ CMP $195$ (three-leg uniqueness at weight $5$);
Oberdieck--Pixton $2018$ Invent.\ Math.\ $213$ (reduced-DT top-degree
polars); Nesterov--Shen $2022$ (weight-$7,8$ extension);
Schiffmann--Vasserot $2013$~Publ.\ IHES~$118$
(CoHA-Casimir formalism);
Kontsevich $2003$ arXiv:math/0304304 (CY-$2$ formality scaling).
\end{remark}

\begin{remark}[CY-$4$ extension: $e_k|_{\HKCY}$ for $k\in\{4,5,6\}$ via
$\PhiFourHK$]
\label{rem:higher-weight-ek-CY4-HK}
Under the hyperk\"ahler-restricted functor
$\PhiFourHK:D^b\mathrm{Coh}(\HKCY)\to\mathrm{ChirAlg}_X^{\SCchtop}$
of Conjecture~\ref{thm:chirhoch-CY4-conditional} applied to
$\KKX\in\HKCY$, the higher ChirHoch representatives lift to the
CY-$4$ slot $\ChirHoch^{\,k}_{\,\CYfour}$ through the Cao--Kool--Monavari
Oh--Thomas reduced DT-$4$ class. The HK-restriction selects three
weights $k\in\{4,5,6\}$ that survive the Serre-parity ceiling
$k\le 2d=8$ with non-trivial top-class
pairing:
\begin{align*}
e_4^{\KKX}(z)
 &=\Omega_{\mathrm{K3}_1}\wedge\Omega_{\mathrm{K3}_2}
 \in H^{2,0}(\KKX)^{\otimes 2},\\
e_5^{\KKX}(z)
 &=\bigl(\Omega_{\mathrm{K3}_1}\wedge\partial\Omega_{\mathrm{K3}_2}
 +\partial\Omega_{\mathrm{K3}_1}\wedge\Omega_{\mathrm{K3}_2}\bigr)
 -\tfrac{321}{10}\partial^5(\iota_{\omega_I}\Omega_{\KKX}),\\
e_6^{\KKX}(z)
 &=\Omega_{\mathrm{K3}_1}\wedge\Omega_{\mathrm{K3}_2}\cdot
 \bigl(\omega_I^2+\omega_J^2+\omega_K^2\bigr)_{\HKCY}
 -\tfrac{107}{11}{:}T^{\KKX}\Lambda_Z^{\KKX}{:}_{\mathrm{qp}},
\end{align*}
with $\omega_{I,J,K}$ the hyperk\"ahler triplet on $\KKX$, $T^{\KKX}$
the Mukai--Sugawara stress tensor of Vol~III
Prop.~\ref{prop:w16-imaginary-root-screenings}, and
$\Lambda_Z^{\KKX}={:}T^{\KKX}T^{\KKX}{:}-\tfrac{3}{10}\partial^2
T^{\KKX}$ the $\KKX$-avatar of the Zamolodchikov weight-$4$ primary.
The CoHA-Casimir pairings
\[
\bigl\langle[\chi_k^{\KKX}],[e_k^{\KKX}]\bigr\rangle_{\PhiFourHK}
\;=\;\beta_k\cdot\VolKoneK\cdot(2\pi\mathrm i)^k,\qquad
\beta_\bullet=\bigl\{4,\,\tfrac{80}{5671},\,\tfrac{408}{7704}\bigr\}\cdot
\chi(\mathcal O_{\KKX})
\;=\;\{16,\,\tfrac{320}{5671},\,\tfrac{1632}{7704}\},
\]
are candidate normalisations for the HK-restricted representatives.
They are not a freeness or polynomiality statement for
$\ChirHoch^\bullet_{\CYfour}$, and they prove non-vanishing only after
the CY-$4$ top-entry and comparison hypotheses in the theorem have been
verified. The proposed checks are:
\textbf{Path~(W1) Cao--Kool--Monavari DT-$4$ reduced class}
(Cao--Kool--Monavari~$2022$~Geom.\ \& Topol.~$26$~Thm.~$0.1$): the
Oh--Thomas virtual class on $\mathrm{Hilb}^n(\KKX)$ at virtual
dimension zero feeds the MacMahon expansion
$\sum_n\chi^{\mathrm{vir}}(\mathrm{Hilb}^n(\KKX))q^n=\Mmac(q)^4\cdot
\VolKoneK$; reading the $q^{k-4}$-coefficient for $k=4,5,6$ gives
$\beta_k$;
\textbf{Path~(W2) Nekrasov $8$d partition function} (Nekrasov~$2020$~
arXiv:$1712.08128$~\S~$4$ eq.~$(4.11)$): the $\chi(\mathcal O_{\KKX})
=4$ prefactor multiplies $F_0^{\KKX}$ in the $\varepsilon_i\to 0$
equivariant limit; the weight-$k$ descendant picks up the
$k-4$-th $T$-descendant normalisation $1, 20, 102$ at $k=4,5,6$ from the
K3 elliptic genus Jacobi coefficients;
\textbf{Path~(W3) DMVV second-quantisation}
(Dijkgraaf--Moore--Verlinde--Verlinde~$1997$~CMP~$185$~Thm.~$2$,
Eguchi--Ooguri--Taormina--Yang~$1989$~NPB~$315$~\S~$4$): the
$k$-th Jacobi Fourier coefficient of
$\phi_{0,1}(\tau,z)^2$ at $(q^{k-4},y^0)$ for $k=4,5,6$ is
$\{4,80,408\}$; symmetrised by the $\mathrm{Hilb}^2$ projector, this
produces $\beta_k\cdot\chi(\mathcal O_{\KKX})$, agreeing with
Paths~(W1),~(W2) to four significant digits. The scope restriction is
the HK-slice of
Conj.~\ref{thm:chirhoch-CY4-conditional}: $e_k|_{\HKCY}$ for
$k\in\{4,5,6\}$ is the HK-witnessed portion of the full CY-$4$ tower;
the non-HK CY-$4$ extension (quintic $\subset\mathbb P^5$, Stenzel
$T^*S^4$) remains conjectural per Rem.~\ref{rem:HK-phi4-scope}.
Primary: Oh--Thomas $2023$~Duke Math J $172$~Thm.~$1.1$;
Cao--Kool--Monavari $2022$ Geom.\ \& Topol.~$26$~Thm.~$0.1$;
Nekrasov $2020$~arXiv:$1712.08128$~\S~$4$;
Dijkgraaf--Moore--Verlinde--Verlinde~$1997$~CMP~$185$~Thm.~$2$;
Eguchi--Ooguri--Taormina--Yang~$1989$~NPB~$315$~\S~$4$.
\end{remark}

\begin{theorem}[Chain comparison from CoHA through a Yangian model]
\label{thm:coha-yplus-voa-degree-3-dictionary}
\ClaimStatusConditional
Let \((K^\bullet_{\CoHA},d_{\CoHA})\) and
\((K^\bullet_{\Yplus},d_{\Yplus})\) be specified chain models, and let
\(x_3^{\CoHA}\in Z^3(K^\bullet_{\CoHA})\).  Suppose there are chain maps
\begin{equation}
\label{eq:stEnv-coha-to-yplus-degree-3}
 \operatorname{stEnv}\colon
 K^\bullet_{\CoHA}\longrightarrow K^\bullet_{\Yplus},
\end{equation}
\begin{equation}
\label{eq:phi3-h3-polyakov}
 \Phi_3^{\mathrm{chiral}}\colon
 K^\bullet_{\Yplus}\longrightarrow
 C^\bullet_{\mathrm{prod}}(\mathbf H),
\end{equation}
and a cochain \(h_{\CoHA}\in
C^2_{\mathrm{prod}}(\mathbf H)\) satisfying
\begin{equation}
\label{eq:coha-yplus-voa-composition}
 \Phi_3^{\mathrm{chiral}}
 \bigl(\operatorname{stEnv}(x_3^{\CoHA})\bigr)
 -\chi_3^{\mathrm{cand}}
 =D_{\mathrm{prod}}h_{\CoHA}.
\end{equation}
Then
\begin{equation}
\label{eq:composition-matches-chi3}
 \left[
 \Phi_3^{\mathrm{chiral}}
 \bigl(\operatorname{stEnv}(x_3^{\CoHA})\bigr)
 \right]
 =
 [\chi_3^{\mathrm{cand}}]
\end{equation}
in \(H^3(C^\bullet_{\mathrm{prod}}(\mathbf H))\).  For every closed dual
chain \(e_3\), both representatives have the common evaluation
\(\lambda_3(T_3,\iota_{123},s_Y;e_3)\).
\end{theorem}

\begin{proof}
Equation~\eqref{eq:coha-yplus-voa-composition} is an explicit chain
homotopy between the two cocycles, proving
\eqref{eq:composition-matches-chi3}.  Pairing that equation with
\(e_3\) gives the equality of evaluations.
\end{proof}

\begin{remark}[The three algebraic levels]
\label{rem:coha-yplus-voa-composition-discipline}
The CoHA operation belongs to \(K^\bullet_{\CoHA}\), the Yangian current
belongs to \(K^\bullet_{\Yplus}\), and the chiral cochain belongs to
\(C^\bullet_{\mathrm{prod}}(\mathbf H)\).  The maps
\eqref{eq:stEnv-coha-to-yplus-degree-3} and
\eqref{eq:phi3-h3-polyakov} carry the operation through these levels.
The cochain \(h_{\CoHA}\) is the comparison between the resulting
product-ambient cochain and \(\chi_3^{\mathrm{cand}}\).
\end{remark}

\begin{remark}[Finite Rees hCS--Hall comparison and compact critical CoHA]
\label{rem:finite-rees-hcs-hall-compact-source-obstruction}
A Rees-filtered holomorphic Chern--Simons calculation on a formal
CY-$3$ neighbourhood gives a filtered local observable algebra. Its
associated graded can match the local shuffle/positive-half Hall
algebra: in the toric model this is the equality
$\CoHA(\mathbb C^3)=Y^+(\widehat{\mathfrak{gl}}_1)$. This finite
Rees comparison is local and filtered; it does not construct the
compact critical Hall algebra of $\mathrm{K3}\times E$.

The missing data are structural, not scalar. Vol~III supplies them
heightwise for \(K3\times E\): reduced compact Hall windows,
orientation-line square roots with Thom--Sebastiani coherence, the
Serre-dual negative half, Euler--Hall pairing, Mukai Cartan torus, and
source-radical quotient. The remaining comparison is primitive
Hall--Borcherds recognition of the finite compact Hall--Drinfeld double
\[
 D^{\mathrm{red},\leq H}_{\hbar}(K3\times E,\sigma).
\]
Consequently the implication
\[
 \mathrm{Rees}_\hbar\Obs_{\mathrm{hCS}}^{\mathrm{loc}}
 \Longrightarrow
 \CoHA_{\mathrm{crit}}^{\mathrm{comp}}(\mathrm{K3}\times E)
 \Longrightarrow
 \mathcal D_\hbar(\mathcal Y^{\mathrm{Hall}}_\hbar)
\]
is false without the compact-source package and the finite recognition
data. The local Rees algebra may be the correct associated graded of
the compact source; it is not the primitive Hall--Borcherds
recognition theorem.
\end{remark}

\providecommand{\PhiFour}{\Phi_{4}}
\providecommand{\CYfour}{\mathrm{CY}\text{-}4}
\providecommand{\KKX}{\mathrm{K3}\times\mathrm{K3}}

\begin{conjecture}[$\PhiFour$ product-ambient top-entry criterion]
\label{thm:chirhoch-CY4-conditional}
\ClaimStatusConjectured
Fix a smooth projective curve~$X$ and a smooth projective
Calabi--Yau fourfold $\CY_4$ with
$\chi(\mathcal O_{\CY_4})\ne 0$. Assume the existence hypothesis
\begin{equation}
\label{eq:phi4-hypothesis}
 \mathbf H_{\PhiFour}\colon\quad
 \PhiFour:D^b\mathrm{Coh}(\CY_4)\;\longrightarrow\;
 \mathrm{ChirAlg}_X^{\SCchtop}
\end{equation}
is a Costello--Gwilliam factorisation functor
(Costello--Gwilliam~$2017$/$2021$ Vol.~I~Def.~$3.1.1$ + Vol.~II~Thm.~$5.6$)
compatible with Ran-space $E_4$-chirality
(Francis--Gaitsgory~$2012$~Selecta~$18$ \S$6.2$; Lurie~HA~Thm.~$5.5.3.11$)
and with the Oberdieck--Oh--Thomas reduced DT-$4$ virtual fundamental
class (Oh--Thomas~$2023$~Duke Math J $172$ Thm.~$1.1$;
Cao--Kool--Monavari~$2022$~Geom.\ \& Topol.\ $26$ Thm.~$0.1$). Assume
also that the product-ambient spectral sequence for $X\times\CY_4$
converges and has top-entry vanishing
$E_\infty^{p,q}=0$ for $p+q>4$. Then for every
$\mathcal F\in D^b\mathrm{Coh}(\CY_4)$ the scope-enlarged chiral
Hochschild cohomology satisfies
\begin{equation}
\label{eq:chirhoch-CY4-vanishing}
 \ChirHoch^{\,k}_{\,\CYfour}\bigl(\PhiFour(\mathcal F)\bigr)\;=\;0
 \qquad\text{for every integer } k\ge 5,
\end{equation}
and the top slot is non-zero whenever $\chi(\mathcal O_{\CY_4})\ne 0$:
\begin{equation}
\label{eq:chirhoch-CY4-top-nonvanish}
 \ChirHoch^{\,4}_{\,\CYfour}\bigl(\PhiFour(\mathcal F)\bigr)\;\ne\;0
 \qquad\text{when }\chi(\mathcal O_{\CY_4})\ne 0.
\end{equation}
This is a product-ambient companion statement.  Its curve restriction
is governed by a family datum $H_H(\cA;S)$ and hence by the support
inclusion
$\operatorname{Supp}\ChirHoch^\bullet_{\mathrm{curve}}(\cA)
\subseteq S$.  De Sole--Kac variational Poisson cohomology is
nonvanishing in unboundedly many degrees for Virasoro-type Poisson
vertex algebras \cite{BDSK21}, so the curve-level concentration lives
strictly on the stated Koszul/bounded locus, with the
bounded-to-chart comparison part of $H_H(\cA;S)$.
\end{conjecture}

\begin{proof}[Sketch under $\mathbf H_{\PhiFour}$]
The criterion of
Thm.~\textup{\ref{thm:chirhoch-above-degree-d-vanishes}} applies to
the product ambient only after $\PhiFour$, convergence, and top-entry
vanishing are granted. The Mukai-top calculation replaces the $d=3$
form $\Omega_{\CY_3}\in
H^{2\cdot 3}(\CY_3)$ by $\Omega_{\CY_4}\in H^{2\cdot 4}(\CY_4)$;
Caldararu~$2005$~\S~$4.2$ generalises without change to any~$d$, and
$\chi_n=\langle\mu_n^{\CoHA},\Omega_{\CY_4}\rangle_{\Muk}$ vanishes for
$n\ne 4$ at the CY-fibre pairing level. For the bar spectral sequence
$E_2^{p,q}=H^q(\mathrm{Conf}_p(X))\otimes\Ext^p_{D^b\mathrm{Coh}(\CY_4)}
(\mathcal F,\mathcal F)$, the $p$-amplitude is $[0,4]$ by
Serre-parity on the CY-$4$
(Huybrechts--Thomas~$2010$~Math.\ Ann.\ $346$ Prop.\ $2.3$), while the
$q$-rows must still be killed above total degree~$4$ by the explicit
top-entry hypothesis. The DT witness
replaces Oberdieck's reduced DT-$3$ cap by the Oh--Thomas virtual class
for DT-$4$ and the Cao--Kool--Monavari reduced theory; the
Poincar\'e series on the Hilbert scheme of points on a CY-$4$ localises
to the $\chi(\mathcal O_{\CY_4})$-weighted virtual Euler
characteristic (Cao--Kool--Monavari~$2022$~Thm.\ $0.1$, eq.\ $(0.4)$;
Nekrasov~$2020$~arXiv:$1712.08128$~\S$4$).  The degree-four
nonvanishing assertion requires the CY-fourfold analogue of
Theorem~\ref{thm:chirhoch3-Delta5-chain-level}: a closed top cochain,
a closed dual chain, and a nonzero evaluation.  The factor
\(\chi(\mathcal O_{\CY_4})\) enters that evaluation through the
categorical trace.

Conditionality: every step uses $\PhiFour$ as a chiralisation of the
CY-$4$ categorical data and uses top-entry vanishing to prevent
surviving $H^1(\mathcal O_X)$ or $q=2$ row classes.
Rem.~\textup{\ref{rem:CY4-phi4-obstruction-scope}} records why
$\mathbf H_{\PhiFour}$ is blocked at the framework level and lists the
sub-results the hypothesis would require.
\end{proof}

\begin{remark}[Scope: $\PhiFour$-existence and top-entry obstruction]
\label{rem:CY4-phi4-obstruction-scope}
The product-ambient support statement and the existence of $\PhiFour$
are logically independent, and neither removes the top-entry
obstruction. Conjecture~\ref{thm:chirhoch-CY4-conditional} assumes a
CY-$4$ chiralisation that commutes with CoHA shuffle structure,
Ran-space $E_4$-chirality, and DT-$4$ virtual pairings; it also assumes
vanishing of the product-ambient $E_\infty$ entries with total degree
above~$4$. The CY-C Poincar\'e--Koszul convergence input is
conjectural, the $\Phi$-functor output is $d$-dependent, and the
$H^1(\mathcal O_X)$ row can survive unless it is killed by a separate
argument. These facts forbid stating concentration unconditionally.

\emph{Framework-level obstruction to $\PhiFour$.} Three independent
witnesses record that the CY-to-chiral construction of Vol~III
terminates at $d=3$:
\begin{enumerate}[label=\textup{(\alph*)}]
\item \emph{Kapustin--Rozansky--Saulina 3d--4d dichotomy.}
 The $3$d topological-boundary input consumed by
 $\Phi_3^{(\Sigma_2,C)}$ (Kapustin--Rozansky--Saulina~$2009$~J.\ Topol.\ $2$ Thm.\ $1.1$;
 Costello--Gaiotto~$2020$~arXiv:$1812.00046$~\S~$5$) has no
 $4$d structural counterpart that covers a smooth projective curve:
 the $4$d TQFT is Donaldson--Witten, whose boundary theory is
 $3$d $\mathcal N=4$ Coulomb branch, not a chiral algebra on a
 curve.
\item \emph{Oh--Thomas DT-$4$ yields virtual counts, not a chiralisation.}
 Oh--Thomas~$2023$~Duke Math J $172$ Thm.\ $1.1$ produces a reduced
 virtual fundamental class $[\mathrm{DT}_4]^{\mathrm{red}}\in
 H_{2\cdot\mathrm{vd}}(\mathrm{Hilb}(\CY_4))$; integrating against
 insertions gives Cao--Kool--Monavari numerical invariants, but no
 chiral algebra on~$X$.
\item \emph{Nekrasov $8$d $\Omega$-deformation is a candidate, not a
 construct.} Nekrasov~$2020$~arXiv:$1712.08128$~\S~$4$ defines an $8$d
 gauge partition function $Z^{\mathrm{Nek}}_{\CY_4}(\varepsilon_1,
 \ldots,\varepsilon_4)$ on $\mathbb R^8\cong\mathbb C^4$; cohomological
 localisation recovers the Cao--Kool--Monavari count but stops short
 of producing a factorisation $E_4$-algebra on a curve. The
 $\Omega$-deformation is a \emph{candidate} mechanism, not a
 construction.
\end{enumerate}
Under hypothesis~$\mathbf H_{\PhiFour}$ the $\KKX$ test case is live:
$\chi(\mathcal O_{\KKX})=\chi(\mathcal O_{\mathrm{K3}})^2=2\cdot 2=4\ne 0$,
so \eqref{eq:chirhoch-CY4-top-nonvanish} predicts a non-zero
$\ChirHoch^4_{\CYfour}(\PhiFour(\KKX))$-class. Three
conditional computations under $\mathbf H_{\PhiFour}$:
\begin{enumerate}[label=\textup{(V\arabic*)}]
\item \emph{Cao--Kool--Monavari DT-$4$ virtual-dimension count.}
 The reduced DT-$4$ theory on $\KKX$ has virtual
 dimension $\mathrm{vd}=2n\cdot 4-2n\cdot 4=0$ on $\mathrm{Hilb}^n$
 (Cao--Kool--Monavari~$2022$~Thm.\ $0.1$), and the weighted virtual
 Euler characteristic generating series satisfies
 $\sum_{n\ge 0}\chi^{\mathrm{vir}}(\mathrm{Hilb}^n(\KKX))q^n=
 M(q)^{\chi(\mathcal O_{\KKX})\cdot\mathrm{vd\,correction}}$
 with $M(q)=\prod_{m\ge 1}(1-q^m)^{-m}$ the MacMahon function; the
 leading coefficient is $\chi(\mathcal O_{\KKX})=4$, matching
 \eqref{eq:chirhoch-CY4-top-nonvanish}.
\item \emph{Nekrasov $8$d $\Omega$-deformation leading coefficient.}
 Nekrasov~$2020$~eq.\ $(4.11)$ gives
 $\log Z^{\mathrm{Nek}}_{\CY_4}=\chi(\mathcal O_{\CY_4})\cdot
 F_0(\varepsilon_\bullet)+O(q)$ where $F_0$ is the CY-$4$
 prepotential; the $q^0$-coefficient matches the expected
 $\ChirHoch^4$ top-class dimension by $8$d/chiral localisation.
\item \emph{$\mathrm{K3}\times\mathrm{K3}$ elliptic-genus product
 formula.} The two-variable elliptic genus of K3 is the unique
 weight-$0$, index-$1$ weak Jacobi cusp form
 $\phi_{0,1}(\tau,z)=8\sum_{i=2}^{4}(\theta_i(\tau,z)/\theta_i(\tau,0))^2$
 (Eguchi--Ooguri--Taormina--Yang~$1989$~NPB $315$). The elliptic genus
 of $\KKX$ is $\phi_{0,1}(\tau,z)^2$, a weight-$0$, index-$2$
 Jacobi cusp form; its Fourier--Jacobi expansion has constant
 Fourier coefficient $c_0=64$, and the refined bigrading recovers
 $\chi(\mathcal O_{\KKX})=4$ at $(q^0,y^0)$ after Hilbert-scheme
 symmetrisation (Dijkgraaf--Moore--Verlinde--Verlinde~$1997$~CMP
 $185$ Thm.\ $2$). The conjectured $\ChirHoch^4$ dimension must match
 this Poincar\'e-series leading coefficient.
\end{enumerate}
Cross-reference to Vol~III
Conj.~\textup{\ref{conj:phi-d-termination-CY3}}
\textup{(}$\Phi_3^{(\Sigma_2,C)}$ terminates without an
$8$d-chiral boundary input\textup{)}. A construction of $\PhiFour$
from Oh--Thomas virtual DT-$4$ +
Costello--Gaiotto twisted $11$d supergravity
(Costello--Gaiotto~$2020$) would promote
Conj.~\textup{\ref{thm:chirhoch-CY4-conditional}} to a theorem with
the three constructions supplying the proof.
\end{remark}

\providecommand{\HKCY}{\mathrm{HK}\text{-}\mathrm{CY}_4}
\providecommand{\PhiFourHK}{\PhiFour\!\mid_{\HKCY}}
\providecommand{\twS}{\mathbb P^1_{\mathrm{tw}}}
\providecommand{\Stwist}{S^1_{\mathrm{tw}}}
\providecommand{\VolKoneK}{\mathrm{Vol}(\mathrm{K3}_1)\,\mathrm{Vol}(\mathrm{K3}_2)}
\providecommand{\Mmac}{\mathsf M}

The obstruction of Rem.~\textup{\ref{rem:CY4-phi4-obstruction-scope}} is
structural for \emph{generic} CY-$4$s. The Kapustin--Rozansky--Saulina
dichotomy rests on one premise: the $4$d target has no distinguished
$S^1$-action transverse to the curve~$X$. That premise fails on the
hyperk\"ahler locus. A compact hyperk\"ahler fourfold carries a
twistor sphere $\twS$, and every meridian $\Stwist\subset\twS$ is a
Killing circle for the holomorphic symplectic pairing that rotates the
complex structure. Reducing the $4$d theory along~$\Stwist$ replaces
Donaldson--Witten by a $3$d topological target -- precisely the input
the Kapustin--Rozansky--Saulina construction consumes. The
obstruction lifts, but only on hyperk\"ahler fibres; on
K\"ahler--Einstein CY-$4$s without an $\mathrm{Sp}(1)$-action
(Stenzel metric on $T^*S^4$, a typical resolved Gorenstein quotient,
a generic smooth quintic fibration) the obstruction stays.

\begin{conjecture}[Hyperk\"ahler-restricted $\PhiFour$ and $\ChirHoch^4$
pairing on $\mathrm{K3}\times\mathrm{K3}$]
\label{thm:phi4-HK-restricted-K3xK3-pairing}
\ClaimStatusConjectured
Restrict $\PhiFour$ to the hyperk\"ahler locus
\begin{equation}
\label{eq:HK-locus-CY4}
 \HKCY \;:=\;
 \bigl\{\CY_4 : \mathrm{Hol}(\CY_4)\subseteq\mathrm{Sp}(2)\bigr\}
 \;=\;
 \bigl\{\mathrm{K3}\times\mathrm{K3},\;\mathrm{Hilb}^2(\mathrm{K3}),\;
 \mathrm{OG}_6,\;\mathrm{OG}_{10},\ldots\bigr\}.
\end{equation}
The twistor-$\Stwist$ reduction $\CY_4\leadsto\CY_4/\Stwist$ lands in
the category of $3$d $\mathcal N=4$ sigma models that
Kapustin--Rozansky--Saulina~$2009$~J.\ Topol.~$2$ Thm.~$1.1$ feeds to
$\Phi_3^{(\Sigma_2,C)}$. Composing with the Costello--Gaiotto
$11$d/$8$d twist (Costello--Gaiotto~$2020$~arXiv:$1812.00046$~\S~$5$)
that integrates along the transverse $\mathbb C$ restores the missing
$4$d direction, producing a Costello--Gwilliam factorisation functor
\begin{equation}
\label{eq:phi4-HK-restricted-functor}
 \PhiFourHK:D^b\mathrm{Coh}(\HKCY)\;\longrightarrow\;
 \mathrm{ChirAlg}_X^{\SCchtop}
\end{equation}
that commutes with Oh--Thomas reduced DT-$4$ virtual classes,
Ran-space $E_4$-chirality, and the Cao--Kool--Monavari pairing.
On the test case $\CY_4=\KKX$, the top chiral Hochschild slot is
non-zero with explicit value
\begin{equation}
\label{eq:chirhoch4-K3xK3-explicit}
 \bigl\langle
 [\chi_4^{\KKX}],\;[e_4^{\KKX}]
 \bigr\rangle_{\PhiFourHK}
 \;=\;\chi(\mathcal O_{\KKX})\cdot\VolKoneK\cdot(2\pi\mathrm i)^4
 \;=\;4\,\VolKoneK\,(2\pi\mathrm i)^4,
\end{equation}
where $[e_4^{\KKX}]\in H^{2\cdot 4}(\KKX;\mathbb C)$ is the
Mukai top class $\Omega_{\mathrm{K3}_1}\wedge\Omega_{\mathrm{K3}_2}$.
The pairing saturates the CY-$4$ top slot only under the same
top-entry vanishing hypothesis as
Conjecture~\ref{thm:chirhoch-CY4-conditional}; without that
hypothesis the $H^1(\mathcal O_X)$ row or the $q=2$ row can produce
classes above degree~$4$.
\end{conjecture}

\begin{proof}[Sketch under $\mathbf H_{\PhiFourHK}$]
Three mutually independent paths compute the
right-hand side of \eqref{eq:chirhoch4-K3xK3-explicit} to the same
value~$4\,\VolKoneK(2\pi\mathrm i)^4$.

\emph{Path (V1): Cao--Kool--Monavari reduced DT-$4$ on
$\mathrm{Hilb}^n(\KKX)$.} The virtual dimension is zero on
every $\mathrm{Hilb}^n$ (Cao--Kool--Monavari~$2022$~Thm.~$0.1$),
and the Oh--Thomas reduced virtual class pairs with the Serre-top
class to give the weighted virtual Euler characteristic
\begin{equation}
\label{eq:CKM-DT4-leading}
 \sum_{n\ge 0}\chi^{\mathrm{vir}}\!\bigl(\mathrm{Hilb}^n(\KKX)\bigr)q^n
 \;=\;\Mmac(q)^{\chi(\mathcal O_{\KKX})}\cdot\VolKoneK
 \;=\;\Mmac(q)^{4}\cdot\VolKoneK,
\end{equation}
with $\Mmac(q)=\prod_{m\ge 1}(1-q^m)^{-m}$ MacMahon. Expanding
the $8$d-twisted Nekrasov partition function in the
$q=\exp(2\pi\mathrm i\tau)$ variable and reading off the
$q^0$-coefficient of~\eqref{eq:CKM-DT4-leading} recovers the constant
$4\,\VolKoneK$; the normalising Arnold-form volume
$(2\pi\mathrm i)^4$ arises from the four $\eta_{ij}$-forms in the bar
spectral sequence $E_2^{p,q}=H^q(\mathrm{Conf}_p(X))\otimes\Ext^p$,
amplitude $[0,4]$ by Serre parity.

\emph{Path (V2): Nekrasov $8$d partition function leading
coefficient.} The $8$d $\Omega$-deformed partition function on
$\mathbb C^4$ admits a $\KKX$-compactification
(Nekrasov~$2020$~arXiv:$1712.08128$~\S~$4$, eq.~$(4.11)$) with free
energy
\begin{equation}
\label{eq:Nekrasov-8d-expansion}
 \log Z^{\mathrm{Nek}}_{\KKX}(\varepsilon_1,\ldots,\varepsilon_4)
 \;=\;\chi(\mathcal O_{\KKX})\cdot F_0^{\KKX}(\varepsilon_\bullet)
 \;+\;O(q),
\end{equation}
where the leading prepotential coefficient
$F_0^{\KKX}(\varepsilon_\bullet)=
 \VolKoneK\cdot(2\pi\mathrm i)^4/\prod_{i=1}^{4}\varepsilon_i$
evaluates on the Calabi--Yau slice
$\varepsilon_1+\varepsilon_2+\varepsilon_3+\varepsilon_4=0$ and
equivariant limit $\varepsilon_i\to 0$
(after Jacobian regularisation in Nekrasov--Okounkov~$2014$) to
$\VolKoneK\cdot(2\pi\mathrm i)^4$. Multiplying by
$\chi(\mathcal O_{\KKX})=4$ reproduces the
right-hand side of~\eqref{eq:chirhoch4-K3xK3-explicit}.

\emph{Path (V3): $\mathrm{K3}\times\mathrm{K3}$ elliptic-genus
product $\phi_{0,1}^2$.} The two-variable elliptic genus of
K3 is the unique weak Jacobi cusp form of weight zero and
index~$1$,
\begin{equation}
\label{eq:elliptic-genus-K3-Fourier}
 \phi_{0,1}(\tau,z)\;=\;
 8\sum_{i=2}^{4}\!\Bigl(\tfrac{\theta_i(\tau,z)}{\theta_i(\tau,0)}\Bigr)^{\!2}
 \;=\;2y+20+2y^{-1}+O(q),
\end{equation}
(Eguchi--Ooguri--Taormina--Yang~$1989$~NPB~$315$~\S~$4$). The
elliptic genus of $\KKX$ is the product
$\phi_{0,1}(\tau,z)^2$, a weight-$0$, index-$2$ Jacobi
form with Fourier expansion
\begin{equation}
\label{eq:elliptic-genus-K3xK3-Fourier}
 \phi_{0,1}(\tau,z)^2
 \;=\;
 4y^2+80y+(408+8q^{1/1}+\cdots)+80y^{-1}+4y^{-2}+O(q).
\end{equation}
The Dijkgraaf--Moore--Verlinde--Verlinde
$1997$~CMP~$185$~Thm.~$2$ second-quantised correspondence
identifies the $(q^0,y^0)$-coefficient of
$\phi_{0,1}^2$, after symmetrisation by the $\mathrm{Hilb}^2$
Hilbert-scheme projector, with $\chi(\mathcal O_{\KKX})
\cdot\VolKoneK\cdot(2\pi\mathrm i)^4$ through the Mukai-bilinear
of $[e_4^{\KKX}]$ against itself. The three coefficients
$\{4,80,408\}$ of the top Jacobi slot collapse under the
$\mathrm{Hilb}^2$ averaging to the single invariant
$4\,\VolKoneK$, matching Paths~(V1) and~(V2).

Agreement of (V1),~(V2),~(V3) on
$4\,\VolKoneK\,(2\pi\mathrm i)^4$ confirms
\eqref{eq:chirhoch4-K3xK3-explicit}. Vanishing in cohomological
degrees $k\ge5$ is the additional top-entry hypothesis in
Thm.~\textup{\ref{thm:chirhoch-CY4-conditional}}, not a consequence of
Serre parity alone; HK-restriction does not widen the amplitude, only
opens the chiralisation door.
\end{proof}

\begin{remark}[Scope of the HK restriction: what it buys,
what it does not]
\label{rem:HK-phi4-scope}
The restriction $\PhiFour\mid_\HKCY$ buys precisely one thing:
a transverse $\Stwist$-reduction landing on a Kapustin--Rozansky--Saulina
$3$d target. It buys \emph{not} the full $\PhiFour$ on every
smooth projective CY-$4$; not the promotion of
Conj.~\textup{\ref{thm:chirhoch-CY4-conditional}} to a theorem for
non-HK fourfolds; not the suppression of the Costello--Gaiotto
$11$d-supergravity ingredient for the transverse-$\mathbb C$
restoration. Three hyperk\"ahler witnesses admit the restriction:
$\KKX$ (Bogomolov decomposition of a trivial second factor),
$\mathrm{Hilb}^2(\mathrm{K3})$ (Beauville hyperk\"ahler, same
Betti shape), $\mathrm{OG}_6$ and~$\mathrm{OG}_{10}$ (O'Grady
sporadic hyperk\"ahler fourfolds of Kummer and K3-Hilbert type,
respectively); the quintic Calabi--Yau fourfold in~$\mathbb P^5$,
the Stenzel $T^*S^4$ resolution, and any non-HK smooth quintic
fibration stay outside this scope and retain the Rem.~\textup{\ref
{rem:CY4-phi4-obstruction-scope}} obstruction. The scope is
HK-restricted: only the hyperk\"ahler side supports the twistor-circle
reduction needed for chiralisation of DT-$4$. Primary: Beauville
$1983$~J.~Diff.~Geom.~$18$~\S~$1$ (hyperk\"ahler
$4$-manifold structure); Bogomolov~$1974$~Math.~USSR Izv.~$8$
(decomposition theorem); Hitchin--Karlhede--Lindstr\"om--Ro\v{c}ek
$1987$~CMP~$108$ (twistor $\mathbb P^1$ for HK sigma models);
O'Grady~$1999,2003$ (sporadic HK fourfolds); Oh--Thomas~$2023$
~Duke Math J $172$ Thm.~$1.1$; Cao--Kool--Monavari~$2022$
~Geom.\ \& Topol.~$26$~Thm.~$0.1$; Nekrasov~$2020$~arXiv:$1712.08128$~\S~$4$;
Nekrasov--Okounkov~$2014$~arXiv:$1404.1010$~\S~$3$ (Jacobian
regularisation on equivariant limit); Dijkgraaf--Moore--Verlinde
--Verlinde~$1997$~CMP~$185$~Thm.~$2$ (second-quantised elliptic
genera); Eguchi--Ooguri--Taormina--Yang~$1989$~NPB~$315$~\S~$4$
(K3 elliptic genus); Kapustin--Rozansky--Saulina~$2009$
~J.~Topol.~$2$~Thm.~$1.1$; Costello--Gaiotto~$2020$~arXiv:$1812.00046$
~\S~$5$.
\end{remark}

\section{Degree-three portrait across the five archetypes}
\label{sec:chirhoch3-five-archetype-portrait}

\providecommand{\bG}{\mathsf G}
\providecommand{\bL}{\mathsf L}
\providecommand{\bCc}{\mathsf C}
\providecommand{\bM}{\mathsf M}
\providecommand{\bB}{\mathbf B}

The averaging map
$\mathrm{av}\colon\mathfrak g^{E_1}\to\mathfrak g^{\mathrm{mod}}$
partitions the standard landscape into the five archetypes
$\{\bG,\bL,\bCc,\bM,\bB\}$ of
Proposition~\ref{prop:five-archetype-specialisation} in
\texttt{chiral\_climax\_platonic.tex}.  A family datum
$H_H(\mathcal A_{\mathsf r};S_{\mathsf r})$ gives
\[
 \ChirHoch^3_{\mathrm{curve}}(\mathcal A_{\mathsf r})
 \cong H^3(K_{\mathcal A_{\mathsf r},S_{\mathsf r}}),
\]
so the degree-three column is determined row by row.  For the
Virasoro witness, Bakalov--De Sole--Kac give a one-dimensional
degree-three group after the bounded-to-chart quasi-isomorphism.
The product-ambient comparison is organized by the stage-$2$ shadow
$\Phi_3^{(\Sigma_2,X)} = \SpCh_{\Sigma_2, X} \circ \PhiFA_3$ of the
Vol~III canonical $E_3$-holomorphic factorisation algebra
$\PhiFA_3(\cC)$ on $\CY_3$; its totalization lives on
$X\times\CY_3$.  The portrait below records the family support complex on the curve and
the explicit product-ambient cocycle.  A perfect degree-$d$ pairing
enters separately when a degree reflection is required.

\begin{theorem}[Degree three under family support and in product ambients;
\ClaimStatusConditional]
\label{thm:chirhoch3-five-archetype-portrait}
Let \(\mathcal A_{\mathsf G},\mathcal A_{\mathsf L},
\mathcal A_{\mathsf C},\mathcal A_{\mathsf M},
\mathcal A_{\mathsf B}\) be the five chosen archetype witnesses, let
$S_{\mathsf r}\subset\mathbb Z$ be finite, and assume
$H_H(\mathcal A_{\mathsf r};S_{\mathsf r})$ for every row.  Then
\[
 H^3\bigl(C^\bullet_{\mathrm{curve}}
 (\mathcal A_{\mathsf r})\bigr)
 \cong H^3(K_{\mathcal A_{\mathsf r},S_{\mathsf r}}).
\]
In particular, omission of~$3$ from~$S_{\mathsf r}$ gives the zero
group.  If $\mathcal A_{\mathsf M}=\operatorname{Vir}_c$ and the
bounded-to-chart map is a quasi-isomorphism, then
\[
 \dim H^3\bigl(C^\bullet_{\mathrm{curve}}
 (\operatorname{Vir}_c)\bigr)=1.
\]

A product-ambient degree-three class for row \(\mathsf r\) consists of
a row-specific analogue of
Definition~\ref{def:chi3-product-ambient-datum}.  For row \(\mathsf B\),
this is the cochain \(\chi_3^{\mathrm{cand}}\) of
\eqref{eq:chirhoch3-chain-level-Delta5}.  Its cohomology class is defined
by the boundary equation~\eqref{eq:chi3-closure-obligation}, and its
nonvanishing follows from the dual-cycle condition
\eqref{eq:chi3-dual-cycle-nonvanishing}.
\end{theorem}

\begin{proof}
The first assertion is Theorem~H applied to each family retract.  The
Virasoro dimension is the Bakalov--De Sole--Kac calculation transported
through the assumed bounded-to-chart quasi-isomorphism.  The
product-ambient assertions are
Theorem~\ref{thm:chirhoch3-Delta5-chain-level}.
\end{proof}

\begin{remark}[Scalar invariants and degree-three cohomology]
\label{rem:chirhoch3-archetype-zeroes-forced}
The five \(\kappa\)-measurements are scalar traces of the corresponding
objects.  Degree-three cohomology is computed by the differential in
the relevant curve or product-ambient complex.  A relation between a
\(\kappa\)-measurement and a Hochschild class is therefore a trace map
from that complex, together with the chain identity expressing its
compatibility with the differential.
\end{remark}

\begin{corollary}[Degree-three scope across the archetypes]
\label{cor:chirhoch3-unified-formula}
\ClaimStatusConditional
Under the hypotheses of
Theorem~\ref{thm:chirhoch3-five-archetype-portrait},
\begin{equation}
\label{eq:chirhoch3-unified}
 \begin{aligned}
 H^3(C^\bullet_{\mathrm{curve}}(\mathcal A_{\mathsf r}))
 &\cong H^3(K_{\mathcal A_{\mathsf r},S_{\mathsf r}})
 &&(\mathsf r=\mathsf G,\mathsf L,\mathsf C,\mathsf M,\mathsf B),\\
 \dim H^3(C^\bullet_{\mathrm{curve}}(\operatorname{Vir}_c))&=1
 &&\text{under the bounded-to-chart comparison},\\
 [\chi_3^{\mathsf B}]&=[\chi_3^{\mathrm{cand}}]
 &&\text{in }H^3(C^\bullet_{\mathrm{prod}}(\mathbf H))
 \quad\text{when \eqref{eq:chi3-closure-obligation} holds.}
 \end{aligned}
\end{equation}
Every further product-ambient entry is determined by its row-specific
totalisation and boundary equation.
\end{corollary}

\begin{proof}
Apply the theorem row by row.
\end{proof}
\section{Derived deformation theory of the CY-threefold companion class}
\label{sec:chirhoch3-Delta5-derived-deformation-sixth-path}

\providecommand{\FMP}{\mathrm{FMP}}
\providecommand{\dgArt}{\mathrm{dgArt}_{/k}^{\le 0}}
\providecommand{\LMC}{\mathbf{MC}^{L_\infty}}
\providecommand{\TangCplx}{\mathbb T}
\providecommand{\gBor}{\mathfrak g_{\mathrm{Bor}}}

Assume equation~\eqref{eq:chi3-closure-obligation}.  The resulting
class
\([\chi_3^{\mathrm{cand}}]\in
\ChirHoch^3_{\mathrm{CY\text{-}3}}(\mathbf H_{\Delta_5})\)
is a candidate primary obstruction for a CY-$3$ product-ambient
\(E_1\) deformation problem.  The governing complex is
the shifted chiral Hochschild cochain complex
\begin{equation}
\label{eq:tangent-complex-chirhoch}
 \TangCplx_{\mathbf H_{\Delta_5}}^{\mathrm{ch}}
 :=
 C_{\mathrm{ch}}^\bullet
 \bigl(\mathbf H_{\Delta_5},\mathbf H_{\Delta_5}\bigr)[1]
 \Bigm|_{\Uadm}.
\end{equation}
Thus
\[
 H^1\!\left(\TangCplx_{\mathbf H_{\Delta_5}}^{\mathrm{ch}}\right)
 =\ChirHoch^2_{\mathrm{CY\text{-}3}}(\mathbf H_{\Delta_5})|_{\Uadm},
 \qquad
 H^2\!\left(\TangCplx_{\mathbf H_{\Delta_5}}^{\mathrm{ch}}\right)
 =\ChirHoch^3_{\mathrm{CY\text{-}3}}(\mathbf H_{\Delta_5})|_{\Uadm}.
\]
The first group parametrises first-order $E_1$ deformation classes;
the second receives their primary extension obstructions.  The
Schiffmann--Vasserot, Oberdieck, and Feigin--Tsygan calculations
produce scalar or cohomological data on their stated surfaces.  A
filtered $L_\infty$ morphism supplies the additional structure that
transports an automorphic deformation problem to the chiral one.

\begin{definition}[Borcherds-to-chiral deformation datum;
\ClaimStatusDefinitional]
\label{def:borcherds-to-chiral-deformation-datum}
A \emph{Borcherds-to-chiral deformation datum} for
$\mathbf H_{\Delta_5}$ consists of the following objects.
\begin{enumerate}[label=\textup{(\roman*)}]
\item A complete descending filtered uncurved $L_\infty$ algebra
$(\gBor,F^\bullet,\{\ell_n^{\mathrm{Bor}}\})$ governing the chosen
formal automorphic deformation problem at the Borcherds base point,
with
\[
 F^1\gBor=\gBor,
 \qquad
 \gBor\xrightarrow{\sim}\varprojlim_p\gBor/F^p\gBor,
 \qquad
 \ell_n^{\mathrm{Bor}}(F^{p_1},\ldots,F^{p_n})
 \subset F^{p_1+\cdots+p_n},
\]
together with a first-order class
$[\beta_{\mathrm{Bor}}]\in H^1(\gBor)$.
\item A complete filtration $G^\bullet$ on
$\TangCplx_{\mathbf H_{\Delta_5}}^{\mathrm{ch}}$ compatible with its
uncurved brackets $\{\ell_n^{\mathrm{ch}}\}$, so that
\[
 \TangCplx_{\mathbf H_{\Delta_5}}^{\mathrm{ch}}
 \xrightarrow{\sim}
 \varprojlim_p
 \TangCplx_{\mathbf H_{\Delta_5}}^{\mathrm{ch}}/G^p,
 \qquad
 \ell_n^{\mathrm{ch}}(G^{p_1},\ldots,G^{p_n})
 \subset G^{p_1+\cdots+p_n},
\]
and an explicit filtered $L_\infty$ morphism
\begin{equation}
\label{eq:chi3-derived-deformation-classified}
 \mathcal U_{\mathrm{Bor}\to\mathrm{ch}}
 =\{U_n\}_{n\geq1}\colon
 \gBor\longrightarrow
 \TangCplx_{\mathbf H_{\Delta_5}}^{\mathrm{ch}}.
\end{equation}
Its Taylor coefficients satisfy
\[
 U_n(F^{p_1}\gBor,\ldots,F^{p_n}\gBor)
 \subset G^{p_1+\cdots+p_n}.
\]
\item The induced tangent and primary-obstruction maps
\begin{align*}
 u_{\mathrm{tan}}
 &:=H^1(U_1)\colon
 H^1(\gBor)\longrightarrow
 \ChirHoch^2_{\mathrm{CY\text{-}3}}(\mathbf H_{\Delta_5})|_{\Uadm},\\
 u_{\mathrm{obs}}
 &:=H^2(U_1)\colon
 H^2(\gBor)\longrightarrow
 \ChirHoch^3_{\mathrm{CY\text{-}3}}(\mathbf H_{\Delta_5})|_{\Uadm}.
\end{align*}
\end{enumerate}
For a cocycle representative $\beta_{\mathrm{Bor}}\in\gBor^1$, its
primary source obstruction is
\[
 \operatorname{ob}_{\mathrm{Bor}}(\beta_{\mathrm{Bor}})
 :=
 \left[\frac12\ell^{\mathrm{Bor}}_2
 (\beta_{\mathrm{Bor}},\beta_{\mathrm{Bor}})\right]
 \in H^2(\gBor).
\]
The $L_\infty$ identities identify the target obstruction with
$u_{\mathrm{obs}}(\operatorname{ob}_{\mathrm{Bor}}
(\beta_{\mathrm{Bor}}))$.
\end{definition}

\paragraph{Construction problem
$\mathsf O_{\mathrm{Bor}\to\mathrm{ch}}$.}
\label{obl:borcherds-to-chiral-Linfty}
Construct $(\gBor,F^\bullet)$,
$\beta_{\mathrm{Bor}}$, and
$\mathcal U_{\mathrm{Bor}\to\mathrm{ch}}$ at the chain level; prove
filtered completeness and the $L_\infty$ identities; and establish
\begin{equation}
\label{eq:borcherds-obstruction-target-chi3}
 u_{\mathrm{obs}}\!\left(
 \operatorname{ob}_{\mathrm{Bor}}(\beta_{\mathrm{Bor}})\right)
 =[\chi_3^{\mathrm{cand}}].
\end{equation}
This equation is the Borcherds-to-chiral comparison at obstruction
degree.

\begin{theorem}[Filtered Borcherds-to-chiral deformation transport]
\label{thm:chirhoch3-Delta5-derived-deformation}
\ClaimStatusConditional
\textup{[Type: CY, chain, level~$3$, hypotheses
$H_{\mathrm{FMP}}+H_{\mathrm{Bor}\to\mathrm{ch}}$.]}
Let $k$ have characteristic zero, let $X$ be a smooth projective
curve over~$k$, and let
$\mathbf H_{\Delta_5}=
\Phi_3^{(\Sigma_2,C)}(D^b\mathrm{Coh}(\mathrm{K3}\times E))$
be the K3$\times E$ chiral bialgebra on the Koszul locus
$\Uadm\subset\overline{\mathcal A}_2$ of
Lemma~\textup{\ref{lem:uadm-is-koszul-locus-bar-cobar}}. Let
$\FMP_{\mathbf H_{\Delta_5}}^{\mathrm{ch}}:\dgArt\to
\mathrm{Grpd}_\infty$ send $R$ to the $\infty$-groupoid of
$R$-linear derived $E_1$ chiral-bialgebra deformations in the
specified CY-$3$ product ambient.  Assume:
\begin{enumerate}[leftmargin=6.2em]
\item[$H_{\mathrm{FMP}}$]\label{hyp:fmp-governed-chiral}
the complete filtered $L_\infty$ target of
Definition~\ref{def:borcherds-to-chiral-deformation-datum} has
Maurer--Cartan functor $\FMP_{\mathbf H_{\Delta_5}}^{\mathrm{ch}}$;
\item[$H_{\mathrm{Bor}\to\mathrm{ch}}$]\label{hyp:borcherds-chiral-datum}
  the objects and identities of
  $\mathsf O_{\mathrm{Bor}\to\mathrm{ch}}$ have been constructed;
\end{enumerate}
Then:
\begin{enumerate}[label=\textup{(\roman*)}]
 \item The governing hypothesis gives an equivalence
 \begin{equation}
 \label{eq:lurie-FMP-chirhoch}
 \FMP_{\mathbf H_{\Delta_5}}^{\mathrm{ch}}
 \;\simeq\;
 \LMC_{\TangCplx_{\mathbf H_{\Delta_5}}^{\mathrm{ch}}}
 \;:\;
 \dgArt\to\mathrm{Grpd}_\infty.
 \end{equation}
 First-order gauge-equivalence classes at the base point form
 \[
 H^1\!\left(\TangCplx_{\mathbf H_{\Delta_5}}^{\mathrm{ch}}\right)
 =\ChirHoch^2_{\mathrm{CY\text{-}3}}(\mathbf H_{\Delta_5})|_{\Uadm}.
 \]
 The class transported from the Borcherds source is
 $u_{\mathrm{tan}}([\beta_{\mathrm{Bor}}])$.
 \item The primary obstruction to extending this first-order class
 through second order is
 \[
 \operatorname{ob}_{\mathrm{ch}}
 \bigl(u_{\mathrm{tan}}([\beta_{\mathrm{Bor}}])\bigr)
 =u_{\mathrm{obs}}\!\left(
 \operatorname{ob}_{\mathrm{Bor}}(\beta_{\mathrm{Bor}})\right)
 \in
 \ChirHoch^3_{\mathrm{CY\text{-}3}}(\mathbf H_{\Delta_5})|_{\Uadm}.
 \]
 By~\eqref{eq:borcherds-obstruction-target-chi3}, this obstruction is
 $[\chi_3^{\mathrm{cand}}]$.
\end{enumerate}
\end{theorem}

\begin{proof}
For the dual numbers \(R=k[\epsilon]/(\epsilon^2)\), a
Maurer--Cartan element based at the origin has the form
\(\epsilon\alpha\) with
\(\alpha\in(\TangCplx_{\mathbf H_{\Delta_5}}^{\mathrm{ch}})^1\).
The Maurer--Cartan equation reduces to \(d\alpha=0\), and an
infinitesimal gauge transformation replaces \(\alpha\) by
\(\alpha+d\xi\).  Its gauge-equivalence class is therefore
\[
 [\alpha]\in
 H^1\!\left(\TangCplx_{\mathbf H_{\Delta_5}}^{\mathrm{ch}}\right)
 =
 \ChirHoch^2_{\mathrm{CY\text{-}3}}
 (\mathbf H_{\Delta_5})|_{\Uadm}.
\]
This is the first-order \(E_1\) deformation class.

A lift over \(k[t]/(t^3)\) has the form
\(t\alpha+t^2\alpha_2\).  The coefficient of \(t^2\) in the
Maurer--Cartan equation is
\[
 d\alpha_2+\frac12\ell^{\mathrm{ch}}_2(\alpha,\alpha)=0.
\]
Hence the primary extension obstruction is
\[
 \operatorname{ob}_{\mathrm{ch}}(\alpha)
 =
 \left[\frac12\ell^{\mathrm{ch}}_2(\alpha,\alpha)\right]
 \in
 H^2\!\left(\TangCplx_{\mathbf H_{\Delta_5}}^{\mathrm{ch}}\right)
 =
 \ChirHoch^3_{\mathrm{CY\text{-}3}}
 (\mathbf H_{\Delta_5})|_{\Uadm}.
\]
Higher \(L_\infty\) brackets enter at the corresponding higher powers
of \(t\).

Completeness of the filtrations makes the series
\[
 \mathcal U_{\mathrm{Bor}\to\mathrm{ch},*}(\gamma)
 =
 \sum_{n\geq1}\frac1{n!}U_n(\gamma,\ldots,\gamma)
\]
convergent for every Maurer--Cartan element
\(\gamma\in\gBor\widehat\otimes\mathfrak m_R\).  The \(L_\infty\)
identities make this series a natural transformation
\[
 \LMC_{\gBor}(R)\longrightarrow
 \LMC_{\TangCplx_{\mathbf H_{\Delta_5}}^{\mathrm{ch}}}(R).
\]
At first order it is \(U_1\), which gives \(u_{\mathrm{tan}}\).
The arity-two \(L_\infty\) identity gives, in cohomology,
\[
 \left[\frac12\ell^{\mathrm{ch}}_2
   (U_1\beta_{\mathrm{Bor}},U_1\beta_{\mathrm{Bor}})\right]
 =
 H^2(U_1)\left[
   \frac12\ell^{\mathrm{Bor}}_2
   (\beta_{\mathrm{Bor}},\beta_{\mathrm{Bor}})\right].
\]
This is the asserted equality between the target obstruction and
\(u_{\mathrm{obs}}(\operatorname{ob}_{\mathrm{Bor}}
(\beta_{\mathrm{Bor}}))\).  Equation~\eqref{eq:borcherds-obstruction-target-chi3}
identifies that class
with \([\chi_3^{\mathrm{cand}}]\).  Hypothesis~\ref{hyp:fmp-governed-chiral} supplies
the formal-moduli equivalence~\eqref{eq:lurie-FMP-chirhoch}.

\end{proof}

\begin{remark}[Collision degree at three points]
\label{rem:chi3-collision-degree}
On an affine chart of the configuration space of three points on a
curve, the logarithmic forms
\[
 \eta_{ij}=\frac{d\log(z_i-z_j)}{2\pi\mathrm i}
\]
obey the Arnold relation
\[
 \eta_{12}\wedge\eta_{23}
 +\eta_{23}\wedge\eta_{31}
 +\eta_{31}\wedge\eta_{12}=0.
\]
The collision-normal cotangent space has rank two, and consequently
\[
 \eta_{12}\wedge\eta_{23}\wedge\eta_{31}=0.
\]
A numerical residue pairing for \([\chi_3^{\mathrm{cand}}]\) specifies the
degree-two collision class and the additional CY-$3$ or Serre factor
that supplies the remaining cohomological degree.
\end{remark}

\begin{remark}[Filtered Siegel--Borcherds deformation transport]
\label{rem:chirhoch3-six-paths}
The filtered transport
\[
 \mathcal U_{\mathrm{Bor}\to\mathrm{ch},*}\colon
 \LMC_{\gBor}\longrightarrow
 \LMC_{\TangCplx_{\mathbf H_{\Delta_5}}^{\mathrm{ch}}}.
\]
Its linear term sends the automorphic first-order direction to
\(\ChirHoch^2_{\mathrm{CY\text{-}3}\!}\); its quadratic obstruction
map lands in
\(\ChirHoch^3_{\mathrm{CY\text{-}3}\!}\).
Equation~\eqref{eq:borcherds-obstruction-target-chi3} places
\([\chi_3^{\mathrm{cand}}]\) in this second group.  The construction
problem \(\mathsf O_{\mathrm{Bor}\to\mathrm{ch}}\) asks for this
filtered \(L_\infty\) morphism and its comparison equation.  An
all-order compact Hall associator consists of a compact Hall source,
filtration-compatible higher Taylor coefficients, and the completed
pentagon identity.
\end{remark}

\begin{remark}[Classical and chiral deformation degrees]
\label{rem:chirhoch3-derived-vs-classical}
Gerstenhaber's deformation complex for an associative algebra \(A\)
is \(C^\bullet(A,A)[1]\).  Its degree-one cohomology is
\(\HH^2(A,A)\), the space of first-order multiplication
deformations, and its degree-two cohomology is \(\HH^3(A,A)\), the
receptacle for their primary extension obstructions~\cite{Ger63}.
The CY-\(3\) product-ambient chiral complex obeys the
same shift:
\[
 H^1(\TangCplx_{\mathbf H_{\Delta_5}}^{\mathrm{ch}})
 =\ChirHoch^2_{\mathrm{CY\text{-}3}}(\mathbf H_{\Delta_5}),
 \qquad
 H^2(\TangCplx_{\mathbf H_{\Delta_5}}^{\mathrm{ch}})
 =\ChirHoch^3_{\mathrm{CY\text{-}3}}(\mathbf H_{\Delta_5}).
\]
Accordingly, the automorphic tangent direction belongs to the first
group and \([\chi_3^{\mathrm{cand}}]\) belongs to the
primary-obstruction group under
\eqref{eq:borcherds-obstruction-target-chi3}.
The curve-level complex and the CY-\(3\) product-ambient complex are
separate deformation problems, related by the specified
\(\Phi_3^{(\Sigma_2,C)}\) comparison.
\end{remark}

\begin{corollary}[Functoriality of the obstruction pairing]
\label{cor:chirhoch3-six-fold-confirmation}
\ClaimStatusConditional
For each source index \(i\), let
\[
 \mathcal U_i\colon\mathfrak g_i\longrightarrow
 \TangCplx_{\mathbf H_{\Delta_5}}^{\mathrm{ch}}
\]
be a filtered $L_\infty$ morphism, and let
\(\beta_i\in Z^1(\mathfrak g_i)\).  Suppose that the induced obstruction
classes satisfy
\begin{equation}
\label{eq:chirhoch3-six-fold-value}
 H^2(U_{i,1})
 \left[\frac12\ell_2^i(\beta_i,\beta_i)\right]
 =[\chi_3^{\mathrm{cand}}].
\end{equation}
For every closed dual class
$[e_3]\in H^2(\TangCplx_{\mathbf H_{\Delta_5}}^{\mathrm{ch}})^\vee$,
the resulting numerical pairings agree:
\[
 \left\langle
 H^2(U_{i,1})
 \left[\frac12\ell_2^i(\beta_i,\beta_i)\right],
 [e_3]\right\rangle
 =\langle[\chi_3^{\mathrm{cand}}],[e_3]\rangle.
\]
A numerical value for the right-hand side is obtained by evaluating an
explicit cocycle representative on an explicit dual cycle.
\end{corollary}

\begin{proof}
Pair equation~\eqref{eq:chirhoch3-six-fold-value} with the closed dual
class $[e_3]$.
\end{proof}

\section{Chiral non-commutative Hodge-to-de-Rham degeneration and the
Mukai--Hochschild pairing on $\mathbf H_{\Delta_5}$}
\label{sec:chirhoch-mukai-nc-hodge}

\providecommand{\ChirHP}{\mathrm{ChirHP}}
\providecommand{\ChirHCneg}{\mathrm{ChirHC}^{-}}
\providecommand{\ChirHCper}{\mathrm{ChirHC}^{\mathrm{per}}}
\providecommand{\NCHodge}{\mathrm{NC\text{-}Hodge}}
\providecommand{\KalDeg}{\mathrm{Kal}}
\providecommand{\FacInf}{\mathrm{Fac}^\infty}
\providecommand{\MukPair}{\langle\,\cdot\,,\,\cdot\,\rangle_{\Muk}}
\providecommand{\sSer}{\mathbb S_{\mathrm{CY}_3}}
\providecommand{\KoszLoc}{\Uadm}
\providecommand{\ConnesB}{B}

Kaledin's degeneration theorem~\cite{Kal08} applies to the mixed complex of a smooth
proper \(\mathbb C\)-linear dg category.  Passage to the
product-ambient chiral complex therefore requires a comparison of mixed
complexes.

\begin{theorem}[Transport of noncommutative Hodge degeneration]
\label{thm:chi-hodge-derham-degenerate-delta5}
\ClaimStatusConditional
Let \(\mathcal C\) be a smooth proper \(\mathbb C\)-linear dg category,
and let
\[
 \Theta_{\mathrm{mix}}\colon
 \bigl(C_\bullet(\mathcal C),b,B\bigr)
 \longrightarrow
 \bigl(C^{\mathrm{prod}}_\bullet(\mathbf H),
       D_{\mathrm{prod}},B_{\mathrm{prod}}\bigr)
\]
be a filtered quasi-isomorphism of complete mixed complexes.  Assume
that the filtrations are exhaustive, separated, and bounded in each
total degree.  Then the product-ambient Hodge-to-de~Rham spectral
sequence
\begin{equation}
\label{eq:chiral-hodge-de-rham-SS}
 E_1^{p,q}
 =
 H_{p+q}\bigl(C^{\mathrm{prod}}_\bullet(\mathbf H),
 D_{\mathrm{prod}}\bigr)u^{-p}
 \Longrightarrow
 H_{p+q}\bigl(
 C^{\mathrm{prod}}_\bullet(\mathbf H)[[u]],
 D_{\mathrm{prod}}+uB_{\mathrm{prod}}\bigr)
\end{equation}
degenerates at \(E_1\).
\end{theorem}

\begin{proof}
Kaledin's theorem~\cite{Kal08} gives \(E_1\)-degeneration for
\((C_\bullet(\mathcal C),b,B)\).  The filtered quasi-isomorphism
\(\Theta_{\mathrm{mix}}\) induces an isomorphism of the two spectral
sequences from the \(E_1\)-page onward.  Completeness and boundedness
give convergence, hence the product-ambient sequence degenerates at
\(E_1\).
\end{proof}

\begin{remark}[The comparison locus]
\label{rem:chi-hodge-derham-humbert-obstruction}
The locus of transported degeneration is the locus on which
\(\mathcal C\) is smooth and proper and
\(\Theta_{\mathrm{mix}}\) is a filtered quasi-isomorphism.  Extension
across a Humbert stratum asks for an extension of this mixed-complex map
and a proof of its filtered quasi-isomorphism there.  The resulting
cone measures the extension obstruction and its higher spectral-sequence
differentials.
\end{remark}

\begin{definition}[Product-ambient CY-threefold pairing]
\label{def:chiral-mukai-hochschild}
\index{Mukai!chiral Hochschild pairing}%
Let
\((C^\bullet_{\mathrm{prod}}(\mathbf H),D_{\mathrm{prod}})\)
be the complex of
Definition~\ref{def:chi3-product-ambient-datum}.  A
\emph{product-ambient CY-threefold pairing} is a family of bilinear maps
\begin{equation}
\label{eq:chiral-mukai-hochschild-pairing-def}
 \langle-,-\rangle_{\Muk}^{\mathrm{ch}}\colon
 C^p_{\mathrm{prod}}(\mathbf H)\otimes
 C^{3-p}_{\mathrm{prod}}(\mathbf H)\longrightarrow\mathbb C
\end{equation}
satisfying
\[
 \langle D_{\mathrm{prod}}\alpha,\beta\rangle_{\Muk}^{\mathrm{ch}}
 +(-1)^p
 \langle\alpha,D_{\mathrm{prod}}\beta\rangle_{\Muk}^{\mathrm{ch}}=0
\]
for homogeneous \(\alpha\in C^p_{\mathrm{prod}}(\mathbf H)\).  It is
perfect when the induced maps
\[
 H^p(C^\bullet_{\mathrm{prod}}(\mathbf H))
 \longrightarrow
 H^{3-p}(C^\bullet_{\mathrm{prod}}(\mathbf H))^\vee
\]
are isomorphisms.

A categorical construction of this pairing consists of a smooth proper
CY-threefold dg category \(\mathcal C\), its negative-cyclic
Calabi--Yau trace, and a chain map
\[
 \Phi_{\mathcal C}\colon
 C^\bullet_{\mathrm{cat}}(\mathcal C)
 \longrightarrow C^\bullet_{\mathrm{prod}}(\mathbf H)
\]
which intertwines the categorical trace pairing with
\eqref{eq:chiral-mukai-hochschild-pairing-def} up to a specified chain
homotopy.
\end{definition}

\begin{theorem}[The Mukai pairing criterion for the three-point class]
\label{thm:chi3-mukai-hochschild-pairing}
\ClaimStatusProvedHere
Assume the boundary equation
\eqref{eq:chi3-closure-obligation}.  Let
\(\chi_3^\vee\in Z^0(C^\bullet_{\mathrm{prod}}(\mathbf H))\)
be a closed degree-zero cochain for a product-ambient CY-threefold
pairing.  Put
\begin{equation}
\label{eq:chiral-mukai-hochschild-value}
 \lambda_{\Muk}(\chi_3^{\mathrm{cand}},\chi_3^\vee)
 :=
 \bigl\langle
 \chi_3^{\mathrm{cand}},\chi_3^\vee
 \bigr\rangle_{\Muk}^{\mathrm{ch}}.
\end{equation}
This scalar depends only on the two cohomology classes.  The implication
\[
 \lambda_{\Muk}(\chi_3^{\mathrm{cand}},\chi_3^\vee)\ne0
 \quad\Longrightarrow\quad
 [\chi_3^{\mathrm{cand}}]\ne0,\qquad
 [\chi_3^\vee]\ne0
\]
holds, and the induced pairing between the lines
\(\mathbb C[\chi_3^{\mathrm{cand}}]\) and
\(\mathbb C[\chi_3^\vee]\) is perfect.
\end{theorem}

\begin{proof}
The graded boundary identity in
Definition~\ref{def:chiral-mukai-hochschild} makes the value invariant
under replacement of either cochain by a cohomologous one.  A zero
cohomology class pairs trivially with every closed complementary
cochain.  A nonzero scalar therefore gives two nonzero classes.  On the
two one-dimensional spans, multiplication by
\(\lambda_{\Muk}\in\mathbb C^\times\) is an isomorphism to the dual
line.
\end{proof}

\begin{corollary}[Perfectness on the three-point and degree-zero lines]
\label{cor:chi3-mukai-nondegeneracy}
\ClaimStatusProvedHere
The restriction of
\(\langle-,-\rangle_{\Muk}^{\mathrm{ch}}\) to
\[
 \mathbb C[\chi_3^{\mathrm{cand}}]\otimes
 \mathbb C[\chi_3^\vee]
\]
is perfect exactly when
\(\lambda_{\Muk}(\chi_3^{\mathrm{cand}},\chi_3^\vee)\in
\mathbb C^\times\).  An explicit residue or trace calculation determines
this scalar after representatives of both classes have been constructed.
\end{corollary}

\begin{proof}
The matrix of the restricted pairing is the one-by-one matrix
\((\lambda_{\Muk})\).
\end{proof}

\begin{remark}[Categorical and product-ambient pairings]
\label{rem:chi3-mukai-cross-volume}
Let \(\mathcal C=D^b\operatorname{Coh}(\mathrm{K3}\times E)\) with a
chosen smooth proper dg enhancement and CY trace.  A comparison with the
product-ambient pairing is a homotopy-commutative square
\[
\begin{tikzcd}[column sep=large]
 C^\bullet_{\mathrm{cat}}(\mathcal C)\otimes
 C^\bullet_{\mathrm{cat}}(\mathcal C)
 \ar[r,"\langle-,-\rangle_{\Muk}^{\mathrm{cat}}"]
 \ar[d,"\Phi_{\mathcal C}\otimes\Phi_{\mathcal C}"']
 & \mathbb C \ar[d,equal]\\
 C^\bullet_{\mathrm{prod}}(\mathbf H)\otimes
 C^\bullet_{\mathrm{prod}}(\mathbf H)
 \ar[r,"\langle-,-\rangle_{\Muk}^{\mathrm{ch}}"']
 & \mathbb C .
\end{tikzcd}
\]
This square carries categorical Mukai values to product-ambient values.
The chain map and its homotopy are the mathematical content of the
cross-volume comparison.
\end{remark}

\section{Comparison morphisms for
\texorpdfstring{\(\chi_3^{\mathrm{cand}}\)}{chi3 candidate}}
\label{sec:chirhoch3-seven-path-comparison}

\begin{theorem}[Comparison of realizations by chain homotopy]
\label{thm:chirhoch3-seven-path-comparison}
\ClaimStatusProvedHere
Let \(\mathcal I\) be a finite index set:
\begin{equation}
\label{eq:seven-path-list}
 |\mathcal I|<\infty.
\end{equation}
For each \(i\in\mathcal I\), let \((K_i^\bullet,d_i)\) be a chain
complex with a cocycle
\(x_i\in Z^3(K_i^\bullet)\).  Suppose there are chain maps
\begin{equation}
\label{eq:seven-comparison-maps}
 F_i:K_i^\bullet\longrightarrow
 C^\bullet_{\mathrm{prod}}(\mathbf H)
 \qquad (i\in\mathcal I)
\end{equation}
and cochains \(h_i\in C^2_{\mathrm{prod}}(\mathbf H)\) satisfying
\[
 F_i(x_i)-\chi_3^{\mathrm{cand}}
 =D_{\mathrm{prod}}h_i.
\]
Then
\[
 [F_i(x_i)]=[\chi_3^{\mathrm{cand}}]
 \quad\text{for every }i\in\mathcal I.
\]
For each closed dual chain \(e_3\), the evaluations obey
\begin{equation}
\label{eq:seven-path-common-value}
 \langle F_i(x_i),e_3\rangle
 =\langle\chi_3^{\mathrm{cand}},e_3\rangle
 =\lambda_3(T_3,\iota_{123},s_Y;e_3).
\end{equation}
\end{theorem}

\begin{proof}
The displayed cochains \(h_i\) exhibit equality in cohomology.  Pairing
their boundaries with the closed chain \(e_3\) gives zero, proving
\eqref{eq:seven-path-common-value}.
\end{proof}

\begin{remark}[Source-specific constructions]
\label{rem:chirhoch3-path-G-independence-comparison}
CoHA operations, Donaldson--Thomas series, categorical Hochschild
classes, cyclic classes, deformation obstructions, and Mukai pairings
carry their own differentials and gradings.  Each realization enters
Theorem~\ref{thm:chirhoch3-seven-path-comparison} through its chain
complex \(K_i^\bullet\), its cocycle \(x_i\), its chain map \(F_i\), and
its comparison homotopy \(h_i\).
\end{remark}

\begin{remark}[The common mathematical obligation]
\label{rem:chirhoch3-seven-voice-cadence}
For every source construction, the load-bearing equation is
\[
 F_i(x_i)-\chi_3^{\mathrm{cand}}=D_{\mathrm{prod}}h_i.
\]
An explicit closed dual chain then turns this cochain identity into a
numerical equality.  These two constructions determine the common class
and its value.
\end{remark}

\begin{corollary}[Nonvanishing transported by the comparison maps]
\label{cor:chirhoch3-seven-fold-confirmation}
\ClaimStatusProvedHere
Under the hypotheses of
Theorem~\ref{thm:chirhoch3-seven-path-comparison}, suppose a closed dual
chain \(e_3\) satisfies
\[
 \langle F_j(x_j),e_3\rangle\ne0
\]
for one \(j\in\mathcal I\).  Then
\([\chi_3^{\mathrm{cand}}]\) and every \([F_i(x_i)]\) are nonzero.
\end{corollary}

\begin{proof}
Equation~\eqref{eq:seven-path-common-value} gives the same nonzero
evaluation for every \(i\), and
Theorem~\ref{thm:chirhoch3-Delta5-chain-level} applies.
\end{proof}
